\documentclass[11pt,a4paper,openany,UTF8]{ctexbook}
% ==============================================================================
% 大模型推理系统与核心原理：从零基础到工业实战
% LaTeX 宏包与排版样式配置文件 (preamble.tex)
% ==============================================================================

\usepackage{geometry}
\geometry{
  a4paper,
  left=2.5cm,
  right=2.5cm,
  top=2.8cm,
  bottom=2.8cm,
  headheight=15pt,
  footskip=25pt
}

% ------------------------------------------------------------------------------
% 字体配置 (针对 Ubuntu / Debian Noto CJK 及 DejaVu)
% ------------------------------------------------------------------------------
\usepackage{fontspec}
\setCJKmainfont[BoldFont=Noto Serif CJK SC Bold]{Noto Serif CJK SC}
\setCJKsansfont[BoldFont=Noto Sans CJK SC Bold]{Noto Sans CJK SC}
\setCJKmonofont{Noto Sans Mono CJK SC}
\setmonofont{DejaVu Sans Mono}

% ------------------------------------------------------------------------------
% 数学公式宏包
% ------------------------------------------------------------------------------
\usepackage{amsmath,amssymb,amsfonts,mathtools,bm}

% ------------------------------------------------------------------------------
% 表格与图形
% ------------------------------------------------------------------------------
\usepackage{booktabs,tabularx,array,colortbl,caption}
\usepackage{graphicx}
\usepackage{tikz}
\usetikzlibrary{arrows.meta,positioning,shapes.geometric,calc,backgrounds,fit}

% ------------------------------------------------------------------------------
% 颜色系统定义
% ------------------------------------------------------------------------------
\usepackage{xcolor}
\definecolor{primary}{RGB}{18, 67, 109}       % 经典海军蓝
\definecolor{secondary}{RGB}{0, 122, 135}     % 科技青
\definecolor{accent}{RGB}{217, 83, 30}        % 强调橙
\definecolor{darktext}{RGB}{33, 37, 41}       % 主文本深灰
\definecolor{lightbg}{RGB}{248, 249, 250}     % 浅灰代码背景
\definecolor{boxborder}{RGB}{220, 224, 230}   % 边框灰
\definecolor{tipgreen}{RGB}{40, 135, 78}      % 技巧绿
\definecolor{warnred}{RGB}{192, 57, 43}       % 警告红
\definecolor{mathpurple}{RGB}{108, 52, 131}   % 定理紫

% 代码语法高亮颜色
\definecolor{codegreen}{RGB}{34, 139, 34}
\definecolor{codegray}{RGB}{110, 115, 120}
\definecolor{codepurple}{RGB}{140, 40, 140}
\definecolor{codeblue}{RGB}{10, 60, 170}

% ------------------------------------------------------------------------------
% Listings 工业级代码高亮配置
% ------------------------------------------------------------------------------
\usepackage{listings}
\usepackage{xeCJK-listings}

\lstset{
  basicstyle=\footnotesize\ttfamily,
  breaklines=true,
  breakatwhitespace=false,
  breakautoindent=true,
  columns=fullflexible,
  keepspaces=true,
  showstringspaces=false,
  frame=single,
  rulecolor=\color{boxborder},
  backgroundcolor=\color{lightbg},
  keywordstyle=\color{codeblue}\bfseries,
  commentstyle=\color{codegray}\itshape,
  stringstyle=\color{codegreen},
  numbers=left,
  numberstyle=\tiny\color{codegray},
  numbersep=7pt,
  stepnumber=1,
  tabsize=4,
  captionpos=t,
  abovecaptionskip=6pt,
  belowcaptionskip=4pt,
  xleftmargin=1.2em,
  framexleftmargin=1.2em
}

% ------------------------------------------------------------------------------
% 彩色提示框 (tcolorbox)
% ------------------------------------------------------------------------------
\usepackage{tcolorbox}
\tcbuselibrary{skins,breakable}

% 1. 背景介绍与导读框
\newtcolorbox{noteBox}[1][导读与背景]{
  enhanced, breakable,
  colback=primary!3!white,
  colframe=primary,
  arc=2mm, boxrule=0.8pt,
  fonttitle=\bfseries\sffamily,
  coltitle=white, colbacktitle=primary,
  title={#1},
  top=3mm, bottom=3mm, left=4mm, right=4mm
}

% 2. 核心直觉与物理意义框
\newtcolorbox{tipBox}[1][核心直觉与思考]{
  enhanced, breakable,
  colback=tipgreen!3!white,
  colframe=tipgreen,
  arc=2mm, boxrule=0.8pt,
  fonttitle=\bfseries\sffamily,
  coltitle=white, colbacktitle=tipgreen,
  title={#1},
  top=3mm, bottom=3mm, left=4mm, right=4mm
}

% 3. 避坑指南与警告框
\newtcolorbox{warningBox}[1][注意事项与避坑指南]{
  enhanced, breakable,
  colback=warnred!3!white,
  colframe=warnred,
  arc=2mm, boxrule=0.8pt,
  fonttitle=\bfseries\sffamily,
  coltitle=white, colbacktitle=warnred,
  title={#1},
  top=3mm, bottom=3mm, left=4mm, right=4mm
}

% 4. 数学定理与证明框
\newtcolorbox{mathBox}[1][数学推导与严格证明]{
  enhanced, breakable,
  colback=mathpurple!3!white,
  colframe=mathpurple,
  arc=2mm, boxrule=0.8pt,
  fonttitle=\bfseries\sffamily,
  coltitle=white, colbacktitle=mathpurple,
  title={#1},
  top=3mm, bottom=3mm, left=4mm, right=4mm
}

% ------------------------------------------------------------------------------
% 页面页眉页脚
% ------------------------------------------------------------------------------
\usepackage{fancyhdr}
\pagestyle{fancy}
\fancyhf{}
\fancyhead[LE,RO]{\small\sffamily\thepage}
\fancyhead[RE]{\small\sffamily\color{primary}\nouppercase{\leftmark}}
\fancyhead[LO]{\small\sffamily\color{secondary}\nouppercase{\rightmark}}
\renewcommand{\headrulewidth}{0.4pt}
\renewcommand{\footrulewidth}{0pt}

\fancypagestyle{plain}{
  \fancyhf{}
  \fancyfoot[C]{\small\sffamily\thepage}
  \renewcommand{\headrulewidth}{0pt}
}

% ------------------------------------------------------------------------------
% 标题与章节样式定制 (ctexset)
% ------------------------------------------------------------------------------
\ctexset{
  part = {
    format = \Huge\bfseries\sffamily\centering\color{primary},
    pagestyle = empty
  },
  chapter = {
    format = \huge\bfseries\sffamily\color{primary},
    name = {第, 章},
    number = \arabic{chapter},
    beforeskip = 15pt,
    afterskip = 25pt
  },
  section = {
    format = \Large\bfseries\sffamily\color{primary},
    beforeskip = 18pt,
    afterskip = 10pt
  },
  subsection = {
    format = \large\bfseries\sffamily\color{secondary},
    beforeskip = 14pt,
    afterskip = 8pt
  },
  subsubsection = {
    format = \normalsize\bfseries\sffamily,
    beforeskip = 10pt,
    afterskip = 6pt
  }
}

% ------------------------------------------------------------------------------
% 超链接与书签
% ------------------------------------------------------------------------------
\usepackage{hyperref}
\usepackage{bookmark}
\hypersetup{
  colorlinks=true,
  linkcolor=primary,
  citecolor=tipgreen,
  urlcolor=secondary,
  pdfborder={0 0 0},
  pdftitle={大模型推理系统与核心原理：从零基础到工业实战},
  pdfauthor={sharp-volta AI System Learning Project},
  pdfsubject={Large Language Model Inference, Systems, CUDA, Attention, Quantization},
  pdfkeywords={LLM, Inference, KV Cache, FlashAttention, Speculative Decoding, PagedAttention, Tensor Parallelism}
}


\begin{document}

% ----------------------------------------------------------------------
% 封面设计
% ----------------------------------------------------------------------
\begin{titlepage}
\centering
\vspace*{2.5cm}
{\Huge\bfseries\sffamily\color{primary} 大模型推理系统与核心原理}\\[0.6cm]
{\LARGE\sffamily\color{secondary} 从零基础数学推导到现代工业工程落地}\\[1.2cm]
{\large\sffamily\color{darktext} 深入自注意力、KV Cache、Roofline 模型、FlashAttention、量化与分布式并行}\\[2.5cm]

\begin{tcolorbox}[enhanced, width=0.88\textwidth, colback=primary!5!white, colframe=primary, arc=3mm, boxrule=1pt, center]
\centering\sffamily\small
\textbf{涵盖 14 篇系统理论全景推导 + 14 组完整可运行实验源码精解}\\[0.3cm]
Transformer $\cdot$ RoPE $\cdot$ SwiGLU $\cdot$ GQA/MLA $\cdot$ PagedAttention $\cdot$ Speculative Decoding $\cdot$ Tensor Parallelism
\end{tcolorbox}

\vfill
{\large\sffamily sharp-volta 开源系统学习工程组 编著}\\[0.3cm]
{\small\sffamily 2026 年 9 月 $\cdot$ 第一版}
\vspace*{1.5cm}
\end{titlepage}

% ----------------------------------------------------------------------
% 前言与目录
% ----------------------------------------------------------------------
\frontmatter
\pagestyle{plain}
\chapter*{前言：大模型推理的底层物理规律}
\addcontentsline{toc}{chapter}{前言：大模型推理的底层物理规律}
\section{[起步] 大模型推理（LLM Inference）系统学习指南}

欢迎开启大模型推理学习之旅！

在现代 AI 系统中，\textbf{模型训练是「一次性成本」，而模型推理是「长期且持续膨胀的在线成本」}。大模型推理的瓶颈通常不是算力（Compute Bound），而是\textbf{显存容量与内存带宽（Memory Bandwidth Bound）}。理解并掌握大模型推理的核心机制，是从事 AI Infra、算法部署和高性能计算（HPC）的关键技能。

\vspace{0.8em}\hrule\vspace{0.8em}

\section{[指南] 从这里开始}

\textbf{--> 先读 \textbf{LEARNING\_PATH.md：完整学习路径}。} 它规定了先学什么、后学什么、每个阶段需要先补哪些外部资源（具体到哪一集、哪一章），以及每个阶段的"出关检验"和"卡住了回去补哪里"。
下面的列表只是索引。

\textbf{环境准备}（RTX 50 系显卡必须使用 CUDA 12.8 版本的 PyTorch）：

\begin{lstlisting}[language=bash]
python3 -m venv .venv && source .venv/bin/activate
pip install -r requirements.txt
pip install torch --index-url https://download.pytorch.org/whl/cu128
\end{lstlisting}

\vspace{0.8em}\hrule\vspace{0.8em}

\section{[重点] 【零基础首选】小白到精通·原理与数学公式完整推导教程}

\begin{mathBox}[核心数学定理推导]
适合人群：\textbf{没有任何深度学习基础、对公式推导和历史背景充满好奇的零基础学习者}。包含从 1950 年代规则系统到 Transformer 的完整演进史，以及所有核心公式的逐行代数证明：
\end{mathBox}

\begin{enumerate}
  \item {}[概览] \textbf{\textbf{第 0 篇（导读）：AI 深度学习全景族谱与零基础入门指南}}
\end{enumerate}

\begin{itemize}
  \item 从感知机到 CNN、ResNet、RNN、Transformer 的科技树全景族谱，各种“Net”到底都是干什么的，不踩坑入门资源与学习路线。
\end{itemize}

\begin{enumerate}
  \item {}[阅读] \textbf{\textbf{第 1 篇：计算机如何读懂语言？从人工规则到大模型演进史}}
\end{enumerate}

\begin{itemize}
  \item 计算机的本质、人工规则为何崩溃、N-gram 语言模型、神经网络与反向传播底座、RNN/LSTM 的长程遗忘与并行死穴、Transformer 的颠覆。
\end{itemize}

\begin{enumerate}
  \item {}[几何] \textbf{\textbf{第 2 篇：零基础必备数学直觉：向量、空间变换、点积与 Softmax 深度推导}}
\end{enumerate}

\begin{itemize}
  \item 标量/向量/张量物理本质、点积与余弦夹角几何意义（为什么代表相似度）、矩阵乘法的空间投影本质、Softmax 的指数来源与防溢出平移不变性证明。
\end{itemize}

\begin{enumerate}
  \item {}[反向] \textbf{\textbf{第 3 篇：神经网络如何学习：导数、链式法则、反向传播与交叉熵}}
\end{enumerate}

\begin{itemize}
  \item 梯度为什么指向下山方向（证明）、反向传播为什么比其它求导方法快、矩阵乘反向公式 $X^TG$ / $GW^T$、交叉熵 = 最大似然、困惑度、\textbf{Softmax 雅可比与 $p - y$ 梯度的完整推导}、梯度消失与"加法通路"、训练 6N vs 推理 2N、训练与推理的显存账本。
\end{itemize}

\begin{enumerate}
  \item {}[推导] \textbf{\textbf{第 4 篇：Transformer 核心公式彻底推导：每一个符号从何而来}}
\end{enumerate}

\begin{itemize}
  \item Q/K/V 检索设计、\textbf{【重点推导】为什么必须除以 $\sqrt{d_k}$（独立变量方差爆炸与 Softmax 梯度消失严格证明）}、因果掩码 $-\infty$、残差连接常数导数高速公路、RoPE 相对位置内积恒等性证明。
\end{itemize}

\begin{enumerate}
  \item {}[架构] \textbf{\textbf{第 5 篇：从原版 Transformer 到 Llama / DeepSeek：每一处改动的历史动机}}
\end{enumerate}

\begin{itemize}
  \item 注意力的真正起源（Bengio 2003 $\rightarrow$ Seq2Seq 瓶颈 $\rightarrow$ Bahdanau）、BPE 分词与词表大小的推理代价、为什么仅解码器胜出、规模定律与"推理成本成为主问题"、Pre-LN / RMSNorm / SwiGLU / RoPE / MQA$\rightarrow$GQA$\rightarrow$MLA / MoE 每一处改动的动机。
\end{itemize}

\begin{enumerate}
  \item {}[加速] \textbf{\textbf{第 6 篇：什么是推理？自回归生成全生命周期与 KV Cache 代数证明}}
\end{enumerate}

\begin{itemize}
  \item 训练与推理的本质鸿沟、Prefill（算力受限）vs Decode（访存受限）、KV Cache 代数展开冗余消除证明、显存爆炸计算公式与 PagedAttention 诞生背景。
\end{itemize}

\begin{enumerate}
  \item {}[决策] \textbf{\textbf{第 7 篇：大模型的决策大脑：Logits、温度系数极限证明与采样算法}}
\end{enumerate}

\begin{itemize}
  \item LM Head 词表投影、\textbf{【定理推导】温度参数 $T \to 0$（唯一最大值时收敛到贪婪）与 $T \to \infty$（均匀分布）严格数学极限证明}、Top-K 截断、Top-P 动态核采样与重复惩罚。
\end{itemize}

\begin{enumerate}
  \item {}[硬件] \textbf{\textbf{第 8 篇：硬件视角：GPU 存储层级、FLOPs、算术强度与浮点格式}}
\end{enumerate}

\begin{itemize}
  \item Roofline 模型、\textbf{推导"Decode 线性层算术强度 $\approx$ batch size"与"Decode 注意力算术强度 $\approx$ GQA 分组大小"}、Decode 耗时模型、算子融合、FP32 / BF16 / FP16 / FP8 / FP4 的位布局与精度陷阱。
\end{itemize}

\begin{enumerate}
  \item {}[加速] \textbf{\textbf{第 9 篇：FlashAttention：Online Softmax 的严格推导与 IO 复杂度}}
\end{enumerate}

\begin{itemize}
  \item Online Softmax 的归纳法证明、分块注意力、IO 复杂度、FlashDecoding 的 split-K 合并公式及其结合律。
\end{itemize}

\begin{enumerate}
  \item {}[采样] \textbf{\textbf{第 10 篇：投机解码：无偏性证明与期望加速比}}
\end{enumerate}

\begin{itemize}
  \item \textbf{投机采样无偏性定理的完整证明}、接受率与总变差距离、期望加速比公式与最优 K、为什么大 batch 时收益消失。
\end{itemize}

\begin{enumerate}
  \item {}[量化] \textbf{\textbf{第 11 篇：量化：从每位 6 dB 到 SmoothQuant、AWQ 与 GPTQ}}
\end{enumerate}

\begin{itemize}
  \item 每位 6 dB 规律、激活离群值、量化权重还是激活（由 Roofline 决定）、SmoothQuant / AWQ / GPTQ 的数学原理（含两变量的 GPTQ 直观例子）。
\end{itemize}

\begin{enumerate}
  \item {}[系统] \textbf{\textbf{第 12 篇：推理服务系统：PagedAttention、前缀缓存与调度}}
\end{enumerate}

\begin{itemize}
  \item TTFT / TPOT / Goodput 指标、Little 定律、显存碎片与分页、写时复制、前缀缓存、Chunked Prefill 与 P/D 分离。
\end{itemize}

\begin{enumerate}
  \item {}[并行] \textbf{\textbf{第 13 篇：多卡并行推理：张量并行、流水线并行与专家并行}}
\end{enumerate}

\begin{itemize}
  \item 通信原语与 Ring All-Reduce、Megatron "先列后行"切法的推导、PP 为什么不降延迟、MoE 的专家并行。
\end{itemize}

\vspace{0.8em}\hrule\vspace{0.8em}

\section{学习路线图 (Learning Roadmap)}

完整的阶段划分、前置条件与外部资源见 \textbf{\textbf{LEARNING\_PATH.md}}。简版：

\vspace{0.8em}\hrule\vspace{0.8em}

\section{[目录] 本学习仓库目录结构}

实验目录 \texttt{labNN\_*} 的 \textbf{NN 就是它配套的那一篇}（一篇配多个实验时用 a/b 区分）。每个目录都有说明文档和可直接运行的代码：

\begin{center}
\small
\begin{tabularx}{\textwidth}{p{2.5cm} p{3.5cm} X X}
\toprule
\textbf{对应篇} & \textbf{目录} & \textbf{核心内容} & \textbf{重点代码} \\
\midrule
\textbf{第 3 篇} & {}\textbf{\texttt{lab03\_autograd\_from\_scratch/}} & 约 350 行纯 Python 自动求导，数值验证 Softmax 雅可比、交叉熵梯度、矩阵乘反向公式、梯度消失、XOR & \texttt{micrograd\_from\_scratch.py} \\
\textbf{第 4 篇} & {}\textbf{\texttt{lab04\_transformer\_microscope/}} & \textbf{【极细基础】}从分词到下一个字的微观全貌、张量维度流动、RoPE 旋转、置换等变性与长程衰减的正确表述 & \texttt{01\_tensor\_trace\_step\_by\_step.py}\newline \texttt{02\_rope\_and\_attention\_microscope.py} \\
\textbf{第 5 篇} & {}\textbf{\texttt{lab05\_train\_tiny\_llama/}} & 用本仓库教程做语料训练一个迷你 Llama（RMSNorm / RoPE / GQA / SwiGLU），再做 KV Cache 对拍与采样 & \texttt{model.py}\newline \texttt{train.py}\newline \texttt{generate.py} \\
\textbf{第 6 篇} & {}\textbf{\texttt{lab06\_kv\_cache/}} & Prefill 与 Decode 阶段区别、自回归生成的重复计算痛点、KV Cache 机制与代码对比 & \texttt{kv\_cache\_benchmark.py} \\
\textbf{第 7 篇} & {}\textbf{\texttt{lab07a\_sampling/}} & Greedy、温度、Top-K、Top-P、重复惩罚的概率分布前后对比 & \texttt{sampling\_strategies.py} \\
\textbf{第 7 篇} & {}\textbf{\texttt{lab07b\_qwen\_from\_scratch/}} & \textbf{【关键】}只读原始权重、手写 Qwen2.5 推理，与 HF 官方逐位对拍 logits 与生成结果 & \texttt{qwen\_from\_scratch.py} \\
\textbf{第 8 篇} & {}\textbf{\texttt{lab08\_memory\_and\_roofline/}} & 模型权重、KV Cache、激活值显存计算，计算密度（Arithmetic Intensity）与 Roofline 模型分析 & \texttt{memory\_calculator.py} \\
\textbf{第 9 篇} & {}\textbf{\texttt{lab09\_flash\_attention/}} & Online Softmax、分块 FlashAttention、FlashDecoding 合并，GPU 上朴素 vs SDPA；\textbf{\textbf{真机 CUDA 附录}}：同一套 online softmax 写成 CUDA C++ kernel，实测 DRAM 带宽与占用率 & \texttt{flash\_attention\_numpy.py}\newline \texttt{cuda\_softmax/softmax\_kernels.cu} \\
\textbf{第 10 篇} & {}\textbf{\texttt{lab10\_speculative\_decoding/}} & 投机采样（Speculative Decoding）核心思想、草稿模型生成与主模型并行验证算法 & \texttt{toy\_speculative\_decoding.py} \\
\textbf{第 11 篇} & {}\textbf{\texttt{lab11\_quantization/}} & 对称与非对称量化、Per-Tensor / Per-Channel / Per-Group 量化、AWQ / GPTQ 原理 & \texttt{quant\_basics.py} \\
\textbf{第 12 篇} & {}\textbf{\texttt{lab12a\_continuous\_batching/}} & 静态 Batching 的「气泡（Padding/Bubble）」问题、迭代级调度（Continuous Batching）仿真 & \texttt{continuous\_batching\_sim.py} \\
\textbf{第 12 篇} & {}\textbf{\texttt{lab12b\_paged\_attention/}} & 块分配器、块表、分页注意力、写时复制、前缀缓存、Chunked Prefill & \texttt{paged\_kv\_cache\_sim.py} \\
\textbf{第 13 篇} & {}\textbf{\texttt{lab13\_parallelism/}} & 张量并行 MLP / 注意力、错误切法反例、Ring All-Reduce、多卡 Decode 延迟模型 & \texttt{tensor\_parallel\_sim.py} \\
\textbf{毕业设计} & {}\textbf{\texttt{capstone\_engines\_practice/}} & vLLM / llama.cpp 部署与压测任务、源码阅读路线（毕业设计） & \texttt{bench\_openai\_server.py} \\
\bottomrule
\end{tabularx}
\end{center}

Lab 08、10、11 还各新增了进阶脚本：\texttt{measure\_gpu\_roofline.py}（在你的显卡上实测 Roofline）、\texttt{verify\_unbiasedness.py}（蒙特卡洛验证投机采样定理）、\texttt{advanced\_quant\_numpy.py}（从零实现 SmoothQuant / AWQ / GPTQ）。

\textbf{全仓库唯一一份真 CUDA C++} 在 \textbf{\texttt{lab09\_flash\_attention/cuda\_softmax/}}：把第 9 篇的 Online Softmax 写成 kernel（\texttt{\_\_global\_\_} / \texttt{\_\_shared\_\_} / \texttt{\_\_syncthreads()}），量出 DRAM 带宽、L2 命中与占用率。用 PyTorch 自带的 NVRTC 编译，\textbf{不需要装 CUDA toolkit}。

\vspace{0.8em}\hrule\vspace{0.8em}

\section{[决策] 核心概念速查手册}

\subsection{Prefill（首字）vs Decode（增量生成）阶段}

\begin{itemize}
  \item \textbf{Prefill 阶段（首字延迟 TTFT: Time to First Token）}：
  \item 输入所有 prompt tokens，能够并行计算。
  \item 特征：\textbf{Compute-bound（算力受限）}，GPU 利用率高，矩阵乘法规模大 ($M \times K \times N$, $M > 1$)。
  \item \textbf{Decode 阶段（每个后续 token 的延迟 TPOT: Time Per Output Token）}：
  \item 串行逐字生成，每次输入仅 1 个 token。
  \item 特征：\textbf{Memory-bound（访存带宽受限）}。每次生成 1 个 token 都需要把数十亿参数和全部历史 KV Cache 从显存搬运进 GPU 计算核心。
\end{itemize}

\subsection{KV Cache 显存占用计算公式}

对于标准 Transformer 结构模型：
\begin{itemize}
  \item 参数量：$P$
  \item 层数：$L$
  \item 隐藏层大小：$H$ 或注意力头数 $n_{\text{heads}} \times d_{\text{head}}$
  \item 上下文总长度（Prompt + Generated）：$S$
  \item 批大小：$B$
  \item 精度字节数：$b$（FP16/BF16 为 2 字节，FP8/INT8 为 1 字节，INT4 为 0.5 字节）
\end{itemize}

每个 Token 产生的 Key 和 Value 向量显存（bytes）：

\[
\text{KV per token} = 2 \times L \times H \times b
\]

\textit{(注：如果采用 MQA/GQA，隐藏层维度 $H$ 需按 KV 组数比例缩减，例如 GQA-8 缩减为原来的 $1/8$)}

批次总 KV Cache 显存（bytes）：

\[
\text{Total KV Cache} = B \times S \times 2 \times L \times H \times b
\]

\begin{noteBox}[核心导读与背景]
\textbf{示例（Llama-3-8B, 32 层, 隐藏层 4096, GQA 8 个 KV 头 vs 32 个 Q 头, FP16）：}
每个 token 仅需要 $2 \times 32 \times (4096/4) \times 2 = 131,072 \text{ bytes} \approx 128 \text{ KB}$。
长度 8192 时，单并发约 1 GB，若并发达 32，则需 32 GB 显存（远超模型权重自身的 16 GB）！
\end{noteBox}

\vspace{0.8em}\hrule\vspace{0.8em}

\section{[工具] 主流推理引擎横向对比}

\begin{center}
\small
\begin{tabularx}{\textwidth}{p{2.5cm} p{3.5cm} X X}
\toprule
\textbf{框架} & \textbf{适用场景} & \textbf{核心亮点} & \textbf{推荐初学者体验} \\
\midrule
\textbf{\href{https://github.com/vllm-project/vllm}{vLLM}} & 工业界高吞吐部署标准 & PagedAttention 显存零浪费、Continuous Batching、广泛支持各类模型 & $\star$$\star$$\star$$\star$$\star$ \\
\textbf{\href{https://github.com/sgl-project/sglang}{SGLang}} & 复杂 Prompt、Agent/结构化生成、超高吞吐 & RadixAttention（自动 KV 复用）、高性能 Runtime & $\star$$\star$$\star$$\star$$\star$ \\
\textbf{\href{https://github.com/NVIDIA/TensorRT-LLM}{TensorRT-LLM}} & NVIDIA 极致性能压榨 & 深度融合 Kernel、In-flight Batching、FP8/FP4 极致调优 & $\star$$\star$$\star$$\star$ \\
\textbf{\href{https://github.com/ggml-org/llama.cpp}{llama.cpp}} & 边缘设备、CPU、Mac、本地轻量推理 & 零依赖纯 C++、极致 GGUF 量化支持 & $\star$$\star$$\star$$\star$$\star$ \\
\textbf{\href{https://github.com/huggingface/text-generation-inference}{TGI (Text Generation Inference)}} & Hugging Face 生态生产部署 & 简单可靠、支持 FlashAttention、生产级运维特性 & $\star$$\star$$\star$$\star$ \\
\bottomrule
\end{tabularx}
\end{center}

\vspace{0.8em}\hrule\vspace{0.8em}

\section{[目标] 推荐第一步}

\begin{enumerate}
  \item 读 \textbf{LEARNING\_PATH.md}，对照"阶段 0"的表格检查自己缺哪块数学基础。
  \item 读第 0、1 篇建立全局地图，然后按路径进入阶段 2（第 2 篇 $\rightarrow$ 第 3 篇 $\rightarrow$ Lab 03）。
  \item 想先看到"真东西"提提兴趣的话，可以先跑一下：
\end{enumerate}

\begin{lstlisting}[language=bash]
   .venv/bin/python lab05_train_tiny_llama/train.py && .venv/bin/python lab05_train_tiny_llama/generate.py
   
\end{lstlisting}

   花一分钟训练出一个会"背诵本教程"的迷你 Llama——然后再回到路径上，把它的每一行都弄懂。

\tableofcontents
\cleardoublepage

% ----------------------------------------------------------------------
% 正文主体
% ----------------------------------------------------------------------
\mainmatter
\pagestyle{fancy}

\part{核心原理与数学推导篇 (Theory \& Derivations)}
\chapter{AI 深度学习全景族谱与零基础入门指南}

\begin{noteBox}[核心导读与背景]
\textbf{前言}：深度学习领域充斥着各种各样的“Net”（如 CNN、RNN、ResNet、Transformer、Diffusion），初学者极易迷失在浩瀚的术语海洋中。本篇为你梳理整张「AI 科技树」，讲清楚各种网络各司其职的关系，并给出最少弯路、最落地的零基础入门路线。
\end{noteBox}

\vspace{0.8em}\hrule\vspace{0.8em}

\section{[概览] 一、深度学习终极演化科技树（全景族谱）}

整个现代 AI 的发展，本质上只由两条主线驱动：\textbf{视觉线（看图像）} 和 \textbf{语言时序线（读文本）}。它们在经历各自数十年的探索后，最终被 \textbf{Transformer} 完成大一统！

\begin{figure}[htbp]
\centering
\begin{tikzpicture}[
  base/.style={draw=primary!80, fill=lightbg, rounded corners=2mm, align=center, font=\footnotesize\sffamily, inner sep=4pt, line width=0.8pt},
  vnode/.style={base, draw=secondary!80, fill=secondary!8, minimum width=3.8cm},
  nnode/.style={base, draw=accent!80, fill=accent!8, minimum width=3.8cm},
  tform/.style={base, draw=primary, fill=primary!15, line width=1.2pt, font=\bfseries\footnotesize\sffamily, minimum width=8cm},
  arr/.style={-{Stealth[scale=0.8]}, draw=gray!70, line width=0.8pt},
  darr/.style={-{Stealth[scale=0.8]}, dashed, draw=primary, line width=0.9pt}
]

% Level 0 & 1
\node[base, minimum width=6.5cm] (p) {1957 感知机 (Perceptron)\\单神经元，无法解决 XOR 异或};
\node[base, below=0.5cm of p, minimum width=6.5cm] (mlp) {1986 多层感知机 (MLP / BP)\\反向传播与现代神经网络基石};

% Level 2: CNN vs RNN
\node[vnode, below left=0.7cm and 0.4cm of mlp] (cnn) {1998 LeNet / CNN\\局部滑窗提取空间特征};
\node[nnode, below right=0.7cm and 0.4cm of mlp] (rnn) {1980s RNN (循环网络)\\时序备忘录处理文本};

% Level 3: AlexNet vs LSTM
\node[vnode, below=0.5cm of cnn] (alex) {2012 AlexNet (GPU 革命)\\算力与深度爆发};
\node[nnode, below=0.5cm of rnn] (lstm) {1997 LSTM / 2014 GRU\\门控机制缓解长程遗忘};

% Level 4: ResNet vs Seq2Seq
\node[vnode, below=0.5cm of alex] (resnet) {2015 ResNet (何恺明)\\残差连接 $x+F(x)$ 破百层};
\node[nnode, below=0.5cm of lstm] (seq2seq) {2014 Seq2Seq + Attention\\机器翻译学会“看重点”};

% Level 5: Transformer (Centered below both branches)
\node[tform, below=1.0cm of $(resnet.south)!0.5!(seq2seq.south)$] (trans) {2017 Transformer (Google 论文)\\彻底废黜循环时序！全矩阵并行点积！统一语言与多模态};

% Level 6: Successors
\node[base, below left=0.8cm and -0.8cm of trans, minimum width=3.2cm] (bert) {2018 BERT\\双向理解完形填空};
\node[base, below=0.8cm of trans, minimum width=3.6cm, fill=primary!10, draw=primary] (gpt) {2018--2023 GPT-1$\sim$4\\单向自回归接龙};
\node[base, below right=0.8cm and -0.8cm of trans, minimum width=3.2cm] (vit) {2020 ViT\\视觉 Transformer};

% Level 7: Modern LLM Era
\node[base, below=0.6cm of gpt, fill=accent!10, draw=accent, line width=1pt, font=\bfseries\footnotesize\sffamily, minimum width=8.5cm] (llm) {2023 至今：现代大模型时代 (Llama-3, Qwen-2.5, DeepSeek)\\核心转向：推理吞吐、显存架构优化、KV Cache、量化};

% Arrows
\draw[arr] (p) -- (mlp);
\draw[arr] (mlp) -| (cnn);
\draw[arr] (mlp) -| (rnn);
\draw[arr] (cnn) -- (alex);
\draw[arr] (alex) -- (resnet);
\draw[arr] (rnn) -- (lstm);
\draw[arr] (lstm) -- (seq2seq);

\draw[darr] (resnet.south) -- ++(0,-0.4) -| node[near start, left, font=\scriptsize] {残差连接思想} ($(trans.north west)+(0.8,0)$);
\draw[darr] (seq2seq.south) -- ++(0,-0.4) -| node[near start, right, font=\scriptsize] {注意力思想飞跃} ($(trans.north east)+(-0.8,0)$);

\draw[arr] (trans) -- (bert);
\draw[arr] (trans) -- (gpt);
\draw[arr] (trans) -- (vit);
\draw[arr] (gpt) -- (llm);

\end{tikzpicture}
\caption{AI 深度学习演化科技树全景族谱}
\end{figure}

\vspace{0.8em}\hrule\vspace{0.8em}

\section{[指南] 二、各种“Net”究竟都是干什么的？}

当你看到不同论文或文章里提到不同的 Net 时，可以对照下表快速定位：

\begin{center}
\small
\begin{tabularx}{\textwidth}{p{2.5cm} p{3.5cm} X X}
\toprule
\textbf{模型代号} & \textbf{全称} & \textbf{诞生初衷与解决的核心痛点} & \textbf{状态与历史地位} \\
\midrule
\textbf{MLP} & Multi-Layer Perceptron (多层感知机 / 全连接网络) & 最基础的万能函数拟合器，全连接矩阵乘法。但输入长度固定，无法直接处理大图像和长文章。 & 所有现代网络的基本积木（如 Transformer 里的 FFN） \\
\textbf{CNN} & Convolutional Neural Network (卷积神经网络) & \textbf{专为图像而生}。一张 1000x1000 的图片直接用全连接会导致参数爆炸，CNN 用一个小滑窗（卷积核）扫过图片提取边缘和纹理。 & 统治计算机视觉领域（2012-2020），经典代表：LeNet, AlexNet, VGG \\
\textbf{ResNet} & Residual Network (残差网络) & 网络堆得太深（超过 20 层）会导致梯度消失退化。何恺明发明\textbf{残差短路连接（$x + F(x)$）}，让网络轻松堆到 152 层甚至 1000 层！ & \textbf{现代所有大模型的标配机制}（Transformer 每层都有残差连接） \\
\textbf{RNN} & Recurrent Neural Network (循环神经网络) & \textbf{专为文本时序而生}。像穿糖葫芦一样逐字传递隐藏状态备忘录，解决文本长度不一的问题。 & 因串行无法利用 GPU 并行，且长程记忆衰减，已被淘汰 \\
\textbf{LSTM} & Long Short-Term Memory (长短期记忆网络) & 在 RNN 基础上引入遗忘门、输入门、输出门，缓解了 RNN 读到后面忘了前面的死穴。 & 在 2017 年前是机器翻译、语音识别的绝对霸主 \\
\textbf{Transformer} & 基于 Self-Attention 的全新网络架构 & \textbf{彻底废弃循环连接}，任意两词距离 $O(1)$，纯矩阵运算极致利用 GPU。不仅统领自然语言，还通过 ViT 杀入计算机视觉，通过 Sora/Diffusion 杀入视频生成。 & \textbf{当今全人类所有大模型的唯一通用底座} \\
\bottomrule
\end{tabularx}
\end{center}

\vspace{0.8em}\hrule\vspace{0.8em}

\section{[目标] 三、新手零基础想真正入门，该看什么？（绝不踩坑三部曲）}

网上的 AI 教程铺天盖地，90\% 都是死记硬背公式或调包调库，初学者往往陷入“懂了但又没完全懂”的痛苦中。

想要扎实、高效、不浪费时间地入门，\textbf{推荐且仅推荐以下三部封神级的黄金资源}：

\subsection{[第一阶段] 阶段 1：建立直观的数学与几何世界观（耗时约 2\textasciitilde{}3 天）}

\begin{itemize}
  \item \textbf{【绝对首选】3Blue1Brown 经典动画教程《深度学习的本质》（Bilibili / YouTube 免费看）}
  \item 搜索关键词：\texttt{3Blue1Brown 深度学习}
  \item 为什么神：全网最好的几何可视化，用生动的动态图形讲清楚：
  \item 什么是神经网络与激活函数？
  \item 什么是梯度下降（走下山坡）？
  \item 反向传播微积分链式法则的几何动画。
  \item 注意力机制（Attention）到底在几何空间中拉伸了什么。
  \item \textbf{建议}：把这 4 集视频当作看电影，看上两遍，你脑海里的数学直觉就彻底打通了！
\end{itemize}

\subsection{[第二阶段] 阶段 2：动手写代码与经典教科书（耗时约 1\textasciitilde{}2 周）}

\begin{itemize}
  \item \textbf{【教材推荐】李沐（前亚马逊首席科学家）团队《动手学深度学习》(Dive into Deep Learning / d2l.ai)}
  \item 中文官网完全免费在线阅读，配备完整可运行代码。
  \item 从基础张量、线性回归、MLP 讲到现代注意力机制，文字极度平实。
  \item \textbf{【硬核进阶】Andrej Karpathy（OpenAI 联合创始人、前特斯拉 AI 总监）的视频系列《Neural Networks: Zero to Hero》}
  \item 他在油管上从一个只有几行的微型框架 \texttt{micrograd} 开始，用纯 Python 逐行写出一个微型 GPT（nanoGPT）。
  \item 跟着他敲一遍，大模型的黑盒彻底透明。
\end{itemize}

\subsection{[第三阶段] 阶段 3：大模型推理与系统实战（面向生产与高性能）}

\begin{itemize}
  \item \textbf{【避坑指南：不要一上来就学“从零预训练大模型”】}
  \item 训练一个 70B 模型需要几千张英伟达 H100 显卡和几千万美元预算，对绝大多数工程师而言，学如何从头训练一个大模型完全脱离现实。
  \item \textbf{【正确的工程路径：深耕大模型推理（LLM Inference）】}
  \item 理解模型的前向传播逻辑与微观张量流（参见本项目的 \texttt{lab04\_transformer\_microscope}）；
  \item 理解为什么推理的核心瓶颈在于\textbf{显存容量与内存带宽（Memory Bandwidth）}（参见 \texttt{lab08\_memory\_and\_roofline}）；
  \item 深入业界工业标准框架（\textbf{vLLM}、\textbf{SGLang}、\textbf{llama.cpp}），学习它们的核心技术：\textbf{PagedAttention}、\textbf{连续批处理（Continuous Batching）}、\textbf{投机采样}、\textbf{量化}；
  \item 这也是当前 AI Infra、大模型私有化部署、高性能边缘计算最紧缺、含金量最高的技术方向！
\end{itemize}

\vspace{0.8em}\hrule\vspace{0.8em}

\begin{noteBox}[核心导读与背景]
{}[指南] \textbf{本仓库的完整路线}：上面三部曲是外部资源的推荐；本仓库自己的 14 篇理论（第 0\textasciitilde{}13 篇）+ 配套实验，以及它们与外部资源如何穿插学习，见 \textbf{LEARNING\_PATH.md}。
族谱里没有展开的细节——注意力在 2014 年的真正起源、从原版 Transformer 到 Llama/DeepSeek 的每一处改动——在 \textbf{第 5 篇}。
\end{noteBox}

\vspace{0.8em}\hrule\vspace{0.8em}

\textbf{下一篇}：\textbf{第 1 篇：计算机如何读懂语言？从人工规则到大模型演进史} ｜ \textbf{总路线}：\textbf{LEARNING\_PATH.md}

\chapter{计算机如何读懂语言？从人工规则到大模型演进史}

\begin{noteBox}[核心导读与背景]
\textbf{阅读前须知}：本篇不需要任何高等数学或编程背景。我们将从人类对计算机的最初探索讲起，带你一步步理解大模型（LLM）究竟是在什么背景下被发明出来的，以及它究竟解决了前人什么无法逾越的难题。
\end{noteBox}

\vspace{0.8em}\hrule\vspace{0.8em}

\section{[代码] 一、计算机的天然缺陷：它只懂数字，不懂人类语言}

首先请思考一个最本质的事实：
\textbf{计算机的硬件核心（CPU 和 GPU），本质上只是无数由晶体管组成的微型电路，它唯一擅长的事情，就是极速计算浮点数的「加法」和「乘法」。}

计算机根本不知道什么是“猫”，什么是“开心”，什么是“我爱你”。
对于计算机而言，屏幕上的汉字或英文字符，只是一串无意义的二进制编码（比如 UTF-8 编码中，“大”这个字的十六进制代码是 \texttt{0xE5 0xA4 0xA7}）。

人类要想让计算机能对话、能写诗、能编程，就必须完成一个根本性的跨越：
\textbf{如何把人类丰富、带有歧义、充满上下文语境的自然语言，转化为数字？并让计算机通过纯粹的数学计算，“理解”语言并生成语言？}

为了解决这个问题，人类科学家经历了整整四代探索。

\vspace{0.8em}\hrule\vspace{0.8em}

\section{[文献] 二、第一代：基于人工规则与专家系统（1950s - 1980s）}

在计算机科学诞生初期，科学家们认为人类语言是有严格语法规则的。

\subsection{核心做法：}

语言学家人工编写成千上万条逻辑规则：
\begin{itemize}
  \item \texttt{句子 = 主语 + 谓语 + 宾语}
  \item 如果主语是“我”，谓语可以是“吃”，宾语可以是“苹果”。
  \item 遇到英语单词 \texttt{"bank"}，如果前文有 \texttt{"river"}，就翻译为“河岸”；如果有 \texttt{"money"}，就翻译为“银行”。
\end{itemize}

\subsection{为什么彻底失败？}

\begin{enumerate}
  \item \textbf{语言规则无穷无尽}：自然语言充斥着成语、倒装句、隐喻、俚语和错别字，规则写得越多，规则之间就越互相冲突。
  \item \textbf{缺乏泛化能力}：人类只要说一句规则库里没有的表达（例如：“今天我被老板炒了鱿鱼”），系统就会瞬间崩溃。
\end{enumerate}

事实证明：\textbf{靠人类肉眼穷举语言规则是死路一条。}

\vspace{0.8em}\hrule\vspace{0.8em}

\section{[图表] 三、第二代：统计语言模型与 N-gram（1990s - 2000s）}

既然人工写规则走不通，统计学家们提出了全新思路：\textbf{不要人工写规则，让计算机去统计真实人类书写的大量文章！}

\subsection{什么是“语言模型（Language Model）”？}

数学家对语言模型的定义极其质朴：
\begin{noteBox}[核心导读与背景]
\textbf{语言模型，就是一个计算“一句话在现实世界中出现的概率”的数学函数。}
\end{noteBox}

例如，对于两句话：
\begin{itemize}
  \item $S_1$ = “今天天气真好”
  \item $S_2$ = “好真天气今天”
\end{itemize}

一个好的语言模型应该计算出：$P(S_1) = 0.05$（很常见），而 $P(S_2) = 0.0000001$（人类几乎不这么说）。

根据概率论的\textbf{链式法则（Chain Rule）}，一句话出现的联合概率，等于逐个词出现的条件概率相乘：

\[
P(w_1, w_2, w_3, \dots, w_n) = P(w_1) \times P(w_2 \mid w_1) \times P(w_3 \mid w_1, w_2) \dots \times P(w_n \mid w_1, \dots, w_{n-1})
\]

这个公式看似完美，但计算机在实际计算时崩溃了：
当句子长达 20 个字时，条件概率 $P(w_{20} \mid w_1, \dots, w_{19})$ 要求计算机去统计：\textbf{在历史上所有文章中，前面 19 个字完全一模一样的情况下，第 20 个字是 $w_{20}$ 的频率}。
现实中，哪怕整个互联网的所有文本加起来，前面 19 个字完全一样的概率也几乎是 \textbf{0}！这就是著名的\textbf{「数据稀疏与维度灾难」}。

\subsection{N-gram 的妥协（马尔可夫假设）}

为了能算出来，统计学家做了一个妥协：\textbf{假定第 $t$ 个词出现的概率，只与它前面的 $N-1$ 个词有关，跟更早以前的词完全无关！}
\begin{itemize}
  \item \textbf{2-gram（Bi-gram）}：每个词只看前 1 个词：$P(w_t \mid w_1, \dots, w_{t-1}) \approx P(w_t \mid w_{t-1})$
  \item \textbf{3-gram（Tri-gram）}：每个词只看前 2 个词：$P(w_t \mid w_1, \dots, w_{t-1}) \approx P(w_t \mid w_{t-2}, w_{t-1})$
\end{itemize}

\subsection{致命死穴：}

\begin{enumerate}
  \item \textbf{短视，无法理解长文本}：如果你写一篇文章，开头说“小明是个法国人\dots{}\dots{}”，在经过 500 字后问“他最擅长的母语是\_\_\_”，N-gram 只看前两个词，根本不可能答对。
  \item \textbf{没有任何概念理解}：在统计模型中，“苹果”和“梨”是两个完全不相关的独立符号。即使模型见过“我吃了一个苹果”，它也无法推断出“我吃了一个梨”是否合理。
\end{enumerate}

\vspace{0.8em}\hrule\vspace{0.8em}

\section{[决策] 四、第三代：分布式词向量与循环神经网络（2010s）}

2013 年，Google 提出了里程碑式的 \textbf{Word2Vec}，彻底颠覆了符号表示法。

\subsection{核心飞跃：词嵌入向量（Word Embedding）}

人类想通了：\textbf{不要用离散的孤立代号表示词，而是把每个词表示为高维空间里的一个坐标点（一个向量，比如 512 个数字）！}
\begin{itemize}
  \item 著名的语言学假设：\textit{“一个词的含义，由它经常与哪些词同时出现所决定（Distributional Hypothesis）。”}
  \item 如果把词映射到多维空间中，在大量文本训练后，你会惊奇地发现数学几何与人类常识奇迹般地重合了：
\end{itemize}

\[
\vec{v}_{\text{国王}} - \vec{v}_{\text{男人}} + \vec{v}_{\text{女人}} \approx \vec{v}_{\text{王后}}
\]

  “苹果”和“梨”在空间中的距离会非常近，而“苹果”和“量子力学”距离极远。

\vspace{0.8em}\hrule\vspace{0.8em}

\subsection{补齐底座：什么是神经网络？它是怎么“学会”知识的？}

\begin{tipBox}[核心直觉与思考]
{}[思考] \textbf{很多教材直接从词向量跳到 RNN，导致零基础同学一头雾水。到底什么是“神经元”？什么是“隐藏层”？它到底是怎么从数据中学习的？我们用最接地气的生活例子彻底讲明白。}
\end{tipBox}

\subsubsection{(1) 单个人工神经元：一个带旋钮的打分器}

生物学中，人脑神经元通过树突接收多个电信号，在细胞核汇总，超过一定阈值后通过轴突激发传递。
数学家把这个过程抽象成了一个极其简单的数学公式：

\begin{figure}[htbp]
\centering
\begin{tikzpicture}[scale=0.88, >=stealth,
  input/.style={rectangle, rounded corners=2pt, draw=secondary, fill=secondary!10, font=\sffamily\footnotesize},
  weight/.style={rectangle, rounded corners=2pt, draw=primary, fill=primary!10, font=\sffamily\footnotesize},
  neuron/.style={circle, draw=warnred, fill=warnred!10, font=\sffamily\footnotesize\bfseries, minimum size=0.9cm}
]
  \node[input] (x1) at (0, 1.5) {特征 $x_1$ (离地铁距离)};
  \node[input] (x2) at (0, 0) {特征 $x_2$ (房屋面积)};
  \node[input] (x3) at (0, -1.5) {特征 $x_3$ (是否带学区)};
  
  \node[weight] (w1) at (3.2, 1.5) {权重 $w_1$};
  \node[weight] (w2) at (3.2, 0) {权重 $w_2$};
  \node[weight] (w3) at (3.2, -1.5) {权重 $w_3$};
  
  \node[neuron] (sum) at (6.0, 0) {$\sum w_i x_i + b$};
  \node[weight, fill=tipgreen!15, draw=tipgreen] (act) at (8.5, 0) {激活函数 $f$};
  \node[input, fill=primary!15, draw=primary] (out) at (10.8, 0) {预测输出 $\hat{y}$};
  
  \draw[->, thick] (x1) -- (w1);
  \draw[->, thick] (x2) -- (w2);
  \draw[->, thick] (x3) -- (w3);
  
  \draw[->, thick] (w1) -- (sum);
  \draw[->, thick] (w2) -- (sum);
  \draw[->, thick] (w3) -- (sum);
  
  \draw[->, thick] (sum) -- (act);
  \draw[->, thick] (act) -- (out);
\end{tikzpicture}
\caption{单神经元（感知机）加权求和与非线性激活计算图}
\end{figure}

\begin{enumerate}
  \item \textbf{输入（$x_1, x_2, x_3$）}：客观事实数据。
  \item \textbf{权重（Weights, $w_1, w_2, w_3$）}：\textbf{每个因素在决策中占多大分量（重要程度）}。
\end{enumerate}

\begin{itemize}
  \item 如果你极度看重学区，学区权重 $w_3$ 就会调得非常大；
  \item 如果你不在乎价格，价格权重 $w_2$ 就会接近 0。
\end{itemize}

\begin{enumerate}
  \item \textbf{偏置（Bias, $b$）}：基础门槛/自身倾向性（比如你天生买房意愿就很强，即使条件一般也想买，$b$ 就是正数）。
  \item \textbf{加权求和}：
\end{enumerate}

\[
z = w_1 x_1 + w_2 x_2 + w_3 x_3 + b = \sum_{i=1}^n w_i x_i + b
\]

\subsubsection{(2) 为什么必须有「激活函数（Activation Function）」？}

如果神经元只做 $z = \sum w_i x_i + b$ 这种加减乘除，它在数学上叫\textbf{线性变换}。
线性变换的致命缺陷是：\textbf{无论你把多少个神经元串联在一起，100 层的线性变换在代数上依然完全等价于仅仅只有 1 层！}

\[
W_2(W_1 x + b_1) + b_2 = (W_2 W_1) x + (W_2 b_1 + b_2) = W_{\text{new}} x + b_{\text{new}}
\]

这意味着，没有非线性激活，多层神经网络就丧失了深度的意义，根本无法拟合现实世界中千变万化的曲折规律！

\textbf{激活函数（如 $\text{ReLU}$、$\text{Sigmoid}$、$\tanh$）}就像是一个“非线性开关”：
\begin{itemize}
  \item 它将输入进来的线性加权和，通过曲线进行扭曲与折叠；
  \item 有了激活函数，多层神经网络就能像折纸一样，在多维空间中折叠出极其复杂的曲面，\textbf{理论上可以逼近宇宙中任意复杂的数学函数}（通用近似定理 Universal Approximation Theorem）！
\end{itemize}

\subsubsection{(3) 什么是多层神经网络与“隐藏层（Hidden Layer）”？}

把成百上千个神经元分层排列，前一层的输出作为后一层的输入，就组成了\textbf{多层感知机（MLP / 神经网络）}：

\begin{figure}[htbp]
\centering
\begin{tikzpicture}[scale=0.9, >=stealth,
  nodeIn/.style={rectangle, rounded corners=2pt, draw=secondary, fill=secondary!10, font=\sffamily\scriptsize, inner sep=3pt},
  nodeHid/.style={circle, draw=primary, fill=primary!10, font=\sffamily\scriptsize, inner sep=2pt},
  nodeOut/.style={rectangle, rounded corners=2pt, draw=warnred, fill=warnred!10, font=\sffamily\scriptsize, inner sep=3pt}
]
  \node[font=\sffamily\footnotesize\bfseries, align=center] at (0, 2.3) {输入层\\[-0.1em]\scriptsize (Input)};
  \node[font=\sffamily\footnotesize\bfseries, align=center] at (3.5, 2.3) {隐藏层 1\\[-0.1em]\scriptsize (Hidden 1)};
  \node[font=\sffamily\footnotesize\bfseries, align=center] at (6.5, 2.3) {隐藏层 2\\[-0.1em]\scriptsize (Hidden 2)};
  \node[font=\sffamily\footnotesize\bfseries, align=center] at (10.0, 2.3) {输出层\\[-0.1em]\scriptsize (Output)};
  
  \foreach \i in {1, 2, 3} {
    \node[nodeIn] (in\i) at (0, 2.0 - \i*1.0) {输入特征 $x_\i$};
    \node[nodeHid] (h1\i) at (3.5, 2.0 - \i*1.0) {$h_{1,\i}$};
    \node[nodeHid] (h2\i) at (6.5, 2.0 - \i*1.0) {$h_{2,\i}$};
    \node[nodeOut] (out\i) at (10.0, 2.0 - \i*1.0) {词表 Logits $\hat{y}_\i$};
  }
  
  \foreach \i in {1, 2, 3} {
    \foreach \j in {1, 2, 3} {
      \draw[->, gray!50] (in\i) -- (h1\j);
      \draw[->, gray!50] (h1\i) -- (h2\j);
      \draw[->, gray!50] (h2\i) -- (out\j);
    }
  }
\end{tikzpicture}
\caption{多层感知机（MLP）特征变换与全连接信息流动图}
\end{figure}

\begin{itemize}
  \item \textbf{为什么叫“隐藏层”？}
\end{itemize}

  因为输入是人类给的文字，输出是最终的预测结果，这两头都是看得见摸得着的。而中间这些层，是\textbf{计算机内部自主形成的高维抽象概念表征（Representation）}，人类很难直接用一两个词去定义中间神经元到底在想什么，因此被称为“隐藏状态（Hidden State）”。

\subsubsection{(4) 神经网络是怎么“学习”的？（损失函数与反向传播）}

网络搭好了，但一开始里面的权重旋钮（$W$ 和 $b$）全是\textbf{随机初始化的乱码}。它怎么学会人类知识的？
这个学习过程完全模拟了人类的“犯错-反省-改正”机制：

\begin{figure}[htbp]
\centering
\begin{tikzpicture}[scale=0.92, >=stealth,
  stepBox/.style={rectangle, rounded corners=3pt, draw=primary, fill=primary!6, font=\sffamily\footnotesize, align=center, inner sep=5pt}
]
  \node[stepBox] (s1) at (0, 0) {输入数据 $X$};
  \node[stepBox, fill=secondary!10, draw=secondary] (s2) at (4.0, 0) {前向传播计算\\(网络前向推理)};
  \node[stepBox] (s3) at (8.5, 0) {模型预测值 $\hat{Y}$};
  \node[stepBox, fill=warnred!10, draw=warnred] (s4) at (8.5, -2.0) {计算损失 (Loss)\\与真实标签 $Y$ 对比};
  \node[stepBox, fill=tipgreen!10, draw=tipgreen] (s5) at (4.0, -2.0) {反向传播 (Backward)\\计算各参数梯度 $\nabla_\theta L$};
  \node[stepBox] (s6) at (0, -2.0) {优化器更新参数\\$\theta \leftarrow \theta - \eta \nabla L$};
  
  \draw[->, very thick, primary] (s1) -- (s2);
  \draw[->, very thick, primary] (s2) -- (s3);
  \draw[->, very thick, warnred] (s3) -- (s4);
  \draw[->, very thick, tipgreen] (s4) -- (s5);
  \draw[->, very thick, tipgreen] (s5) -- (s6);
  \draw[->, thick, dashed, primary] (s6) -- (s1) node[midway, left, font=\scriptsize\sffamily] {迭代下一轮};
\end{tikzpicture}
\caption{深度学习模型训练闭环：前向计算、损失衡量与反向梯度传播}
\end{figure}

\begin{enumerate}
  \item \textbf{第一步（前向瞎猜）}：给模型输入“太阳从东边”，模型一顿乱算，吐出下一个字是“西瓜”。
  \item \textbf{第二步（算差距 Loss）}：损失函数发现标准答案是“升”，预测概率和事实差了十万八千里，给出一个很大的惩罚分数（Loss）。
  \item \textbf{第三步（反向算责任）}：利用高等数学的\textbf{链式求导法则（Chain Rule）}，从后向前计算输出误差对每一个权重参数的偏导数（称为\textbf{梯度 Gradient}）。梯度告诉我们：“如果把第 354 号旋钮稍微拧大一点，总误差就会减小！”
  \item \textbf{第四步（更新旋钮）}：按照梯度的相反方向，微调所有的权重参数（这个过程叫\textbf{梯度下降}）。
  \item \textbf{重复数十亿次}：当把互联网上几万亿个词喂给模型，历经成百上千卡 GPU 几个月的日夜调整，所有的旋钮都达到了极致精密的和谐状态——大模型就此诞生！
\end{enumerate}

\vspace{0.8em}\hrule\vspace{0.8em}

\subsection{为什么传统神经网络处理文本会崩溃？循环神经网络（RNN）的诞生}

搞懂了基础神经网络，我们立刻面临一个致命挑战：
\textbf{传统的神经网络，输入大小必须是固定的！}
\begin{itemize}
  \item 比如设计一个能输入 3 个词的网络，它就只能吃 3 个词；
  \item 但人类说话有长有短，有的问句只有 2 个字（“在吗？”），有的文章长达 1 万字！
  \item 怎么让同一个神经网络，既能读短句，又能读长文呢？
\end{itemize}

科学家为此设计了\textbf{循环神经网络（Recurrent Neural Network, RNN）}。
RNN 的核心构想是：\textbf{网络结构固定不变，但配一个随时更新的“备忘录小本子（隐藏状态 $h_t$）”！}

\begin{figure}[htbp]
\centering
\begin{tikzpicture}[scale=0.9, >=stealth,
  cell/.style={rectangle, rounded corners=3pt, draw=primary, fill=primary!10, font=\sffamily\footnotesize\bfseries, minimum width=2.4cm, minimum height=1.0cm},
  io/.style={rectangle, rounded corners=2pt, draw=secondary, fill=secondary!8, font=\sffamily\scriptsize}
]
  \node[io] (x1) at (0, 0) {$x_1$ (\texttt{"我"})};
  \node[cell] (rnn1) at (0, 1.8) {RNN 单元};
  \node[io, fill=tipgreen!10, draw=tipgreen] (h1) at (0, 3.6) {隐状态 $h_1$};
  
  \node[io] (x2) at (4.0, 0) {$x_2$ (\texttt{"喜欢"})};
  \node[cell] (rnn2) at (4.0, 1.8) {RNN 单元};
  \node[io, fill=tipgreen!10, draw=tipgreen] (h2) at (4.0, 3.6) {隐状态 $h_2$};
  
  \node[io] (x3) at (8.0, 0) {$x_3$ (\texttt{"大模型"})};
  \node[cell] (rnn3) at (8.0, 1.8) {RNN 单元};
  \node[io, fill=tipgreen!10, draw=tipgreen] (h3) at (8.0, 3.6) {隐状态 $h_3$};
  
  \draw[->, thick] (x1) -- (rnn1);
  \draw[->, thick] (rnn1) -- (h1);
  \draw[->, thick] (x2) -- (rnn2);
  \draw[->, thick] (rnn2) -- (h2);
  \draw[->, thick] (x3) -- (rnn3);
  \draw[->, thick] (rnn3) -- (h3);
  
  \draw[->, very thick, warnred] (rnn1) -- (rnn2) node[midway, above, font=\tiny\sffamily] {时序记忆传递};
  \draw[->, very thick, warnred] (rnn2) -- (rnn3) node[midway, above, font=\tiny\sffamily] {语境累加};
\end{tikzpicture}
\caption{循环神经网络（RNN）沿时序展开的信息流动机制}
\end{figure}

每个时间步 $t$，同一个 RNN 网络只接收两样东西：
\begin{enumerate}
  \item \textbf{上一步传过来的备忘录} $h_{t-1}$
  \item \textbf{当前进来的这 1 个新词} $x_t$
\end{enumerate}

通过矩阵乘法与激活函数更新备忘录：

\[
h_t = \tanh(W_h h_{t-1} + W_x x_t + b)
\]

其中：
\begin{itemize}
  \item $W_x$ 是针对新词的权重矩阵；
  \item $W_h$ 是针对历史备忘录的权重矩阵；
  \item $\tanh$ 是激活函数，把所有数字压缩在 $[-1, 1]$ 之间，防止数值无休止放大。
\end{itemize}

\vspace{0.8em}\hrule\vspace{0.8em}

\subsection{为什么强大的 RNN 和 LSTM 最终被历史无情淘汰？}

尽管 LSTM 和 GRU 引入了复杂的门控机制，但在大模型时代，RNN 撞上了一堵无法逾越的物理高墙：

\begin{enumerate}
  \item \textbf{信息串行的「传话游戏」缺陷（长程遗忘与梯度消失）}：
\end{enumerate}

\begin{itemize}
  \item 文本必须一个词接一个词地顺序传递。如果一句话长达 1000 个词，第 1 个词的信息必须经历 1000 次矩阵乘法才能到达末尾。
  \item 在数学上，多次连乘会导致浮点数的\textbf{梯度消失（数值缩小到 0）}或\textbf{梯度爆炸（数值溢出为 NaN）}。哪怕最先进的 LSTM，在超过 100\textasciitilde{}200 个词后，对前文的记忆也基本荡然无存。
\end{itemize}

\begin{enumerate}
  \item \textbf{算力杀手：完全无法在 GPU 上并行训练！}
\end{enumerate}

\begin{itemize}
  \item 现代 GPU 内部有上万个计算核心，GPU 最强大的地方在于\textbf{同一毫秒同时做数千个矩阵相乘}。
  \item 但是 RNN 要求：\textbf{必须先算完第 1 个词，拿到 $h_1$ 之后，才能算第 2 个词；算完第 2 个词才能算第 3 个词！}
  \item 意味着 99\% 的 GPU 核心只能闲置排队等待！训练速度慢如蜗牛，根本无法利用海量互联网数据进行缩放（Scale）。
\end{itemize}

\vspace{0.8em}\hrule\vspace{0.8em}

\section{[加速] 五、第四代：Transformer 的横空出世（2017 至今）}

2017 年，Google 发表了震动整个 AI 领域的著名论文：
\textbf{《Attention Is All You Need》（注意力就是你所需要的一切）}。

Transformer 做出了一个在当时看似离经叛道、甚至疯狂的决定：
\begin{noteBox}[核心导读与背景]
\textbf{彻底抛弃所有的循环（RNN）和时间步串行连接！一个都不留！}
\end{noteBox}

\subsection{Transformer 为什么能一统天下？}

\begin{enumerate}
  \item \textbf{任意两词之间距离永远是 $O(1)$}：
\end{enumerate}

\begin{itemize}
  \item 不管一篇文章有 10 个词还是 10 万个词，Transformer 的\textbf{自注意力机制（Self-Attention）}让文章中任意一个词与另一个词之间，直接通过矩阵点积相连。
  \item 第 1 个词和第 10000 个词可以直接打分产生关联，不再有任何“传话游戏”的信息衰减！
\end{itemize}

\begin{enumerate}
  \item \textbf{天然契合 GPU 的暴力并行计算}：
\end{enumerate}

\begin{itemize}
  \item 一整篇文章输入进来，所有词的 Q、K、V 投影矩阵\textbf{在一瞬间同时相乘完成}！
  \item GPU 的成千上万个核心全负荷拉满运转，训练吞吐量相比 RNN 提升了几个数量级。
\end{itemize}

\begin{enumerate}
  \item \textbf{奠定了如今大模型（LLM）的基石}：
\end{enumerate}

\begin{itemize}
  \item 正是因为能在海量 GPU 集群上并行计算，才使得人类能够把整个维基百科、GitHub 源码、电子书全部喂给模型，催生了 GPT-3、GPT-4、Llama、Qwen 等如今震撼世界的超级大模型。
\end{itemize}

\vspace{0.8em}\hrule\vspace{0.8em}

\section{[目标] 总结与下一章预告}

回顾人类探寻机器理解语言的 70 年历程：
\begin{enumerate}
  \item \textbf{人工规则}：因语言的无限歧义性与复杂性而崩溃；
  \item \textbf{统计 N-gram}：因缺乏语义理解与上下文数据稀疏而停滞；
  \item \textbf{循环神经网络 RNN}：引入连续词向量，但因信息串行传递无法长久记忆，且无法利用现代 GPU 并行算力而遭淘汰；
  \item \textbf{Transformer}：以完全并行的自注意力机制一统天下，成为当今所有大模型的唯一通用底座。
\end{enumerate}

在下一篇中，我们将\textbf{从最基础的初等数学起步}，推导并理解：
\begin{itemize}
  \item 什么是向量与空间变换？
  \item 两个词向量为什么做「点积」就能代表语义相似度？
  \item 深度学习中最常用的归一化函数 Softmax 是怎么来的？（它的导数为什么那么优雅，在第 3 篇推导）
\end{itemize}

\begin{warningBox}[避坑指南与核心警示]
{}[重点] 本篇为了控制篇幅，跳过了两段重要的历史：\textbf{2003 年第一个神经语言模型}，以及 \textbf{2014 年 Seq2Seq 的信息瓶颈如何催生了注意力机制}。另外本篇 §4.2 讲神经网络如何"学习"只是比喻，严格的数学在第 3 篇。它们分别在：
- \textbf{第 3 篇：导数、链式法则、反向传播与交叉熵}
- \textbf{第 5 篇：从原版 Transformer 到 Llama / DeepSeek} §1
\end{warningBox}

\vspace{0.8em}\hrule\vspace{0.8em}

\textbf{上一篇}：\textbf{第 0 篇（导读）：AI 深度学习全景族谱与零基础入门指南} ｜ \textbf{下一篇}：\textbf{第 2 篇：零基础必备数学直觉：向量、空间变换、点积与 Softmax 深度推导} ｜ \textbf{总路线}：\textbf{LEARNING\_PATH.md}

\chapter{零基础必备数学直觉：向量、空间变换、点积与 Softmax 深度推导}

\begin{tipBox}[核心直觉与思考]
\textbf{阅读前须知}：许多同学被 AI 算法劝退，往往不是因为逻辑理解不了，而是被一堆看似高深莫测的数学符号吓退。本篇将把大模型涉及的所有核心数学工具（向量、矩阵乘法、点积、Softmax）用最通俗的几何直觉和一步步代数推导彻底讲透。
\end{tipBox}

\vspace{0.8em}\hrule\vspace{0.8em}

\section{[重点] 一、从标量到张量：数据的物理载体}

在代码和论文中，你经常会看到这四个名词：

\begin{figure}[htbp]
\centering
\begin{tikzpicture}[scale=0.92, >=stealth,
  tensor/.style={rectangle, rounded corners=3pt, align=center, minimum width=2.3cm, minimum height=1.1cm, text centered, font=\sffamily\footnotesize, draw=primary, fill=primary!6}
]
  \node[tensor] (T0) at (0,0) {\textbf{标量 (Scalar)}\\[0.2em]\scriptsize 0 维 $\cdot$ 单个数值};
  \node[tensor] (T1) at (3.2,0) {\textbf{向量 (Vector)}\\[0.2em]\scriptsize 1 维 $\cdot$ 一组有序数值};
  \node[tensor] (T2) at (6.4,0) {\textbf{矩阵 (Matrix)}\\[0.2em]\scriptsize 2 维 $\cdot$ 二维平面网格};
  \node[tensor, fill=tipgreen!10, draw=tipgreen] (T3) at (9.8,0) {\textbf{张量 (Tensor)}\\[0.2em]\scriptsize 3 维及以上 $\cdot$ 多维数据体};
  
  \draw[->, very thick, primary] (T0) -- (T1);
  \draw[->, very thick, primary] (T1) -- (T2);
  \draw[->, very thick, primary] (T2) -- (T3);
\end{tikzpicture}
\caption{数据物理维度的层级跃迁（标量 $\to$ 向量 $\to$ 矩阵 $\to$ 张量）}
\end{figure}

\begin{enumerate}
  \item \textbf{标量（Scalar）}：单个数值。
\end{enumerate}

\begin{itemize}
  \item 例如：今天的气温 $25.5$℃。
\end{itemize}

\begin{enumerate}
  \item \textbf{向量（Vector）}：一列按顺序排列的数字，通常表示高维空间里的一个\textbf{坐标点}或一个\textbf{带方向的箭头}。
\end{enumerate}

\begin{itemize}
  \item 例如：描述一个人：$[180\text{cm}, 75\text{kg}, 25\text{岁}]$，这是一个 3 维向量。
  \item 在大模型中，一个词（如“大”）会被表示为一个包含 4096 个数字的向量，这 4096 个数字就是该词在人类语义空间中的精细坐标。
\end{itemize}

\begin{enumerate}
  \item \textbf{矩阵（Matrix）}：由行和列构成的二维表格。
\end{enumerate}

\begin{itemize}
  \item 形状写作 $[M, N]$，表示有 $M$ 行、$N$ 列。
  \item 矩阵在大模型里通常有两种角色：
  \item \textbf{批量数据}：比如一句话包含 10 个词，每个词是 4096 维，整句话就是 $[10, 4096]$ 的矩阵；
  \item \textbf{空间变换器（权重参数）}：用来把向量从一个空间变换到另一个空间。
\end{itemize}

\begin{enumerate}
  \item \textbf{张量（Tensor）}：超过二维的多维数组。
\end{enumerate}

\begin{itemize}
  \item 例如在推理系统中，常见形状 \texttt{[Batch, Seq\_Len, Hidden\_Dim]}：
  \item \texttt{Batch = 4}（同时处理 4 个用户的提问）
  \item \texttt{Seq\_Len = 512}（每句话有 512 个词）
  \item \texttt{Hidden\_Dim = 4096}（每个词的维度）
  \item 这是一个三维张量，形状为 $[4, 512, 4096]$。
\end{itemize}

\vspace{0.8em}\hrule\vspace{0.8em}

\section{[几何] 二、点积（Dot Product）：大模型如何计算两个词的关联度？}

在自注意力机制中，最重要的核心计算就是\textbf{点积}。为什么计算两个向量的点积，就能代表两个词的语义相关性？

\subsection{代数定义（怎么算）}

假设有两个 $n$ 维向量 $\vec{a} = [a_1, a_2, \dots, a_n]$ 和 $\vec{b} = [b_1, b_2, \dots, b_n]$，它们的点积等于\textbf{对应位置元素相乘再全部相加}：

\[
\vec{a} \cdot \vec{b} = a_1 b_1 + a_2 b_2 + \dots + a_n b_n = \sum_{i=1}^n a_i b_i
\]

\begin{noteBox}[核心导读与背景]
\textbf{例题}：$\vec{a} = [1, 2]$, $\vec{b} = [3, 4]$
点积 $\vec{a} \cdot \vec{b} = 1 \times 3 + 2 \times 4 = 3 + 8 = 11$。
\end{noteBox}

\subsection{几何意义与余弦定理（为什么代表相似度？）}

在欧氏空间几何学中，点积还有一个极其重要的等价公式：

\[
\vec{a} \cdot \vec{b} = |\vec{a}| |\vec{b}| \cos(\theta)
\]

其中：
\begin{itemize}
  \item $|\vec{a}|$ 是向量 $\vec{a}$ 的长度（模长）；
  \item $|\vec{b}|$ 是向量 $\vec{b}$ 的长度；
  \item $\theta$ 是两个向量之间的夹角（从 $0^\circ$ 到 $180^\circ$）；
  \item $\cos(\theta)$ 是夹角的余弦值。
\end{itemize}

\begin{figure}[htbp]
\centering
\begin{tikzpicture}[scale=1.2, >=stealth]
  \coordinate (O) at (0,0);
  \coordinate (A) at (3.5,0);
  \coordinate (B) at (2.8,2.2);
  
  \draw[->, very thick, primary] (O) -- (A) node[right, font=\small\sffamily\bfseries] {向量 $\vec{a}$};
  \draw[->, very thick, secondary] (O) -- (B) node[above right, font=\small\sffamily\bfseries] {向量 $\vec{b}$};
  
  \draw[thick, warnred, ->] (0.8,0) arc (0:38:0.8) node[midway, right, font=\small\sffamily] {$\theta$};
  
  \node[font=\footnotesize\sffamily, align=left, color=darktext] at (4.2, 1.2) {
    $\vec{a} \cdot \vec{b} = |\vec{a}| |\vec{b}| \cos\theta$\\
    $\theta \to 0^\circ \implies \cos\theta \to 1$ (最相似)\\
    $\theta = 90^\circ \implies \cos\theta = 0$ (正交无关)
  };
\end{tikzpicture}
\caption{高维向量内积与夹角余弦相似度直观几何解释}
\end{figure}

让我们来看看 $\cos(\theta)$ 随夹角的变化规律：
\begin{enumerate}
  \item \textbf{方向完全相同（$\theta = 0^\circ$）}：
\end{enumerate}

   $\cos(0^\circ) = 1$，点积取得最大正值！\textbf{（表示两个词的含义高度一致、强烈共鸣）}
\begin{enumerate}
  \item \textbf{方向垂直/正交（$\theta = 90^\circ$）}：
\end{enumerate}

   $\cos(90^\circ) = 0$，点积精确等于 0！\textbf{（表示两个词风马牛不相及，毫不相关）}
\begin{enumerate}
  \item \textbf{方向完全相反（$\theta = 180^\circ$）}：
\end{enumerate}

   $\cos(180^\circ) = -1$，点积为负数！\textbf{（表示语义冲突或对立）}

\begin{tipBox}[核心直觉与思考]
{}[思考] \textbf{核心结论}：
在大模型中，Query（查询）向量与 Key（键）向量做点积：
\end{tipBox}

\[
\text{Score} = Q \cdot K^T
\]

\begin{warningBox}[避坑指南与核心警示]
由 $\vec{q} \cdot \vec{k} = |\vec{q}| |\vec{k}| \cos\theta$，它同时由\textbf{方向（夹角）}和\textbf{长度（模长）}决定：模长相近时，夹角越小点积越大、分配的注意力越高；但模长不同时，一个夹角更大却更"长"的 Key 也可能拿到更高的分数。
所以点积是"\textbf{未归一化的}相似度"，并不等于余弦相似度。这一点很重要：第 4 篇会证明，正是因为模长会随维度 $d_k$ 增长，点积的方差才会变大，所以必须除以 $\sqrt{d_k}$。
\end{warningBox}

\vspace{0.8em}\hrule\vspace{0.8em}

\section{三、矩阵乘法：空间的旋转、拉伸与投影}

很多人高中或大学学过矩阵乘法，只记得“左边行乘以右边列”，却不知道它在物理世界中意味着什么。

\subsection{物理直觉：矩阵是一个“空间投影仪”}

当一个向量 $\vec{x}$（比如代表一个词的 4096 维特征）乘以一个权重矩阵 $W$ 时：

\[
\vec{y} = \vec{x} W
\]

\textbf{这绝不仅是一堆数字的加减乘除，而是在把向量 $\vec{x}$ 投射到一个全新的语义子空间中！}

在 Transformer 中：
\begin{itemize}
  \item 相同的输入词向量 $X$，乘以矩阵 $W_q$ 就被投射到了\textbf{“提问空间（Query Space）”}；
  \item 乘以矩阵 $W_k$ 就被投射到了\textbf{“特征被检索空间（Key Space）”}；
  \item 乘以矩阵 $W_v$ 就被投射到了\textbf{“实际语义内容空间（Value Space）”}。
\end{itemize}

这三个矩阵参数，是大模型在预训练过程中学习到的最重要的核心知识载体。

\vspace{0.8em}\hrule\vspace{0.8em}

\section{[推导] 四、Softmax 函数：从任意实数打分到合法的概率分布}

当注意力机制计算完点积后，得到了一组得分（称为 Logits 或 Attention Scores），例如：

\[
z = [2.5, -1.0, 5.0]
\]

这些数字有正有负，求和不等于 1，计算机无法把它们当作概率或者注意力加权权重使用。

我们需要一个函数，将任意实数数组 $z$ 转换为合法的概率分布 $p$：
\begin{itemize}
  \item 条件 1：所有元素必须严格大于 0（$p_i > 0$）；
  \item 条件 2：所有元素的和必须严格等于 1（$\sum_{i} p_i = 1$）。
\end{itemize}

\subsection{为什么不用最简单的“除以总和”？}

你可能会想，为什么不直接取绝对值再除以总和呢？例如：

\[
p_i = \frac{|z_i|}{\sum |z_j|}
\]

\textbf{为什么不行？}
\begin{enumerate}
  \item \textbf{失去符号判别}：$-5.0$ 和 $+5.0$ 会得到相同的权重，但负数代表不相关甚至反对，绝对不能和正数等同。
  \item \textbf{缺乏数学可导性}：绝对值函数在 $x=0$ 处不可导，神经网络无法进行梯度反向传播。
  \item \textbf{无法拉开差距}：简单的线性归一化不能突出高分项的绝对优势。
\end{enumerate}

\subsection{指数函数 $\exp(x) = e^x$ 的登场与推导}

数学家引入了自然常数的指数函数 $e^x$（$e \approx 2.71828$），它拥有极其完美的数学特质：

\begin{figure}[htbp]
\centering
\begin{tikzpicture}[scale=1.0, >=stealth]
  \draw[->, thick, darktext] (-3.5,0) -- (2.5,0) node[right, font=\small\sffamily] {$x$};
  \draw[->, thick, darktext] (0,-0.3) -- (0,3.8) node[above, font=\small\sffamily] {$y = \exp(x)$};
  
  \draw[domain=-3.2:1.3, smooth, variable=\x, very thick, primary] plot ({\x}, {exp(\x)});
  
  \fill[warnred] (0,1) circle (2pt) node[above left, font=\footnotesize\sffamily] {$(0, 1)$};
  \draw[dashed, codegray] (-3.2, 0) -- (2.2, 0);
  
  \node[font=\footnotesize\sffamily, align=left, draw=boxborder, fill=lightbg, rounded corners=2pt, inner sep=4pt] at (-1.5, 2.2) {
    $e^x > 0$ 恒成立 (严格正定)\\
    单调严格递增\\
    导数 $\frac{d}{dx}e^x = e^x$
  };
\end{tikzpicture}
\caption{$\exp(x) = e^x$ 的严格正定性与单调性}
\end{figure}

利用这一性质，对每一项先取指数确保严格为正，再除以所有指数项的总和完成归一化，这就诞生了著名的 \textbf{Softmax 公式}：

\[
p_i = \text{Softmax}(z)_i = \frac{e^{z_i}}{\sum_{j=1}^n e^{z_j}}
\]

\subsubsection{实例计算：}

设 $z = [1.0, 2.0, 3.0]$：
\begin{enumerate}
  \item 分别取指数：
\end{enumerate}

\begin{itemize}
  \item $e^{1.0} \approx 2.718$
  \item $e^{2.0} \approx 7.389$
  \item $e^{3.0} \approx 20.086$
\end{itemize}

\begin{enumerate}
  \item 求和：$\sum = 2.718 + 7.389 + 20.086 = 30.193$
  \item 归一化概率：
\end{enumerate}

\begin{itemize}
  \item $p_1 = 2.718 / 30.193 \approx 9.0\%$
  \item $p_2 = 7.389 / 30.193 \approx 24.5\%$
  \item $p_3 = 20.086 / 30.193 \approx 66.5\%$
  \item 总和：$9.0\% + 24.5\% + 66.5\% = 100\%$！
\end{itemize}

\begin{tipBox}[核心直觉与思考]
{}[思考] \textbf{Softmax 的直观名字含义}：
传统的 $\max(1, 2, 3)$ 是“硬最大值（Hard-max）”，只会输出 3（概率为 $[0, 0, 1]$）。
而 \textbf{Softmax} 是“平滑软化版的最大值”，最大项拿到绝大部分概率（66.5\%），但同时也给次高项保留了一定机会，且整体函数处处光滑可导！
\end{tipBox}

\vspace{0.8em}\hrule\vspace{0.8em}

\section{[防护] 五、工程避坑细节：Softmax 的平移不变性与防止数值溢出}

在工业级大模型推理引擎中，如果直接按照数学公式 \texttt{math.exp(z)} 编写代码，你的服务会在运行几秒钟后彻底崩溃（报错 \texttt{OverflowError: math range error} 或输出 \texttt{NaN}）。

\subsection{为什么会溢出？}

浮点数（如 float32）能表示的最大数字约为 $3.4 \times 10^{38}$。
当大模型某一步算出的未缩放点积较大时（例如 $z_i = 100$），计算 $e^{100} \approx 2.68 \times 10^{43}$，直接突破硬件浮点数表示上限，导致\textbf{指数爆炸溢出（Overflow）}！

\subsection{数学拯救：平移不变性证明}

科学家们发现了一个美妙的代数恒等式：\textbf{给输入向量的所有元素同时减去一个常数 $C$，Softmax 的计算结果完全不发生任何改变！}

\textbf{【数学严格证明】}：
设新的输入为 $z_i' = z_i - C$，代入 Softmax 公式：

\[
\text{Softmax}(z - C)_i = \frac{e^{z_i - C}}{\sum_{j=1}^n e^{z_j - C}}
= \frac{e^{z_i} \cdot e^{-C}}{\sum_{j=1}^n (e^{z_j} \cdot e^{-C})}
\]

提取公因式 $e^{-C}$：

\[
= \frac{e^{z_i} \cdot e^{-C}}{e^{-C} \cdot \sum_{j=1}^n e^{z_j}}
= \frac{e^{z_i}}{\sum_{j=1}^n e^{z_j}} = \text{Softmax}(z)_i \quad \blacksquare
\]

\subsection{工业级标准实现（Safe Softmax）：}

令常数 $C = \max(z)$（整个向量中的最大值）。
此时每一个被计算的数值 $z_i - C \le 0$，其指数 $e^{z_i - C} \in (0, 1]$，最大项为 $e^0 = 1$。
\textbf{无论输入的数字多么大，指数永远被锁死在 0 到 1 之间，彻底根除浮点数溢出 Bug！}

\vspace{0.8em}\hrule\vspace{0.8em}

\section{[目标] 总结与下一章预告}

在本篇中，我们建立了大模型所需的全部底层数学直觉：
\begin{enumerate}
  \item \textbf{向量}：语义空间的几何坐标点；
  \item \textbf{点积}：由\textbf{夹角与模长共同}决定的未归一化相似度（$\vec a\cdot\vec b = |\vec a||\vec b|\cos\theta$，\textbf{不等于}余弦相似度）；
  \item \textbf{矩阵乘法}：将特征投射到不同的语义子空间；
  \item \textbf{Softmax}：利用指数函数的单调性与正值性，将打分安全映射为加权概率。
\end{enumerate}

掌握了这些基石，下一篇（第 3 篇）我们先回答"神经网络到底怎么学习"：导数、反向传播、交叉熵，以及 \textbf{Softmax 的导数为什么那么优雅}。
再下一篇（第 4 篇）直面 Transformer 最核心的王牌公式——它"为什么必须除以 $\sqrt{d_k}$"的论证，正要用到第 3 篇推导的 Softmax 导数：

\[
\text{Attention}(Q, K, V) = \text{softmax}\left(\frac{Q K^T}{\sqrt{d_k}}\right) V
\]

我们将推导：\textbf{为什么必须除以 $\sqrt{d_k}$？如果没有这根小小的平方根号，大模型为什么会彻底崩溃？}

\vspace{0.8em}\hrule\vspace{0.8em}

\textbf{上一篇}：\textbf{第 1 篇：计算机如何读懂语言？从人工规则到大模型演进史} ｜ \textbf{下一篇}：\textbf{第 3 篇：神经网络如何学习：导数、链式法则、反向传播与交叉熵} ｜ \textbf{总路线}：\textbf{LEARNING\_PATH.md}

\chapter{神经网络如何学习：导数、链式法则、反向传播与交叉熵}

\begin{mathBox}[核心数学定理推导]
\textbf{阅读前须知}：第 1 篇用"拧旋钮"的比喻讲了神经网络怎么学习，但没有真正推导；下一篇（第 4 篇）讲 Transformer 时，还要用到 Softmax 的导数 $p_i(\delta_{ij} - p_j)$。本篇把"学习"这件事从数学上彻底讲透：导数到底是什么、为什么梯度指向下山方向、反向传播为什么比别的求导方法快几十亿倍、损失函数为什么偏偏是交叉熵、以及 Sigmoid / RNN 的梯度为什么会消失。

\textbf{前置知识}：会求 $x^2$、$e^x$、$\ln x$ 的导数即可（高中水平）。如果对"导数是什么"本身没有直觉，先看 \href{https://www.3blue1brown.com/topics/calculus}{3Blue1Brown《微积分的本质》} 第 1\textasciitilde{}4 集，再回来。

\textbf{配套实验}：\textbf{\texttt{lab03\_autograd\_from\_scratch/micrograd\_from\_scratch.py}}——用约 350 行纯 Python 实现一个自动求导引擎，并用数值差分验证本篇的每一个公式。
\end{mathBox}

\vspace{0.8em}\hrule\vspace{0.8em}

\section{[重点] 一、导数：函数在一点附近的"最佳线性近似"}

很多人记得导数是"切线斜率"，但在深度学习里更有用的理解是：

\[
f(x + h) \approx f(x) + f'(x) \cdot h \qquad (h \text{ 很小})
\]

\textbf{导数告诉你：输入挪动一点点 $h$，输出大约挪动 $f'(x) \cdot h$。} 它是一个"灵敏度"。

多元函数 $f(x_1, \dots, x_n)$ 同理，每个输入各有一个灵敏度，称为\textbf{偏导数} $\frac{\partial f}{\partial x_i}$，把它们排成向量就是\textbf{梯度}：

\[
\nabla f = \left[\frac{\partial f}{\partial x_1}, \dots, \frac{\partial f}{\partial x_n}\right], \qquad f(\vec{x} + \vec{h}) \approx f(\vec{x}) + \nabla f \cdot \vec{h}
\]

\subsection{为什么梯度的反方向就是"下山最快"的方向？（严格证明）}

设我们沿单位方向 $\vec{u}$（$|\vec{u}| = 1$）走一小步 $\epsilon$，函数变化量约为：

\[
f(\vec{x} + \epsilon\vec{u}) - f(\vec{x}) \approx \epsilon \, (\nabla f \cdot \vec{u})
\]

由第 2 篇的点积几何公式 $\nabla f \cdot \vec{u} = |\nabla f| \, |\vec{u}| \cos\theta = |\nabla f| \cos\theta$：
\begin{itemize}
  \item $\theta = 0$（沿梯度方向）时增长最快，增量 $= \epsilon |\nabla f|$；
  \item $\theta = 180°$（沿\textbf{负梯度}方向）时下降最快，减量 $= \epsilon |\nabla f|$。 $\blacksquare$
\end{itemize}

于是得到\textbf{梯度下降}（Gradient Descent）更新规则，$\eta$ 叫学习率：

\[
\theta \leftarrow \theta - \eta \, \nabla_\theta \mathcal{L}
\]

\begin{tipBox}[核心直觉与思考]
{}[思考] \textbf{为什么学习率不能太大？} 因为上面整个推导依赖"$f(x+h) \approx f(x) + f'(x)h$"这个\textbf{一阶近似}，只在 $h$ 很小时成立。步子一大，近似失效，损失反而可能上升甚至发散。
\end{tipBox}

\vspace{0.8em}\hrule\vspace{0.8em}

\section{[链接] 二、链式法则：误差如何"顺藤摸瓜"找到每个旋钮}

\subsection{单变量链式法则}

若 $y = g(x)$，$L = f(y)$，则：

\[
\frac{dL}{dx} = \frac{dL}{dy} \cdot \frac{dy}{dx}
\]

用"灵敏度"理解：$x$ 动一点 $\rightarrow$ $y$ 动 $\frac{dy}{dx}$ 倍 $\rightarrow$ $L$ 再动 $\frac{dL}{dy}$ 倍，\textbf{灵敏度沿路径相乘}。

\subsection{多变量链式法则：多条路径要相加}

若 $x$ 同时影响了 $y_1$ 和 $y_2$，而 $L$ 依赖 $y_1, y_2$：

\[
\frac{\partial L}{\partial x} = \frac{\partial L}{\partial y_1}\frac{\partial y_1}{\partial x} + \frac{\partial L}{\partial y_2}\frac{\partial y_2}{\partial x}
\]

\textbf{规则只有两条：同一条路径上相乘，不同路径之间相加。} 这就是反向传播的全部数学。

\subsection{手算一个完整例子}

模型 $\hat{y} = w x + b$，损失 $L = (\hat{y} - y)^2$。取 $x = 2, \ w = 3, \ b = 1, \ y = 10$：

\begin{center}
\small
\begin{tabularx}{\textwidth}{p{3cm} p{4.5cm} X}
\toprule
\textbf{步骤} & \textbf{前向（从左到右算值）} & \textbf{反向（从右到左算梯度）} \\
\midrule
$u = w x$ & $u = 6$ & $\frac{\partial L}{\partial w} = \frac{\partial L}{\partial u} \cdot x = -6 \times 2 = \mathbf{-12}$ \\
$\hat{y} = u + b$ & $\hat{y} = 7$ & $\frac{\partial L}{\partial u} = \frac{\partial L}{\partial \hat{y}} \cdot 1 = -6$，$\frac{\partial L}{\partial b} = -6 \times 1 = \mathbf{-6}$ \\
$e = \hat{y} - y$ & $e = -3$ & $\frac{\partial L}{\partial \hat{y}} = \frac{\partial L}{\partial e} \cdot 1 = -6$ \\
$L = e^2$ & $L = 9$ & $\frac{\partial L}{\partial e} = 2e = -6$ \\
\bottomrule
\end{tabularx}
\end{center}

梯度 $\frac{\partial L}{\partial w} = -12 < 0$：说明把 $w$ 调大，损失会下降。取 $\eta = 0.01$，$w \leftarrow 3 - 0.01 \times (-12) = 3.12$。预测值从 7 向 10 靠近了。
\textbf{你可以在配套实验里用代码逐行复现这张表。}

\vspace{0.8em}\hrule\vspace{0.8em}

\section{[加速] 三、反向传播：为什么它能让训练百亿参数成为可能？}

"对每个参数求偏导"有很多种做法，为什么全世界都用反向传播（Backpropagation，本质是\textbf{反向模式自动微分 Reverse-mode AD}）？

\subsection{方案 A：数值差分（最笨）}

\[
\frac{\partial L}{\partial \theta_i} \approx \frac{L(\theta + h e_i) - L(\theta - h e_i)}{2h}
\]

每个参数要跑 2 次前向。一个 7B 模型就要跑 140 亿次前向才能得到\textbf{一次}梯度——完全不可行。（但它是\textbf{检验}梯度实现是否写对的金标准，配套实验里会用到。）

\subsection{方案 B：前向模式（Forward-mode AD）}

从输入往输出推"每个输入对当前节点的灵敏度"。代价与\textbf{输入个数}成正比——对参数有几十亿、输出只有 1 个标量损失的神经网络，依然要几十亿次。

\subsection{方案 C：反向模式（Backpropagation）[OK]}

从唯一的输出 $L$ 出发，反向计算 $\frac{\partial L}{\partial (\text{每个中间节点})}$。每个节点的梯度只需要它"下游"节点的梯度（已算好，直接复用）乘以局部导数，这就是\textbf{动态规划}。
\textbf{一次反向传播就得到全部参数的梯度，代价只是前向的约 2 倍。}

\begin{noteBox}[核心导读与背景]
{}[文献] \textbf{历史}：反向模式自动微分在 1970 年由 Linnainmaa 提出；1986 年 Rumelhart、Hinton、Williams 在《Nature》上把它用于多层神经网络并展示了隐藏层能学到有意义的表示，这才让 MLP 真正可训练（第 0 篇族谱里的"1986 BP 算法"）。
\end{noteBox}

\subsection{一个对推理工程师很重要的推论：训练 \texorpdfstring{$\approx$}{≈} 6N，推理 \texorpdfstring{$\approx$}{≈} 2N}

\begin{itemize}
  \item 一个参数量为 $N$ 的模型，处理 1 个 token 的前向，每个参数参与\textbf{一次乘法 + 一次加法} $\rightarrow$ 约 $2N$ FLOPs。
  \item 反向传播要对"激活"和"权重"各求一次梯度（见下一节的 $G W^T$ 与 $X^T G$），约 $4N$ FLOPs。
  \item 所以训练每个 token 约 $6N$ FLOPs，\textbf{推理每个 token 约 $2N$ FLOPs}（不计注意力的 $O(\text{上下文长度})$ 项）。第 8 篇会用这个数字来估算 Prefill 的耗时。
\end{itemize}

\vspace{0.8em}\hrule\vspace{0.8em}

\section{[推导] 四、矩阵形式的反向传播：只需记住一个公式}

神经网络里 99\% 的计算是矩阵乘法 $Y = XW$（$X: [B, K]$，$W: [K, N]$，$Y: [B, N]$）。
设上游已经传来了 $G = \frac{\partial L}{\partial Y}$（形状同 $Y$，$[B, N]$），求 $\frac{\partial L}{\partial W}$ 与 $\frac{\partial L}{\partial X}$。

\textbf{逐元素推导}：$Y_{ij} = \sum_k X_{ik} W_{kj}$，所以 $W_{kj}$ 影响了第 $j$ 列所有的 $Y_{ij}$（$i = 1..B$，多条路径相加）：

\[
\frac{\partial L}{\partial W_{kj}} = \sum_i \frac{\partial L}{\partial Y_{ij}} \frac{\partial Y_{ij}}{\partial W_{kj}} = \sum_i G_{ij} X_{ik} = (X^T G)_{kj}
\]

同理 $X_{ik}$ 影响了第 $i$ 行所有的 $Y_{ij}$：

\[
\frac{\partial L}{\partial X_{ik}} = \sum_j G_{ij} W_{kj} = (G W^T)_{ik}
\]

\[
\boxed{\ \frac{\partial L}{\partial W} = X^T G, \qquad \frac{\partial L}{\partial X} = G W^T\ }
\]

\begin{noteBox}[核心导读与背景]
{}[工具] \textbf{形状对齐记忆法}：$\frac{\partial L}{\partial W}$ 必须和 $W$ 同形状 $[K, N]$，手上只有 $X: [B, K]$ 和 $G: [B, N]$，唯一能拼出 $[K, N]$ 的写法就是 $X^T G$。

{}[思考] \textbf{一个会在推理里反复出现的事实}：$\frac{\partial L}{\partial W} = X^T G$ 需要前向时的输入 $X$——这就是训练时必须\textbf{把每层激活值存下来}的原因（训练显存大头之一）。推理不做反向，所以激活用完即扔，这也是第 6 篇对比表中"训练显存杀手是激活值，推理显存杀手是 KV Cache"的根源。
\end{noteBox}

\vspace{0.8em}\hrule\vspace{0.8em}

\section{[目标] 五、损失函数为什么是交叉熵？——最大似然估计}

\subsection{语言模型的训练目标}

回到第 1 篇的概率链式法则：一段文本的概率 $P(x_1, \dots, x_T) = \prod_t P(x_t \mid x_{<t})$。
模型 $p_\theta$ 在每个位置输出一个词表上的分布。\textbf{我们希望模型给真实出现过的文本尽可能高的概率}，这就是\textbf{最大似然估计（MLE）}：

\[
\max_\theta \prod_{t=1}^{T} p_\theta(x_t \mid x_{<t})
\]

\subsection{为什么取对数、再取负号？}

\begin{enumerate}
  \item \textbf{连乘 $\rightarrow$ 连加}：几千个小于 1 的概率相乘会\textbf{下溢}成 0（float32 最小正数约 $10^{-45}$），取对数后变成求和，数值稳定；
  \item $\log$ 单调递增，最大化 $\prod p$ 等价于最大化 $\sum \log p$；
  \item 优化器习惯"最小化"，于是加负号：
\end{enumerate}

\[
\mathcal{L}(\theta) = -\frac{1}{T}\sum_{t=1}^{T} \log p_\theta(x_t \mid x_{<t})
\]

这就是\textbf{负对数似然（NLL）}，也等于"真实分布（one-hot）与模型分布之间的\textbf{交叉熵}"：$H(y, p) = -\sum_i y_i \log p_i = -\log p_{\text{正确词}}$。

\subsection{困惑度（Perplexity）：一个有物理含义的指标}

\[
\text{PPL} = \exp(\mathcal{L})
\]

\begin{itemize}
  \item 如果模型对词表 $V$ 个词\textbf{完全均匀地瞎猜}，$p = 1/V$，$\mathcal{L} = \ln V$，$\text{PPL} = V$。
  \item 所以 PPL 可以理解为"\textbf{模型平均在多少个词之间犹豫}"。PPL = 10 意味着模型的不确定性相当于每步从 10 个等可能的词里挑。
\end{itemize}

\begin{noteBox}[核心导读与背景]
{}[工具] \textbf{调试黄金法则}（\textbf{Lab 05} 的 \texttt{train.py} 会真的断言它）：一个刚随机初始化的模型，初始损失应该约等于 $\ln V$。字符级小模型 $V \approx 1500$（\textbf{会随语料增删而变}，因为语料就是本仓库的文档），初始损失应 $\approx \ln 1500 \approx 7.3$。如果远大于这个值，说明初始化让模型"自信地瞎猜"，训练前就有 bug。

{}[链接] \textbf{与推理的联系}：量化（第 11 篇）后要评估精度损失，工业界最常用的指标就是量化前后 PPL 的变化。
\end{noteBox}

\vspace{0.8em}\hrule\vspace{0.8em}

\section{[循环] 六、Softmax 与交叉熵的梯度：本系列最重要的一组推导}

\subsection{Softmax 的雅可比矩阵}

$p_i = \frac{e^{z_i}}{S}$，其中 $S = \sum_k e^{z_k}$。注意 $\frac{\partial S}{\partial z_j} = e^{z_j}$。用商的求导法则 $\left(\frac{u}{v}\right)' = \frac{u'v - uv'}{v^2}$：

\textbf{情况 $i = j$}：

\[
\frac{\partial p_i}{\partial z_i} = \frac{e^{z_i} S - e^{z_i} e^{z_i}}{S^2} = \frac{e^{z_i}}{S} - \left(\frac{e^{z_i}}{S}\right)^2 = p_i - p_i^2 = p_i(1 - p_i)
\]

\textbf{情况 $i \neq j$}（分子 $e^{z_i}$ 不含 $z_j$，导数为 0）：

\[
\frac{\partial p_i}{\partial z_j} = \frac{0 \cdot S - e^{z_i} e^{z_j}}{S^2} = -p_i p_j
\]

合并成一个式子（$\delta_{ij}$ 在 $i=j$ 时为 1，否则为 0）：

\[
\boxed{\ \frac{\partial p_i}{\partial z_j} = p_i(\delta_{ij} - p_j)\ }
\]

第 4 篇解释"为什么必须除以 $\sqrt{d_k}$"时，用的就是这个公式。\textbf{当某个 $p_i \to 1$、其余 $p_j \to 0$ 时，每一项 $p_i(\delta_{ij} - p_j)$ 都 $\to 0$}：$i=j$ 时是 $1 \times (1-1)$，$i \neq j$ 时含因子 $p_j \to 0$ 或 $p_i \to 0$。这就严格说明了"Softmax 饱和 $\rightarrow$ 梯度消失"，第 4 篇会据此推出"必须除以 $\sqrt{d_k}$"。

\subsection{Softmax + 交叉熵：一个美得不可思议的结果}

设正确类别为 $y$，$\mathcal{L} = -\log p_y$。对 logit $z_j$ 求导（链式法则，经过所有 $p_i$ 的路径相加——但 $\mathcal{L}$ 只依赖 $p_y$，所以只有一条路径）：

\[
\frac{\partial \mathcal{L}}{\partial z_j} = -\frac{1}{p_y} \cdot \frac{\partial p_y}{\partial z_j} = -\frac{1}{p_y} \cdot p_y(\delta_{yj} - p_j) = p_j - \delta_{yj}
\]

\[
\boxed{\ \nabla_z \mathcal{L} = \vec{p} - \vec{y}_{\text{one-hot}}\ }
\]

\textbf{梯度就是"模型预测的分布减去真实分布"。} 预测对了多少，梯度就小多少；而且它永远不会因为 Softmax 饱和而消失——只要预测错了（$p_y$ 小），梯度就大。这正是分类与语言模型都用"Softmax + 交叉熵"组合的深层原因。

\subsection{工程细节：永远不要先算 Softmax 再取 log}

$\log(\text{softmax}(z))_y = z_y - \log \sum_k e^{z_k}$。后一项叫 \textbf{LogSumExp}，用第 2 篇的平移技巧稳定计算：

\[
\text{LSE}(z) = m + \log \sum_k e^{z_k - m}, \quad m = \max_k z_k
\]

先算 softmax 再取 log，当某个 $p$ 下溢成 0 时会得到 $\log 0 = -\infty$。PyTorch 的 \texttt{F.cross\_entropy} 直接吃 logits、内部用 LSE，就是这个原因。

\vspace{0.8em}\hrule\vspace{0.8em}

\section{[下降] 七、梯度消失与爆炸：从 Sigmoid、RNN 到残差与 LSTM 的一条暗线}

\subsection{激活函数的导数决定了梯度能走多远}

\begin{center}
\small
\begin{tabularx}{\textwidth}{p{2.5cm} p{3.5cm} X X}
\toprule
\textbf{激活函数} & \textbf{公式} & \textbf{导数} & \textbf{导数上界} \\
\midrule
Sigmoid & $\sigma(x) = \frac{1}{1 + e^{-x}}$ & $\sigma(x)(1 - \sigma(x))$ & $\mathbf{1/4}$（在 $x=0$） \\
Tanh & $\tanh(x)$ & $1 - \tanh^2(x)$ & $1$ \\
ReLU & $\max(0, x)$ & $0$ 或 $1$ & $1$ \\
SiLU / Swish & $x \sigma(x)$ & $\sigma(x)(1 + x(1 - \sigma(x)))$ & 约 $1.1$ \\
\bottomrule
\end{tabularx}
\end{center}

用 Sigmoid 堆 10 层，即使权重完美，梯度也至少被乘上 $(1/4)^{10} \approx 10^{-6}$。\textbf{这就是 2012 年以前深层网络训练不动、AlexNet 改用 ReLU 后一举成功的重要原因之一。} 现代大模型的 SwiGLU 里用的是 SiLU（第 5 篇讲它的来历）。

\subsection{RNN 为什么会忘记（第 1 篇"传话游戏"的数学版）}

$h_t = \tanh(W_h h_{t-1} + W_x x_t)$，对前一步求导：

\[
\frac{\partial h_t}{\partial h_{t-1}} = \text{diag}(\tanh'(\cdot)) \, W_h
\]

从第 $T$ 步传回第 1 步要连乘 $T-1$ 个这样的矩阵：

\[
\frac{\partial h_T}{\partial h_1} = \prod_{t=2}^{T} \text{diag}(\tanh'_t) \, W_h
\]

粗略地说，若 $W_h$ 的最大奇异值 $< 1$，连乘指数级趋于 0（\textbf{梯度消失}，远处的词学不到）；若 $> 1$ 则可能指数级爆炸（Pascanu et al., 2013）。

\subsection{LSTM 与残差连接：同一个思想——"加法通路"}

\begin{itemize}
  \item \textbf{LSTM（1997）}：细胞状态更新为 $c_t = f_t \odot c_{t-1} + i_t \odot g_t$，于是 $\frac{\partial c_t}{\partial c_{t-1}} = \text{diag}(f_t)$。遗忘门 $f_t$ 可以学到接近 1，梯度就能几乎无衰减地沿 $c$ 流过很多步。
  \item \textbf{ResNet / Transformer 残差（2015/2017）}：$x_l = x_{l-1} + F(x_{l-1})$，$\frac{\partial x_l}{\partial x_{l-1}} = I + \frac{\partial F}{\partial x_{l-1}}$（第 4 篇 §5 会在 Transformer 里再推导一次）。
\end{itemize}

\textbf{两者都是把"乘法连乘的通路"换成"加法直通的通路"。} 理解了这一点，你就理解了为什么深度学习史上几次关键突破长得如此相似。

\vspace{0.8em}\hrule\vspace{0.8em}

\section{八、初始化与归一化：与 $\sqrt{d_k}$ 同一个方差推导}

\textbf{独立随机变量之和的方差等于各自方差之和}，所以 $d$ 个独立的单位方差项相加，方差变为 $d$（第 4 篇推导"除以 $\sqrt{d_k}$"时用的也是这一条）。把它用在一层线性变换 $y_j = \sum_{i=1}^{n} w_{ij} x_i$ 上（$w, x$ 独立、零均值）：

\[
\text{Var}(y_j) = n \cdot \text{Var}(w) \cdot \text{Var}(x)
\]

想让信号逐层既不放大也不衰减，就需要 $\text{Var}(w) = 1/n$。这就是 \textbf{Xavier 初始化（2010）} 的核心；配合 ReLU（砍掉一半信号）时要 $\text{Var}(w) = 2/n$，即 \textbf{He 初始化（2015）}。
即使初始化得当，训练过程中分布也会漂移，于是有了 \textbf{BatchNorm（2015）$\rightarrow$ LayerNorm（2016）$\rightarrow$ RMSNorm（2019）}，把每层输入的尺度重新"拉回来"（演进原因见第 5 篇）。

\begin{mathBox}[核心数学定理推导]
{}[反向] \textbf{同一个推导出现了三次}：$\sqrt{d_k}$ 缩放、Xavier/He 初始化、归一化层，本质都是"\textbf{控制方差，让数值待在激活函数和 Softmax 的敏感区间里}"。
\end{mathBox}

\vspace{0.8em}\hrule\vspace{0.8em}

\section{[演练] 九、优化器与"为什么训练比推理贵这么多显存"}

\begin{itemize}
  \item \textbf{SGD}：$\theta \leftarrow \theta - \eta g$。
  \item \textbf{Momentum}：$m \leftarrow \beta m + g$，$\theta \leftarrow \theta - \eta m$（在峡谷形的损失面上抑制来回震荡）。
  \item \textbf{Adam（2014）}：同时维护一阶矩 $m$ 和二阶矩 $v$：
\end{itemize}

\[
m \leftarrow \beta_1 m + (1-\beta_1) g, \quad v \leftarrow \beta_2 v + (1-\beta_2) g^2, \quad \theta \leftarrow \theta - \eta \frac{\hat{m}}{\sqrt{\hat{v}} + \epsilon}
\]

（$\hat{m}, \hat{v}$ 是偏差修正后的值。）相当于给每个参数一个自适应的学习率。

\textbf{显存账本（混合精度 + Adam，每个参数）}：

\begin{center}
\small
\begin{tabularx}{\textwidth}{p{3cm} p{4.5cm} X}
\toprule
\textbf{项目} & \textbf{训练} & \textbf{推理} \\
\midrule
BF16 权重 & 2 B & 2 B \\
BF16 梯度 & 2 B & — \\
FP32 主权重副本 & 4 B & — \\
Adam $m$（FP32） & 4 B & — \\
Adam $v$（FP32） & 4 B & — \\
\textbf{合计} & \textbf{16 B / 参数}（还没算激活） & \textbf{2 B / 参数}（+ KV Cache） \\
\bottomrule
\end{tabularx}
\end{center}

所以 7B 模型训练至少 112 GB，推理只要 14 GB。这就是为什么你在一张 16GB 的 5060 Ti 上\textbf{能}做推理研究、\textbf{不能}从头训练 7B——也是本仓库路线把重心放在推理上的原因。

\vspace{0.8em}\hrule\vspace{0.8em}

\section{[OK] 本篇自测（答不出来就回去重读对应小节）}

\begin{enumerate}
  \item 为什么学习率过大时损失会上升？（提示：一阶近似）
  \item 一个 $[B, K] \times [K, N]$ 的矩阵乘，反向传播需要做哪两个矩阵乘？各自的形状是什么？总 FLOPs 是前向的几倍？
  \item 用 $p_i(\delta_{ij} - p_j)$ 解释：为什么 Softmax 各输入之间的\textbf{差距}越大，梯度越小？（为什么不能说"输入数值越大"？提示：给所有输入同时加 100，雅可比变不变？）
  \item 字符级模型词表 $V = 65$，随机初始化时损失大约应是多少？PPL 呢？
  \item LSTM 的细胞状态和 Transformer 的残差连接，在"梯度流"意义上共享了哪个设计？
  \item 为什么 7B 模型推理只要约 14GB，训练却要 100GB 以上？
\end{enumerate}

\vspace{0.8em}\hrule\vspace{0.8em}

\section{[资料] 外部资源（按推荐顺序）}

\begin{enumerate}
  \item {}\href{https://www.3blue1brown.com/topics/neural-networks}{3Blue1Brown $\cdot$ 神经网络系列} 第 2\textasciitilde{}4 集（梯度下降、反向传播直觉、反向传播微积分）——先建立图像。
  \item {}\href{https://karpathy.ai/zero-to-hero.html}{Andrej Karpathy $\cdot$ The spelled-out intro to neural networks and backpropagation: building micrograd}（Zero to Hero 第 1 集，约 2.5 小时）——\textbf{强烈建议跟着敲}，配套实验就是它的中文注释精简版。
  \item {}\href{https://cs231n.github.io/optimization-2/}{CS231n 反向传播讲义} ——计算图与"门"的局部梯度视角，含矩阵求导的形状技巧。
  \item {}\href{https://explained.ai/matrix-calculus/}{The Matrix Calculus You Need For Deep Learning} ——想把矩阵求导彻底弄懂时看。
  \item {}\href{https://zh.d2l.ai/chapter\_multilayer-perceptrons/backprop.html}{《动手学深度学习》4.7 前向传播、反向传播和计算图}、\href{https://zh.d2l.ai/chapter\_multilayer-perceptrons/numerical-stability-and-init.html}{4.8 数值稳定性和模型初始化}。
  \item 论文（选读）：Rumelhart et al. 1986 \textit{Learning representations by back-propagating errors}；Pascanu et al. 2013 \href{https://arxiv.org/abs/1211.5063}{On the difficulty of training RNNs}；Kingma \& Ba 2014 \href{https://arxiv.org/abs/1412.6980}{Adam}。
\end{enumerate}

\vspace{0.8em}\hrule\vspace{0.8em}

\textbf{上一篇}：\textbf{第 2 篇：零基础必备数学直觉：向量、空间变换、点积与 Softmax 深度推导} ｜ \textbf{下一篇}：\textbf{第 4 篇：Transformer 核心公式彻底推导：每一个符号从何而来} ｜ \textbf{总路线}：\textbf{LEARNING\_PATH.md}

\chapter{Transformer 核心公式彻底推导：每一个符号从何而来}

\begin{mathBox}[核心数学定理推导]
\textbf{阅读前须知}：在这篇教程中，我们将直面大模型皇冠上的明珠——Transformer 的数学核心。我们不会像很多教材那样直接甩给你公式，而是像侦探破案一样，带你一步一步把每一个符号、每一个分母、每一个负号的数学来源推导得清清楚楚。
\end{mathBox}

\vspace{0.8em}\hrule\vspace{0.8em}

\section{[架构] 一、终极核心公式全貌}

在所有大模型（GPT-4、Llama-3、Qwen-2.5）的心脏位置，跳动着的都是这行经典的注意力公式：

\[
\text{Attention}(Q, K, V) = \text{softmax}\left(\frac{Q K^T}{\sqrt{d_k}}\right) V
\]

请盯住这个公式 3 秒钟。它包含 5 个核心要素：
\begin{enumerate}
  \item \textbf{$Q$（Query，查询）}
  \item \textbf{$K$（Key，键）}
  \item \textbf{$V$（Value，值）}
  \item \textbf{$Q K^T$（点积打分矩阵）}
  \item \textbf{分母 $\sqrt{d_k}$（缩放因子，最精妙的数学设计）}
\end{enumerate}

下面我们逐一拆解。

\vspace{0.8em}\hrule\vspace{0.8em}

\section{[观察] 二、为什么需要拆分成 Q、K、V？（搜索引擎直觉）}

很多人初学时都有一个巨大的困惑：\textbf{明明输入的只有一句话的词嵌入 $X$，为什么偏偏要用三个不同的矩阵把 $X$ 变成 $Q$、$K$、$V$？直接拿 $X$ 算不行吗？}

\subsection{类比：YouTube / 百度搜索引擎的底层逻辑}

想象你在 YouTube 搜索框里搜视频：
\begin{enumerate}
  \item \textbf{你的搜索词}：\texttt{"如何学习大模型推理"} —— 这就是 \textbf{Query（查询）}。它代表当前词“迫切想寻找什么信息”。
  \item \textbf{库里每个视频的标题/标签}：\texttt{"大模型实战教程"}, \texttt{"美食烹饪视频"} —— 这就是 \textbf{Key（索引键）}。它代表每个词用来被别人检索的“特征标签”。
  \item \textbf{视频真正的画面和声音}：几十分钟的真实视频内容 —— 这就是 \textbf{Value（真实值）}。它代表每个词真正承载的信息。
\end{enumerate}

\begin{figure}[htbp]
\centering
\begin{tikzpicture}[scale=0.88, >=stealth,
  box/.style={rectangle, rounded corners=3pt, minimum width=2.6cm, minimum height=0.75cm, text centered, font=\sffamily\footnotesize\bfseries, draw=primary, fill=primary!8}
]
  \node[box, fill=secondary!15, draw=secondary] (X) at (0, 0) {输入向量序列 $X$};
  
  \node[box] (Q) at (4.5, 1.6) {Query ($Q = X W_q$)};
  \node[box] (K) at (4.5, 0) {Key ($K = X W_k$)};
  \node[box] (V) at (4.5, -1.6) {Value ($V = X W_v$)};
  
  \node[anchor=west, font=\sffamily\scriptsize, color=darktext] at (6.2, 1.6) {当前词想要探寻什么信息？};
  \node[anchor=west, font=\sffamily\scriptsize, color=darktext] at (6.2, 0) {当前词有什么特征可被检索？};
  \node[anchor=west, font=\sffamily\scriptsize, color=darktext] at (6.2, -1.6) {当前词承载的实际内容是什么？};
  
  \draw[->, thick, primary] (X.east) -- ++(1.2, 0) |- (Q.west);
  \draw[->, thick, primary] (X.east) -- (K.west);
  \draw[->, thick, primary] (X.east) -- ++(1.2, 0) |- (V.west);
\end{tikzpicture}
\caption{自注意力机制三元组 $Q, K, V$ 投影与语义物理意义}
\end{figure}

\subsection{为什么不能直接用 $X \cdot X^T$？}

如果直接拿原始特征自己和自己做点积：
\begin{itemize}
  \item 相似度矩阵就变成了对称矩阵：词 A 对词 B 的注意力，必须强制等于词 B 对词 A 的注意力。
  \item 但在人类语言中，\textbf{关联往往是极其不对称的}！
  \item 例如句子：“小猫跳上了桌子，因为它饿了”。
  \item 当模型读到代词“它”时，“它”需要极其强烈地关注前文的“小猫”（消歧代词）；
  \item 但前面的“小猫”在被输入时，根本不需要关注后面很远的“它”。
  \item \textbf{引入独立的 $W_q$ 和 $W_k$ 投影，彻底打破了对称性束缚，让信息流动的指向性更加灵活、自由！}
\end{itemize}

\vspace{0.8em}\hrule\vspace{0.8em}

\section{[爆炸] 三、【重磅数学推导】：为什么必须除以 $\sqrt{d_k}$？}

这是面试中最高频、最考验底层功底的一道经典推导题：\textbf{论文作者为什么不除以 $d_k$，也不除以 10，偏偏要除以根号 $\sqrt{d_k}$？}

我们来进行完整的\textbf{方差演化数学推导}。

\subsection{假设输入是标准正态分布}

设 Query 向量的一个分量为 $q_i$，Key 向量的一个分量为 $k_i$。
在经过层归一化（LayerNorm）后，模型内部特征通常满足均值为 0、方差为 1 的独立标准正态分布：
\begin{itemize}
  \item 期望（均值）：$E[q_i] = 0, \quad E[k_i] = 0$
  \item 方差：$\text{Var}(q_i) = 1, \quad \text{Var}(k_i) = 1$
  \item 单个维度的长度为 $d_k$（例如 $d_k = 64$ 或 $128$）。
\end{itemize}

\subsection{计算点积 $Z = \vec{q} \cdot \vec{k}$ 的均值与方差}

点积公式展开为：

\[
Z = \sum_{i=1}^{d_k} q_i k_i = q_1 k_1 + q_2 k_2 + \dots + q_{d_k} k_{d_k}
\]

\subsubsection{(1) 计算点积的数学期望 $E[Z]$：}

根据独立变量期望的可加性与可乘性：

\[
E[Z] = \sum_{i=1}^{d_k} E[q_i k_i] = \sum_{i=1}^{d_k} E[q_i] \cdot E[k_i] = \sum_{i=1}^{d_k} (0 \times 0) = 0
\]

均值仍然是 0，没有偏置。

\subsubsection{(2) 计算单项乘积 $q_i k_i$ 的方差 $\text{Var}(q_i k_i)$：}

根据方差定义公式 $\text{Var}(X) = E[X^2] - (E[X])^2$：

\[
\text{Var}(q_i k_i) = E[(q_i k_i)^2] - (E[q_i k_i])^2 = E[q_i^2] E[k_i^2] - 0
\]

因为 $E[q_i^2] = \text{Var}(q_i) + (E[q_i])^2 = 1 + 0 = 1$，同理 $E[k_i^2] = 1$：

\[
\text{Var}(q_i k_i) = 1 \times 1 = 1
\]

\subsubsection{(3) 计算 $d_k$ 项累加后的总方差 $\text{Var}(Z)$：}

由于这 $d_k$ 个维度相互独立，总方差等于各项方差相加：

\[
\text{Var}(Z) = \sum_{i=1}^{d_k} \text{Var}(q_i k_i) = \underbrace{1 + 1 + \dots + 1}_{d_k \text{ 个}} = d_k
\]

\textbf{根据标准差是方差的平方根}：

\[
\sigma(Z) = \sqrt{\text{Var}(Z)} = \sqrt{d_k}
\]

\vspace{0.8em}\hrule\vspace{0.8em}

\subsection{方差爆炸的灾难后果：Softmax 梯度消失}

请看当隐藏维度变大时发生的事情：
\begin{itemize}
  \item 现代大模型每个头的维度 $d_k$ 通常是 \textbf{$128$}。
  \item 点积未缩放前的标准差就是 $\sqrt{128} \approx \mathbf{11.31}$！
  \item 按照正态分布 $3\sigma$ 原则，大部分点积数值会在 $[-34, +34]$ 之间剧烈震荡。
\end{itemize}

如果把 $[+30, -10, 0]$ 送入 Softmax 会发生什么？
\begin{itemize}
  \item $e^{30} \approx 1.06 \times 10^{13}$
  \item $e^{0} = 1$
  \item 最大的一项占据了 $99.99999999999\%$ 的概率，其他所有项几乎全被压为 $0$！
\end{itemize}

\textbf{更致命的是梯度反向传播公式}：
Softmax 的导数公式为：

\[
\frac{\partial p_i}{\partial z_j} = p_i (\delta_{ij} - p_j)
\]

当某个 $p_i \approx 1$ 而其他 $p_j \approx 0$ 时，这个导数项\textbf{直接逼近于 0}！
\textbf{梯度彻底消失，网络所有参数锁死，无法学习任何深层知识！}

\vspace{0.8em}\hrule\vspace{0.8em}

\subsection{救世主：除以 $\sqrt{d_k}$}

根据随机变量的方差性质：常数 $c$ 提出来要平方，即 $\text{Var}(c X) = c^2 \text{Var}(X)$。
如果我们将点积 $Z$ 除以 $\sqrt{d_k}$：

\[
\text{Var}\left(\frac{Z}{\sqrt{d_k}}\right) = \left(\frac{1}{\sqrt{d_k}}\right)^2 \cdot \text{Var}(Z) = \frac{1}{d_k} \cdot d_k = \mathbf{1}
\]

\textbf{完美的代数闭环！}
无论模型的维度 $d_k$ 做得多大（64、128、256），除以 $\sqrt{d_k}$ 之后，点积结果的方差被\textbf{死死压制在 1.0}！
使得输入进 Softmax 的数值始终落在 $[-3, +3]$ 的黄金活跃区间，Softmax 既能分出高下，又保持了充沛的梯度流动！

\vspace{0.8em}\hrule\vspace{0.8em}

\section{[禁用] 四、因果掩码（Causal Mask）：为什么必须填充 $-\infty$？}

在生成文本时，第 1 个词绝对不能看到第 2 个词的内容。
我们通过一个掩码矩阵 $M$ 与原始点积得分相加：

\[
S_{\text{masked}} = \frac{Q K^T}{\sqrt{d_k}} + M
\]

其中上三角区域（代表未来时刻）全部被填入 $-\infty$（代码中通常用 $-10000.0$ 或 \texttt{-1e9}）：

\[
M = \begin{bmatrix}
0 & -\infty & -\infty \\
0 & 0 & -\infty \\
0 & 0 & 0
\end{bmatrix}
\]

\subsection{数学逻辑：}

根据极限：

\[
\lim_{x \to -\infty} e^x = 0
\]

在计算第 1 行的 Softmax 时：

\[
p_{1, 2} = \frac{e^{S_{1, 2} + (-\infty)}}{\sum e^{\dots}} = \frac{0}{\sum e^{\dots}} = \mathbf{0}
\]

\textbf{加法叠加负无穷，经指数变换后精确转化为乘法意义上的 0 权重，彻底斩断对未来信息的窥探！}

\vspace{0.8em}\hrule\vspace{0.8em}

\section{[路径] 五、残差连接（Residual Connection）：信息的高速公路}

在 Transformer 的每一层后面，都有这样一行代码：

\[
X_{\text{output}} = X_{\text{input}} + \text{SubLayer}(X_{\text{input}})
\]

\subsection{为什么不能直接 $X_{\text{output}} = \text{SubLayer}(X_{\text{input}})$？}

假设一个大模型有 80 层（如 Llama-3-70B）：
在反向传播计算梯度时，根据微积分链式法则：

\[
\frac{\partial \mathcal{L}}{\partial X_0} = \frac{\partial \mathcal{L}}{\partial X_{80}} \times \prod_{l=1}^{80} \frac{\partial X_l}{\partial X_{l-1}}
\]

80 个矩阵导数连乘，只要每一层的雅可比矩阵范数稍微小于 1（例如 0.9），$0.9^{80} \approx 0.0002$，底层的参数梯度瞬间衰减归零，大模型根本训练不起来！

\subsection{加入残差后的神奇变化：}

因为 $X_l = X_{l-1} + F(X_{l-1})$，对 $X_{l-1}$ 求导：

\[
\frac{\partial X_l}{\partial X_{l-1}} = \mathbf{I} + \frac{\partial F(X_{l-1})}{\partial X_{l-1}}
\]

展开到前向损失：

\[
\frac{\partial \mathcal{L}}{\partial X_{l-1}} = \frac{\partial \mathcal{L}}{\partial X_l} \left( \mathbf{I} + \frac{\partial F}{\partial X_{l-1}} \right) = \frac{\partial \mathcal{L}}{\partial X_l} + \frac{\partial \mathcal{L}}{\partial X_l} \frac{\partial F}{\partial X_{l-1}}
\]

看到了吗？式子中永远出现了一个干净纯粹的单位矩阵 $\mathbf{I}$！
\textbf{它宛如一条畅通无阻的高速公路，无论模型叠到 100 层还是 1000 层，梯度都能沿着这条高速路无损地一路冲回第一层！}

\vspace{0.8em}\hrule\vspace{0.8em}

\section{[循环] 六、RoPE（旋转位置编码）：复平面上的几何芭蕾}

最后我们来看现代开源模型（LLaMA/Qwen/DeepSeek）一致采用的 \textbf{RoPE（Rotary Position Embedding）}。

\subsection{核心目标}

我们希望找到一个变换函数 $R(x, m)$，给向量 $x$ 注入位置序号 $m$。
并且满足一个梦幻般的特性：\textbf{两个向量注入位置 $m$ 和 $n$ 后的点积，只和它们的相对距离 $(m - n)$ 有关！}

\[
\langle R(q, m), R(k, n) \rangle = g(q, k, m - n)
\]

\subsection{二维复平面的几何推导}

在复数几何中，将一个二维向量 $(x_0, x_1)$ 视为复数 $z = x_0 + i x_1$。
根据欧拉公式，旋转角度 $\theta$ 相当于乘以 $e^{i \theta}$：

\[
z_{\text{rot}} = z \cdot e^{i \theta} = (x_0 + i x_1)(\cos\theta + i\sin\theta) = (x_0\cos\theta - x_1\sin\theta) + i(x_0\sin\theta + x_1\cos\theta)
\]

转换为实数矩阵乘法形式：

\[
R_{\theta} = \begin{bmatrix}
\cos\theta & -\sin\theta \\
\sin\theta & \cos\theta
\end{bmatrix}
\]

在位置 $m$ 处，旋转角度为 $m \theta$：

\[
q_m = R_{m \theta} q
\]

在位置 $n$ 处，旋转角度为 $n \theta$：

\[
k_n = R_{n \theta} k
\]

\subsection{为什么内积只取决于相对距离？}

让我们计算旋转后的内积：

\[
\langle q_m, k_n \rangle = (R_{m \theta} q)^T (R_{n \theta} k) = q^T (R_{m \theta}^T R_{n \theta}) k
\]

注意旋转矩阵的正交性质：逆时针旋转 $m\theta$ 的转置矩阵，就等于顺时针旋转，即 $R_{m \theta}^T = R_{-m \theta}$！
根据旋转矩阵的相乘性质（先转 $-m\theta$ 再转 $n\theta$）：

\[
R_{m \theta}^T R_{n \theta} = R_{-m \theta} R_{n \theta} = R_{(n - m) \theta}
\]

代入内积式：

\[
\langle q_m, k_n \rangle = q^T R_{(n - m) \theta} k = g(q, k, n - m) \quad \blacksquare
\]

\textbf{数学证明完毕！}
通过在二维平面上纯粹旋转一个角度 $m\theta$：
\begin{enumerate}
  \item 向量本身的长度（模长）没有被拉伸或破坏；
  \item 两个词点积时，两者的绝对位置 $m$ 和 $n$ 自动相消，\textbf{唯独留下了相对距离 $n - m$}！
  \item 当高维空间拆成许多个 2D 坐标对，每个对使用不同的旋转频率 $\theta_i = 10000^{-2(i-1)/d}$ 时，会出现一种\textbf{有条件的}"长程衰减"——这一点经常被说错，要精确表述：
\end{enumerate}

\begin{itemize}
  \item 把第 $i$ 对分量的贡献写成 $A_i \cos(\Delta\theta_i) + B_i \sin(\Delta\theta_i)$（$\Delta$ 是相对距离）。如果 $q$ 与 $k$ \textbf{是相关的}（语义上"匹配"，$A_i$ 平均为正），那么点积的\textbf{期望}是一堆不同频率余弦的正权重叠加，$\Delta$ 变大后各频率相位错开，\textbf{期望值随距离衰减}：同样相似的两个词，离得越远得到的注意力加成越小。
  \item 但如果 $q$ 与 $k$ \textbf{不相关}，点积的分布与距离\textbf{完全无关}（旋转不改变各向同性随机向量的分布），没有任何衰减。对某一对具体的向量，点积随距离剧烈震荡，远处的值完全可能比近处更大。
  \item RoFormer 论文 §3.4.3 证明的是点积的一个\textbf{上界}随距离衰减，并不是点积本身一定变小。
  \item 所以 RoPE 的衰减是"对匹配信号的软性偏置"，不是硬性的距离惩罚。长上下文扩展（PI / NTK / YaRN，第 5 篇 §5.4）要解决的是另一个问题：超过训练长度后，低频维度会转到训练时从未见过的角度（分布外）。
\end{itemize}

\vspace{0.8em}\hrule\vspace{0.8em}

\section{[目标] 总结与下一章预告}

在本篇中，我们逐一完成了 4 大核心公式的严格数学推导：
\begin{enumerate}
  \item \textbf{$Q, K, V$ 拆分}：打破对称性，建立检索信息流；
  \item \textbf{除以 $\sqrt{d_k}$}：证明了独立变量点积方差随维度线性增长，通过平方根归一化彻底拯救了 Softmax 的梯度消失；
  \item \textbf{因果掩码 $-\infty$}：利用 $e^{-\infty}=0$ 斩断未来信息泄露；
  \item \textbf{残差连接}：通过常数导数高速公路根治深度退化；
  \item \textbf{RoPE 旋转矩阵}：通过二维复平面旋转证明了相对位置内积恒等性。
\end{enumerate}

至此，Transformer 的核心公式你已经推导了一遍！

\begin{tipBox}[核心直觉与思考]
{}[重点] \textbf{但注意}：本篇讲的是 2017 年原版论文的公式。现代大模型（Llama、Qwen、DeepSeek）在此基础上做了许多改动——Pre-LN、RMSNorm、SwiGLU、GQA、MLA、MoE——每一处都有明确的历史动机，见 \textbf{第 5 篇}。
想亲手验证本篇的 RoPE 性质（包括"衰减只发生在 q、k 相关时的期望值上，不相关时完全不衰减"这一容易被误解的点），运行 \texttt{lab04\_transformer\_microscope/02\_rope\_and\_attention\_microscope.py}。
\end{tipBox}

在下一篇（第 5 篇）中，我们先看这些公式在 2017 年之后被改了哪些地方、为什么改；再下一篇（第 6 篇）跨越到实际运行：\textbf{到底什么是大模型推理（Inference）？为什么它在首字（Prefill）和后续生成（Decode）表现出完全不同的运行规律？KV Cache 在代数上是如何推演出来的？}

\vspace{0.8em}\hrule\vspace{0.8em}

\textbf{上一篇}：\textbf{第 3 篇：神经网络如何学习：导数、链式法则、反向传播与交叉熵} ｜ \textbf{下一篇}：\textbf{第 5 篇：从原版 Transformer 到 Llama / DeepSeek：每一处改动的历史动机} ｜ \textbf{总路线}：\textbf{LEARNING\_PATH.md}

\chapter{从原版 Transformer 到 Llama / DeepSeek：每一处改动的历史动机}

\begin{mathBox}[核心数学定理推导]
\textbf{阅读前须知}：第 1 篇讲到"2017 年 Transformer 横空出世"就结束了，第 4 篇讲的是 Transformer 的公式，Lab 04 里的代码却已经是 \textbf{RMSNorm + RoPE + SwiGLU}——这些东西 2017 年的原版 Transformer 里一个都没有。它们是怎么来的？本篇补上两段缺失的历史：
1. \textbf{Transformer 之前}：注意力机制到底是为了解决什么问题被发明的？（它并不是 2017 年凭空出现的）
2. \textbf{Transformer 之后}：从原版到 Llama-3、Qwen2.5、DeepSeek-V3，每一处改动\textbf{被什么问题逼出来}，以及它\textbf{对推理意味着什么}。

\textbf{前置}：第 1、3、4 篇。特别是第 3 篇的"加法通路"和"方差控制"，本篇会反复用到。
\end{mathBox}

\vspace{0.8em}\hrule\vspace{0.8em}

\section{[指南] 一、注意力的真正起源（2003 \texorpdfstring{$\rightarrow$}{->} 2017）}

\subsection{2003：第一个神经语言模型（Bengio NPLM）}

第 1 篇说 N-gram 的死穴是"苹果"和"梨"毫无关系。2003 年 Bengio 等人提出\textbf{神经概率语言模型}：

\[
P(w_t \mid w_{t-n+1}, \dots, w_{t-1}) = \text{softmax}\big(b + W x + U \tanh(d + H x)\big), \qquad x = [e_{t-n+1}; \dots; e_{t-1}]
\]

（$x$ 是前 $n-1$ 个词向量的拼接；$U\tanh(d + Hx)$ 是一个单隐藏层 MLP；$Wx$ 是论文特意加入的\textbf{从输入直连输出}的线性项，可以设为 0。）
\begin{itemize}
  \item 每个词先查表得到一个\textbf{可学习的向量} $e$（\textbf{这就是词嵌入的真正起点}，比 Word2Vec 早 10 年）；
  \item 相似的词学到相似的向量，于是"我吃了苹果"的统计经验能\textbf{泛化}到"我吃了梨"。
  \item 但它仍然只看\textbf{固定窗口}的 $n-1$ 个词，并且 softmax 覆盖整个词表，在当时的硬件上训练极慢。
\end{itemize}

之后 2010 年 Mikolov 用 RNN 做语言模型（RNNLM）去掉了固定窗口；2013 年 Word2Vec 把"只学词向量"这件事做到极致地便宜。

\subsection{2014：Seq2Seq 与"信息瓶颈"}

Sutskever 等人用两个 LSTM 做翻译：\textbf{编码器}读完整个源句子，把它压成\textbf{一个固定长度的向量} $c$；\textbf{解码器}从 $c$ 出发逐词生成译文。
问题马上暴露：无论源句子 5 个词还是 50 个词，都要塞进同一个几百维的向量里。Cho 等人（2014）的实验显示\textbf{句子越长，翻译质量下降越明显}。这个固定长度向量就是\textbf{信息瓶颈}。

\subsection{2014：Bahdanau 注意力——"别压缩了，每一步回头看"}

Bahdanau、Cho、Bengio 的想法是：解码器生成第 $t$ 个词时，\textbf{不要只看那个压缩向量，而是回头看编码器每个位置的隐藏状态 $h_j$，按需加权}：

\[
e_{tj} = a(s_{t-1}, h_j), \qquad \alpha_{tj} = \frac{\exp(e_{tj})}{\sum_k \exp(e_{tk})}, \qquad c_t = \sum_j \alpha_{tj} h_j
\]

\textbf{请对照第 4 篇的注意力公式}：

\begin{center}
\small
\begin{tabularx}{\textwidth}{p{3.5cm} X}
\toprule
\textbf{Bahdanau 2014} & \textbf{Transformer 2017} \\
\midrule
解码器当前状态 $s_{t-1}$（"我现在要翻译什么"） & Query \\
编码器各位置 $h_j$（用来算匹配分数） & Key \\
编码器各位置 $h_j$（被加权求和的内容） & Value \\
打分函数 $a(\cdot)$ 是一个小 MLP（加性注意力） & 打分函数是点积 $q \cdot k / \sqrt{d_k}$ \\
\bottomrule
\end{tabularx}
\end{center}

2015 年 Luong 等人把打分函数换成\textbf{点积}（乘性注意力）——更便宜，而且整批计算恰好是一次矩阵乘，GPU 最擅长。\textbf{注意}：2014 年 Bahdanau 已经有了「打分 + softmax 加权求和」这个两段结构，但那时 Key 和 Value 是\textbf{同一个} $h_j$（没有 $W_k/W_v$ 两套投影），Query 就是解码器状态而非可学习的投影。\textbf{Q/K/V 三套独立投影的分工是 2017 年 Transformer 才定型的}——上面表格里的 Q/K/V 是后人回填的框架。

\subsection{2017：Transformer 的真正创新是"只要注意力"}

既然注意力这么有用，而 RNN 又串行、又会遗忘（第 1、3 篇），Vaswani 等人问：\textbf{能不能把 RNN 整个扔掉，只用注意力？}
\begin{itemize}
  \item \textbf{自注意力}：Query、Key、Value 都来自同一个序列自身（而不是解码器看编码器）；
  \item \textbf{位置编码}：扔掉 RNN 后，模型丧失了顺序感（\texttt{lab04\_transformer\_microscope/02\_rope\_and\_attention\_microscope.py} 实验 1 证明了这一点），必须显式注入位置；
  \item \textbf{多头}：多组独立的 $W_q, W_k, W_v$，让不同的头关注不同类型的关系；
  \item \textbf{缩放点积}：除以 $\sqrt{d_k}$（第 4 篇的推导）。
\end{itemize}

原版 Transformer 是\textbf{编码器-解码器}结构，用来做翻译，6 层编码器 + 6 层解码器，$d_{\text{model}} = 512$，FFN 用 ReLU，\textbf{Post-LN}，\textbf{正弦位置编码}。

\vspace{0.8em}\hrule\vspace{0.8em}

\section{二、分词（Tokenization）：为什么"词表大小"会影响推理速度}

\subsection{为什么不按"词"或"字"切？}

\begin{itemize}
  \item \textbf{按词切}：词表爆炸（英文有各种变形，中文没有天然空格），总会遇到没见过的词（OOV, Out-Of-Vocabulary）只能标记为 \texttt{$<$unk$>$}；
  \item \textbf{按字符 / 字节切}：永远不会 OOV，但序列变得很长——注意力是 $O(n^2)$，而且\textbf{Decode 步数 = token 数}，每多一个 token 就要多读一遍全部权重（第 6 篇）。
\end{itemize}

\subsection{BPE：从数据压缩借来的算法}

2016 年 Sennrich 等人把 1994 年的数据压缩算法 \textbf{BPE（Byte Pair Encoding）} 用到翻译上：
\begin{enumerate}
  \item 从单个字符（或字节）出发作为初始词表；
  \item 统计语料里\textbf{最常相邻出现的一对符号}，合并成一个新符号加入词表；
  \item 重复第 2 步，直到词表达到目标大小（例如 32000、128000、151936）。
\end{enumerate}

结果：常见词是一个 token，罕见词被拆成几个子词，\textbf{永远不会 OOV}。GPT-2 进一步在\textbf{字节}上做 BPE（Byte-level BPE），任何 UTF-8 文本都能编码。

\subsection{词表大小是一个推理工程上的权衡}

\begin{center}
\small
\begin{tabularx}{\textwidth}{p{3.5cm} X}
\toprule
\textbf{词表更大} & \textbf{词表更小} \\
\midrule
{}[OK] 同样的文本 token 更少 $\rightarrow$ Decode 步数更少、KV Cache 更小 & {}[OK] Embedding 与 LM Head 更小 \\
{}[X] Embedding 表和 LM Head 更大：LM Head 每个 token 都要算 $d \times V$ 的矩阵乘 & {}[X] 同样文本 token 更多 \\
\bottomrule
\end{tabularx}
\end{center}

例子：Qwen2.5-0.5B 的 $V = 151936$、$d = 896$，\textbf{光是词嵌入矩阵就有 1.36 亿参数，占全模型约 28\%}（Lab 07b 会打印出来）。Llama-2 词表只有 32000，对中文很不友好（一个汉字常被切成多个 token）；Llama-3 扩到 128256、Qwen2.5 小尺寸（0.5B/1.5B/3B）用 151936（且词嵌入与 LM Head 共享权重），大尺寸（7B 及以上）用 152064（不共享），中文同样的内容 token 数显著减少——\textbf{这直接等于中文推理变快、变便宜}。

\begin{noteBox}[核心导读与背景]
{}[链接] 深入：Karpathy 的视频 \textit{Let's build the GPT Tokenizer} 与代码 \href{https://github.com/karpathy/minbpe}{minbpe}。
\end{noteBox}

\vspace{0.8em}\hrule\vspace{0.8em}

\section{[架构] 三、三条路线：为什么最后是"只有解码器"的 GPT 赢了？}

\begin{center}
\small
\begin{tabularx}{\textwidth}{p{2.5cm} p{3.5cm} X X}
\toprule
\textbf{路线} & \textbf{代表} & \textbf{训练目标} & \textbf{能力特点} \\
\midrule
仅编码器 & BERT（2018） & 完形填空（遮住 15\% 的词，双向看上下文去猜） & 理解任务极强，但\textbf{不能自回归生成} \\
编码器-解码器 & 原版 Transformer、T5（2019） & 输入 $\rightarrow$ 输出 & 适合翻译、摘要等"有明确输入输出"的任务 \\
仅解码器 & GPT 系列、Llama、Qwen、DeepSeek & \textbf{预测下一个词} & 什么文本都能当训练数据；生成与理解统一成"续写" \\
\bottomrule
\end{tabularx}
\end{center}

主流的解释（不是定理，而是历史经验）：
\begin{enumerate}
  \item \textbf{目标最简单、数据最多}：任何一段文本都是"预测下一个词"的训练数据，无需标注；
  \item \textbf{一切任务都能写成续写}：GPT-3（2020）展示了\textbf{上下文学习（In-Context Learning）}——在提示里给几个例子，不改参数就能做新任务；
  \item \textbf{规模化最顺}：同样算力下，仅解码器结构简单、易于扩展；
  \item \textbf{对推理友好}：只有一种因果注意力，KV Cache 结构统一、实现简单。
\end{enumerate}

\vspace{0.8em}\hrule\vspace{0.8em}

\section{[增长] 四、规模定律：为什么"推理成本"成了整个行业的主问题}

\begin{itemize}
  \item \textbf{2019 GPT-2（1.5B）}：发现模型够大后，不微调也能做一些任务（zero-shot）。
  \item **2020 Kaplan 等 \textit{Scaling Laws}\textbf{：损失随参数量 $N$、数据量 $D$、算力 $C$ 呈}幂律**下降，例如 $L(N) \propto N^{-0.076}$。只要按规律加大规模，效果就可预测地变好。
  \item \textbf{2020 GPT-3（175B）}：上下文学习能力涌现。
  \item \textbf{2022 Chinchilla（Hoffmann 等）}：在\textbf{固定训练算力}下，最优做法是参数量与训练 token 数同比例增长，约 \textbf{每个参数 20 个 token}。此前的大模型普遍"参数太多、数据太少"。
  \item \textbf{2022 InstructGPT $\rightarrow$ ChatGPT}：用指令微调 + RLHF 让模型"听话"，大模型从研究品变成亿级用户产品。
\end{itemize}

\textbf{关键转折}：Chinchilla 只优化了\textbf{训练}成本。但一个模型训练一次，却要被调用成千上万亿次。2023 年的 Llama 明确把目标改成"\textbf{给定推理预算下效果最好}"——宁可用远超 Chinchilla 比例的数据去\textbf{过度训练一个小模型}（Llama-3-8B 用了 15T token $\approx$ 1870 token/参数，是 Chinchilla 比例 20 的约 93 倍；
早期的 Llama-1-7B 只用了 1T $\approx$ 143 token/参数，约 7 倍）。
\textbf{从这一刻起，推理效率成为模型设计本身的一等公民}：下面第五节的 MQA/GQA/MLA、MoE、词表扩大，全部带着"推理更便宜"的动机。

\vspace{0.8em}\hrule\vspace{0.8em}

\section{[工具] 五、从原版 Transformer 到现代 LLM：逐个改动的"为什么"}

\subsection{Post-LN \texorpdfstring{$\rightarrow$}{->} Pre-LN：把归一化挪出"高速公路"}

\begin{itemize}
  \item \textbf{原版 Post-LN}：$x_{l} = \text{LN}(x_{l-1} + F(x_{l-1}))$。LN 挡在残差路径上，第 3 篇说的"$I$ 直通路径"被每一层的 LN 打断，梯度必须穿过 $L$ 个 LN 的雅可比。深层时训练不稳定，必须用学习率预热（warmup）小心伺候。
  \item \textbf{Pre-LN}：$x_{l} = x_{l-1} + F(\text{LN}(x_{l-1}))$。残差路径上\textbf{什么都没有}，梯度可以原封不动流回第一层。Xiong 等（2020）从理论上分析了这一点；实际上 GPT-2（2019）就已经改成了 Pre-LN。
  \item \textbf{代价}：Pre-LN 的残差流数值会随层数累积变大，所以最后要再加一个 \textbf{Final Norm}（你在 \texttt{lab04\_transformer\_microscope} 里看到的 Step 8）。
\end{itemize}

\subsection{LayerNorm \texorpdfstring{$\rightarrow$}{->} RMSNorm：去掉"减均值"}

\[
\text{LayerNorm}(x) = \frac{x - \mu}{\sqrt{\sigma^2 + \epsilon}} \odot \gamma + \beta \qquad \longrightarrow \qquad \text{RMSNorm}(x) = \frac{x}{\sqrt{\frac{1}{d}\sum_i x_i^2 + \epsilon}} \odot \gamma
\]

Zhang \& Sennrich（2019）发现 LayerNorm 起作用的主要是"\textbf{缩放}"（re-scaling），"\textbf{减均值}"（re-centering）贡献很小。RMSNorm 少一次求均值的归约、没有 $\beta$，效果相当。
\textbf{对推理的意义}：归一化是典型的\textbf{访存受限}小算子（读一遍、写一遍，算得很少），少一次归约就少一次同步，也更容易和前后算子\textbf{融合（fuse）}成一个 kernel（第 8 篇）。

\subsection{ReLU \texorpdfstring{$\rightarrow$}{->} GELU \texorpdfstring{$\rightarrow$}{->} SwiGLU：门控前馈网络}

\begin{itemize}
  \item 原版 FFN：$\text{ReLU}(x W_1) W_2$，隐藏维度 $4d$。
  \item BERT / GPT-2 换成更平滑的 \textbf{GELU}。
  \item Shazeer（2020）系统比较了多种\textbf{门控线性单元（GLU）}变体，SwiGLU 在同等计算量下效果最好：
\end{itemize}

\[
\text{SwiGLU}(x) = \big(\text{SiLU}(x W_{\text{gate}}) \odot x W_{\text{up}}\big) W_{\text{down}}
\]

  门控的直觉：一路计算"内容"，另一路计算"让多少内容通过"，两者\textbf{相乘}，比单纯的逐元素非线性表达力更强。有趣的是，论文作者自己写道：\textit{"我们不解释这些架构为何有效，把它们的成功和其他一切一样，归功于上天的仁慈。"}——深度学习很多改动的"为什么"首先是\textbf{实验结果}，理论解释往往滞后。
\begin{itemize}
  \item \textbf{一个你在配置文件里会看到的数字}：SwiGLU 有 3 个矩阵而不是 2 个。为了保持参数量不变，隐藏维度从 $4d$ 缩成 $\frac{2}{3} \times 4d = \frac{8}{3}d$。Llama-7B 的 $d = 4096$，$\frac{8}{3} \times 4096 \approx 10923$，向上取整到 256 的倍数就是 \textbf{11008}——这就是 Llama 配置里那个"奇怪数字"的来历。
\end{itemize}

\subsection{绝对位置编码 \texorpdfstring{$\rightarrow$}{->} RoPE：为了"外推"与"相对位置"}

\begin{itemize}
  \item 原版用正弦函数的绝对位置编码，加在输入上；GPT-2 改成\textbf{可学习的}绝对位置向量——但训练时最长 1024，就\textbf{根本没有}第 1025 个位置的向量，无法外推。
  \item T5（相对位置偏置）、ALiBi（按距离线性惩罚注意力分数）走向了"相对位置"。
  \item \textbf{RoPE}（苏剑林等，2021）：对 Q、K 按位置旋转，内积只依赖相对距离（第 4 篇已证明），\textbf{没有额外参数}。
\end{itemize}

\textbf{RoPE 对推理的一个关键影响}：K 在\textbf{写入 KV Cache 之前}就已经按它自己的位置旋转好了。之后无论新 token 走到哪个位置，缓存里的 K 都不用再动——只需要旋转新来的 q 和 k。Lab 07b 里你会亲手实现这一点。

\textbf{长上下文扩展}：RoPE 的频率 $\theta_i = b^{-2i/d}$ 中的底数 $b$ 决定了"最慢的那一对分量转多快"。训练长度为 $L$ 时，低频维度只转过了很小的角度；推理时位置一旦超过 $L$，这些维度会转到\textbf{训练中从未见过的角度}（分布外），模型行为随之失控——这才是需要扩展方法的原因（RoPE 对"匹配信号"的期望衰减是另一回事，见第 4 篇 §6）。要把训练时的 4K 上下文扩到 32K、128K，常见做法是\textbf{调整频率}：位置插值（PI, 2023）、NTK-aware 缩放、YaRN（2023）。这也是为什么 Llama-3 把 $b$ 从 10000 调到 500000，Qwen2.5 用 1000000。

\subsection{MHA \texorpdfstring{$\rightarrow$}{->} MQA \texorpdfstring{$\rightarrow$}{->} GQA \texorpdfstring{$\rightarrow$}{->} MLA：一部"为 KV Cache 减肥"的历史}

这一条演进\textbf{完全由推理驱动}，是本篇和第 6 篇之间最直接的桥梁。

\begin{center}
\small
\begin{tabularx}{\textwidth}{l X X X X }
\toprule
\textbf{方案} & \textbf{年份} & \textbf{KV 头数} & \textbf{每 token 每层 KV 存储} & \textbf{动机与代价} \\
\midrule
MHA 多头注意力 & 2017 & $= n_h$ & $2 n_h d_h$ & 原版 \\
\textbf{MQA} 多查询注意力 & 2019 & $1$ & $2 d_h$ & Shazeer 的论文标题就叫 \textit{One Write-Head is All You Need}，\textbf{明确为了解决增量解码的显存带宽瓶颈}。代价：质量下降、训练不稳 \\
\textbf{GQA} 分组查询注意力 & 2023 & $g$（如 8） & $2 g d_h$ & Ainslie 等人的折中：每 $n_h / g$ 个 Q 头共享一组 KV。质量接近 MHA，KV 缩小 $n_h / g$ 倍。Llama-2-70B、Llama-3、Qwen2.5 均采用 \\
\textbf{MLA} 多头潜在注意力 & 2024 & — & $d_c + d_h^R$（如 512 + 64） & DeepSeek-V2：不存 K、V 本身，而是存一个低秩的\textbf{压缩潜向量} $c^{KV}$，用时再"解压"；解压矩阵还能代数上吸收进 $W_q$ 与 $W_o$。RoPE 与低秩压缩不兼容，所以额外存一个 64 维的带位置小 key \\
\bottomrule
\end{tabularx}
\end{center}

用数字感受一下（BF16，每 token，全部层）：
\begin{itemize}
  \item Llama-3-8B（32 层，GQA-8，$d_h = 128$）：$2 \times 32 \times 8 \times 128 \times 2 = 131{,}072$ 字节 = \textbf{128 KB}
  \item DeepSeek-V3（671B 总参数，61 层，MLA）：$(512 + 64) \times 61 \times 2 \approx 70{,}272$ 字节 $\approx$ \textbf{69 KB}
  \item 如果 DeepSeek-V3 用 MHA（128 头 $\times$ 128 维）：$2 \times 128 \times 128 \times 61 \times 2 \approx$ \textbf{3.8 MB}
\end{itemize}

一个 671B 的模型，每 token 的 KV Cache 比 8B 的 Llama-3 还小。\textbf{这就是"模型架构为推理服务"的最佳例证。}

\subsection{稠密 \texorpdfstring{$\rightarrow$}{->} MoE（混合专家）：参数多 \texorpdfstring{$\neq$}{!=} 算得多}

把每层的 FFN 换成 $E$ 个"专家" FFN，再加一个\textbf{路由器}：每个 token 只挑得分最高的 $k$ 个专家计算（例如 $E = 256$ 选 $k = 8$）。
\begin{itemize}
  \item Mixtral 8x7B：总参数约 47B，每 token 激活约 13B；
  \item DeepSeek-V3：总参数 671B，每 token 激活 37B。
\end{itemize}

\textbf{对推理的意义（很反直觉）}：
\begin{itemize}
  \item \textbf{算力}只跟"激活参数"有关 $\rightarrow$ Prefill 很便宜；
  \item \textbf{显存}要装下\textbf{全部}参数 $\rightarrow$ 必须多卡；
  \item \textbf{Decode 带宽}：batch 很小时每步只读被选中专家的权重；但 batch 一大，几乎所有专家都会被某个 token 选中，每一步又要读全部权重——于是出现了\textbf{专家并行（EP）}等专门的部署方案（第 13 篇）。
\end{itemize}

\subsection{其它细节改动}

\begin{itemize}
  \item \textbf{去掉偏置项}：Llama 系列去掉了线性层的 bias（Qwen2 保留了 Q/K/V 投影的 bias，Lab 07b 实现时要注意）。
  \item \textbf{词嵌入与 LM Head 共享权重（tied embeddings）}：小模型里词表矩阵占比极大，共享可省下几亿参数。Qwen2.5-0.5B 就是共享的。
  \item \textbf{滑动窗口注意力}（Mistral-7B, 2023）：每个 token 只看最近 $W$ 个 token，KV Cache 大小封顶。
\end{itemize}

\vspace{0.8em}\hrule\vspace{0.8em}

\section{六、一张总表：同一个"Transformer"，七年里变了什么}

\begin{center}
\small
\begin{tabularx}{\textwidth}{l X X X X X }
\toprule
\textbf{组件} & \textbf{原版 Transformer（2017）} & \textbf{GPT-2/3（2019-20）} & \textbf{Llama-1（2023）} & \textbf{Llama-3 / Qwen2.5（2024）} & \textbf{DeepSeek-V3（2024）} \\
\midrule
结构 & 编码器-解码器 & 仅解码器 & 仅解码器 & 仅解码器 & 仅解码器 \\
归一化位置 & Post-LN & \textbf{Pre-LN} & Pre-LN & Pre-LN & Pre-LN \\
归一化类型 & LayerNorm & LayerNorm & \textbf{RMSNorm} & RMSNorm & RMSNorm \\
FFN 激活 & ReLU & GELU & \textbf{SwiGLU} & SwiGLU & SwiGLU（\textbf{MoE}） \\
位置编码 & 正弦绝对 & 可学习绝对 & \textbf{RoPE} & RoPE（大底数） & RoPE（解耦 + YaRN） \\
注意力 & MHA & MHA & MHA & \textbf{GQA} & \textbf{MLA} \\
词表 & \textasciitilde{}37K（BPE） & 50257（字节 BPE） & 32000 & 128256 / 151936（Qwen2.5 小尺寸）或 152064（7B+） & 129280 \\
\bottomrule
\end{tabularx}
\end{center}

\begin{noteBox}[核心导读与背景]
{}[目标] \textbf{读完本篇你应该能做到}：打开任意一个模型的 \texttt{config.json}（例如 Hugging Face 上 Qwen2.5 的配置），逐个字段说出它是什么、为什么有它、对推理的显存和速度有什么影响。Lab 07b 的第一个任务就是这个。
\end{noteBox}

\vspace{0.8em}\hrule\vspace{0.8em}

\section{[OK] 本篇自测}

\begin{enumerate}
  \item Bahdanau 注意力要解决 Seq2Seq 的什么问题？它和 Transformer 的注意力，Q/K/V 分别对应什么？
  \item 为什么 Transformer 必须有位置编码，而 RNN 不需要？
  \item 为什么 Llama-7B 的 FFN 中间维度是 11008？
  \item Pre-LN 比 Post-LN 好训练，用第 3 篇的"加法通路"解释。
  \item GQA-8 的 70B 模型比 MHA 版本省多少 KV Cache？为什么说 MQA 的动机来自推理而不是训练？
  \item MoE 模型"参数 671B、激活 37B"，在 batch=1 与 batch=256 时，Decode 每一步需要读的权重量分别大约是多少？
  \item Llama-3 为什么要用远超 Chinchilla 最优比例的数据训练 8B 模型？
\end{enumerate}

\vspace{0.8em}\hrule\vspace{0.8em}

\section{[资料] 外部资源}

\textbf{入门可视化}
\begin{itemize}
  \item Jay Alammar：\href{https://jalammar.github.io/illustrated-transformer/}{The Illustrated Transformer}、\href{https://jalammar.github.io/visualizing-neural-machine-translation-mechanics-of-seq2seq-models-with-attention/}{Visualizing A Neural Machine Translation Model (Seq2Seq + Attention)}
  \item Lilian Weng：\href{https://lilianweng.github.io/posts/2018-06-24-attention/}{Attention? Attention!}（从 Bahdanau 到 Transformer 的注意力家谱）
  \item 3Blue1Brown：\href{https://www.3blue1brown.com/topics/neural-networks}{神经网络系列} 第 5\textasciitilde{}7 集（GPT、注意力、MLP 如何存储知识）
\end{itemize}

\textbf{跟着写代码}
\begin{itemize}
  \item Karpathy：\textit{Let's build GPT: from scratch, in code, spelled out} 与 \textit{Let's build the GPT Tokenizer}（\href{https://karpathy.ai/zero-to-hero.html}{Zero to Hero} 第 7、8 集）
  \item Harvard NLP：\href{https://nlp.seas.harvard.edu/annotated-transformer/}{The Annotated Transformer}（原版论文逐段对照代码）
  \item Sebastian Raschka：\href{https://github.com/rasbt/LLMs-from-scratch}{LLMs-from-scratch}（含 GPT $\rightarrow$ Llama $\rightarrow$ Qwen 的逐步改造笔记本）
\end{itemize}

\textbf{中文深度}
\begin{itemize}
  \item 苏剑林（RoPE 作者）博客：\href{https://kexue.fm/archives/8265}{Transformer 升级之路：2、博采众长的旋转式位置编码}
\end{itemize}

\textbf{论文清单（按本篇出现顺序）}
\begin{itemize}
  \item Bengio et al. 2003 \href{https://www.jmlr.org/papers/volume3/bengio03a/bengio03a.pdf}{A Neural Probabilistic Language Model}
  \item Mikolov et al. 2013 \href{https://arxiv.org/abs/1301.3781}{Word2Vec}
  \item Sutskever et al. 2014 \href{https://arxiv.org/abs/1409.3215}{Seq2Seq}；Cho et al. 2014 \href{https://arxiv.org/abs/1409.1259}{On the Properties of NMT}
  \item Bahdanau et al. 2014 \href{https://arxiv.org/abs/1409.0473}{Neural Machine Translation by Jointly Learning to Align and Translate}；Luong et al. 2015 \href{https://arxiv.org/abs/1508.04025}{Effective Approaches to Attention-based NMT}
  \item Vaswani et al. 2017 \href{https://arxiv.org/abs/1706.03762}{Attention Is All You Need}
  \item Sennrich et al. 2016 \href{https://arxiv.org/abs/1508.07909}{BPE for NMT}
  \item Devlin et al. 2018 \href{https://arxiv.org/abs/1810.04805}{BERT}；Raffel et al. 2019 \href{https://arxiv.org/abs/1910.10683}{T5}
  \item Radford et al. 2018 \href{https://cdn.openai.com/research-covers/language-unsupervised/language\_understanding\_paper.pdf}{GPT-1}；2019 \href{https://cdn.openai.com/better-language-models/language\_models\_are\_unsupervised\_multitask\_learners.pdf}{GPT-2}；Brown et al. 2020 \href{https://arxiv.org/abs/2005.14165}{GPT-3}
  \item Kaplan et al. 2020 \href{https://arxiv.org/abs/2001.08361}{Scaling Laws}；Hoffmann et al. 2022 \href{https://arxiv.org/abs/2203.15556}{Chinchilla}；Ouyang et al. 2022 \href{https://arxiv.org/abs/2203.02155}{InstructGPT}
  \item Xiong et al. 2020 \href{https://arxiv.org/abs/2002.04745}{On Layer Normalization in the Transformer Architecture}；Zhang \& Sennrich 2019 \href{https://arxiv.org/abs/1910.07467}{RMSNorm}
  \item Hendrycks \& Gimpel 2016 \href{https://arxiv.org/abs/1606.08415}{GELU}；Shazeer 2020 \href{https://arxiv.org/abs/2002.05202}{GLU Variants Improve Transformer}
  \item Su et al. 2021 \href{https://arxiv.org/abs/2104.09864}{RoFormer}；Press et al. 2021 \href{https://arxiv.org/abs/2108.12409}{ALiBi}；Chen et al. 2023 \href{https://arxiv.org/abs/2306.15595}{Position Interpolation}；Peng et al. 2023 \href{https://arxiv.org/abs/2309.00071}{YaRN}
  \item Shazeer 2019 \href{https://arxiv.org/abs/1911.02150}{MQA}；Ainslie et al. 2023 \href{https://arxiv.org/abs/2305.13245}{GQA}；DeepSeek-AI 2024 \href{https://arxiv.org/abs/2405.04434}{DeepSeek-V2 (MLA)}、\href{https://arxiv.org/abs/2412.19437}{DeepSeek-V3}
  \item Fedus et al. 2021 \href{https://arxiv.org/abs/2101.03961}{Switch Transformer}；Jiang et al. 2024 \href{https://arxiv.org/abs/2401.04088}{Mixtral}
  \item Touvron et al. 2023 \href{https://arxiv.org/abs/2302.13971}{Llama}、\href{https://arxiv.org/abs/2307.09288}{Llama 2}；Llama Team 2024 \href{https://arxiv.org/abs/2407.21783}{The Llama 3 Herd of Models}
\end{itemize}

\vspace{0.8em}\hrule\vspace{0.8em}

\textbf{上一篇}：\textbf{第 4 篇：Transformer 核心公式彻底推导：每一个符号从何而来} ｜ \textbf{下一篇}：\textbf{第 6 篇：什么是推理？自回归生成全生命周期与 KV Cache 代数证明} ｜ \textbf{总路线}：\textbf{LEARNING\_PATH.md}

\chapter{什么是推理？自回归生成全生命周期与 KV Cache 代数证明}

\begin{mathBox}[核心数学定理推导]
\textbf{阅读前须知}：许多人以为大模型“训练”和“推理”是一回事，只是输入不同而已。本篇我们将彻底厘清大模型推理（Inference）的物理本质，推导为什么模型只能逐字生成，以及为什么一个小小的 KV Cache 会成为大模型显存爆炸的万恶之源与工程关键。
\end{mathBox}

\vspace{0.8em}\hrule\vspace{0.8em}

\section{[对比] 一、训练（Training）vs 推理（Inference）的物理鸿沟}

在讨论推理之前，我们必须先对比两者在底层运行模式上的根本差异：

\begin{center}
\small
\begin{tabularx}{\textwidth}{p{3cm} p{4.5cm} X}
\toprule
\textbf{对比维度} & \textbf{训练阶段 (Training)} & \textbf{推理阶段 (Inference)} \\
\midrule
\textbf{已知信息} & \textbf{全部已知}（整篇文章的所有字都摆在眼前） & \textbf{完全未知}（只知道用户的 Prompt，答复一无所知） \\
\textbf{参数权重} & \textbf{频繁变动}（通过反向传播不断更新梯度） & \textbf{完全冻结}（只读参数，绝不更改） \\
\textbf{计算模式} & \textbf{全序列完全并行}（所有 Token 一瞬间同时算完） & \textbf{自回归逐字串行}（像挤牙膏一样一个词一个词蹦） \\
\textbf{显存杀手} & \textbf{激活值与优化器状态}（Adam 动量、权重梯度） & \textbf{模型权重与 KV Cache}（历史上下文缓存） \\
\bottomrule
\end{tabularx}
\end{center}

\subsection{训练时的“上帝视角”（Teacher Forcing）}

训练时，我们给模型一段完整的文字：“人工智能彻底改变了世界”。
因为整句话都在显卡里，模型可以在\textbf{一次前向计算中，同时预测所有位置的下一个词}：
\begin{itemize}
  \item 位置 1 的特征去预测“智能”；
  \item 位置 2 的特征去预测“彻底”；
  \item 位置 3 的特征去预测“改变”\dots{}\dots{}
\end{itemize}

成千上万个词的预测误差在这一瞬间同时算出，然后梯度反向传播一次性更新全部参数。这就是为什么训练高度并行。

\subsection{推理时的“摸黑前行”（自回归 Autoregressive Loop）}

推理时，用户发来一个问题：“北京的特产是什么？”
模型此时根本不知道后续几百个字是什么，它只能像盲人摸象一样，玩一个\textbf{接龙游戏}：

\begin{center}
\small
\begin{tabularx}{\textwidth}{p{1.8cm} X l}
\toprule
\textbf{解码步数} & \textbf{模型前向输入上下文（Context Tokens）} & \textbf{预测输出下一个 Token} \\
\midrule
第 1 步 & \texttt{["北京", "的", "特产", "是", "什么", "？"]} & \textbf{"烤"} \\
第 2 步 & \texttt{["北京", "的", "特产", "是", "什么", "？", "烤"]} & \textbf{"鸭"} \\
第 3 步 & \texttt{["北京", "的", "特产", "是", "什么", "？", "烤", "鸭"]} & \textbf{"。"}（触发 EOS 停止） \\
\bottomrule
\end{tabularx}
\end{center}

\textbf{上一步自己吐出来的词，拼入原句作为下一步的输入，周而复始，直到吐出结束标记 \texttt{$<$eos$>$} 为止。这就是大模型自回归推理的全部物理本质！}

\vspace{0.8em}\hrule\vspace{0.8em}

\section{[加速] 二、推理的两大完全不同阶段：Prefill 与 Decode}

在大模型推理系统（如 vLLM、TensorRT-LLM）中，你会每天听到两个术语：\textbf{Prefill（预填充）} 与 \textbf{Decode（解码）}。它们两者的硬件瓶颈截然相反！

\begin{center}
\small
\begin{tabularx}{\textwidth}{p{3.0cm} X X}
\toprule
\textbf{特征维度} & \textbf{阶段 1：Prefill 阶段 (首字计算)} & \textbf{阶段 2：Decode 阶段 (增量生成)} \\
\midrule
\textbf{输入形态} & 一次性输入全部 Prompt (如 1000 字) & 每次仅输入最新生成的 \textbf{1 个 Token} \\
\textbf{计算模式} & 大规模矩阵乘法 (GEMM) & 极小规模矩阵-向量乘法 (GEMV) \\
\textbf{硬件瓶颈} & \textbf{算力受限 (Compute-Bound)}，核心满载 & \textbf{访存受限 (Memory-Bound)}，带宽跑满 \\
\textbf{核心考核指标} & \textbf{TTFT} (首字时间，用户等待延迟) & \textbf{TPOT} (每字时间，流式打字流畅度) \\
\bottomrule
\end{tabularx}
\end{center}

\vspace{0.8em}\hrule\vspace{0.8em}

\section{[观察] 三、KV Cache 的代数展开与冗余消除证明}

现在我们进入推理系统中最核心的概念：\textbf{KV Cache}。

\subsection{朴素自回归计算中的重复劳动（无 KV Cache）}

假设当前已经生成到了第 $t$ 个词，序列为 $[x_1, x_2, \dots, x_{t-1}, x_t]$。
我们要预测第 $t+1$ 个词。

根据 Transformer 注意力机制：
\begin{enumerate}
  \item 计算整个序列的 Query、Key、Value：
\end{enumerate}

\[
Q_{1:t} = [x_1 W_q; x_2 W_q; \dots; x_t W_q]
\]

\[
K_{1:t} = [x_1 W_k; x_2 W_k; \dots; x_t W_k]
\]

\[
V_{1:t} = [x_1 W_v; x_2 W_v; \dots; x_t W_v]
\]

\begin{enumerate}
  \item 预测下一个词时，我们\textbf{只关心最后一个位置的 Query}：$q_t = x_t W_q$。
  \item 最后一个位置的注意力加权计算为：
\end{enumerate}

\[
\text{Output}_t = \text{softmax}\left(\frac{q_t K_{1:t}^T}{\sqrt{d_k}}\right) V_{1:t}
\]

\subsubsection{仔细看这里的 $K_{1:t}$ 和 $V_{1:t}$：}

\[
K_{1:t} = [ \underbrace{x_1 W_k, x_2 W_k, \dots, x_{t-1} W_k}_{\text{历史前 } t-1 \text{ 个词的 Key}}, \quad x_t W_k ]
\]

\[
V_{1:t} = [ \underbrace{x_1 W_v, x_2 W_v, \dots, x_{t-1} W_v}_{\text{历史前 } t-1 \text{ 个词的 Value}}, \quad x_t W_v ]
\]

\textbf{发现了什么惊人的事实？}
\begin{itemize}
  \item 权重矩阵 $W_k$ 和 $W_v$ 是固定冻结的；
  \item 历史上的词 $x_1, x_2, \dots, x_{t-1}$ 是已经说出来的话，绝不会再改变！
  \item 这意味着：\textbf{前面 $t-1$ 个词的 Key 和 Value 向量，在第 $t-1$ 步时就已经原模原样算过一次了！}
  \item 如果不保存它们，每生成一个新词，都要把前面所有历史词重新投影一遍。当文本长达 4000 字时，第 4001 步要把前 4000 个词的投影全量重算一遍！总计算量呈二次方 $O(T^2)$ 暴增，生成速度越来越慢！
\end{itemize}

\vspace{0.8em}\hrule\vspace{0.8em}

\subsection{KV Cache 的工程解法：以显存换算力}

既然历史的 Key 和 Value 永远不会变，那我们\textbf{在第一次算出它们时，就直接保存在显存的缓存池中（KV Cache）}！

\begin{mathBox}[第 $t$ 步推理：KV Cache 增量运算过程]
\begin{enumerate}
  \item \textbf{单向量投影}：仅对最新生成的单个 Token 向量 $x_t \in \mathbb{R}^{1 \times d_{\mathrm{model}}}$ 计算投影：
  \[
  q_t = x_t W_q, \quad k_t = x_t W_k, \quad v_t = x_t W_v
  \]
  \item \textbf{更新增量缓存}：将新产生的单步 $k_t, v_t$ 追加到已有显存缓存中：
  \[
  K_{\mathrm{all}} = \begin{bmatrix} K_{\mathrm{cache}} \\ k_t \end{bmatrix}, \quad V_{\mathrm{all}} = \begin{bmatrix} V_{\mathrm{cache}} \\ v_t \end{bmatrix}
  \]
  \item \textbf{注意力打分}：单向量 $q_t$ 与整个长序列 $K_{\mathrm{all}}$ 做矩阵乘法，得到 $[1, t]$ 注意力分布，再加权 $V_{\mathrm{all}}$ 得出输出向量！
\end{enumerate}
\end{mathBox}

\subsection{复杂度代数对比：}

\begin{itemize}
  \item \textbf{无 KV Cache}：每次前向需要做整个矩阵的投影 $O(t \cdot d_{\text{model}}^2)$，总共生成 $T$ 步需要累计 $O(T^2 \cdot d_{\text{model}}^2)$ 次计算。
  \item \textbf{有 KV Cache}：每次前向只需做单个向量的投影 $O(1 \cdot d_{\text{model}}^2)$，总计算量降为 $O(T \cdot d_{\text{model}}^2)$！
\end{itemize}

\textbf{投影与 FFN 这类"逐 token"矩阵乘法的总计算量少了一个因子 $T$（而不是常数倍）。}

\begin{warningBox}[避坑指南与核心警示]
{}[注意] \textbf{严谨补充（容易被忽略）}：KV Cache 消除的是"对历史 token 重复做投影 / FFN"的冗余，\textbf{并没有消除注意力本身的二次方}：
第 $t$ 步里，新的 $q_t$ 仍然要和缓存中全部 $t$ 个 $k$ 做点积、再对 $t$ 个 $v$ 加权，代价是 $O(t \cdot d)$。累加 $T$ 步后注意力部分仍是 $O(T^2 d)$。
所以"有 Cache 后每步耗时恒定"只在 $t \cdot d \ll d_{\text{model}}^2$（上下文不长）时近似成立；上下文变长后，每一步都要从显存\textbf{读一遍整个 KV Cache}，这正是长上下文 Decode 变慢、以及 FlashDecoding / GQA / MLA / KV 量化出现的原因（见第 8、9、12 篇）。
\end{warningBox}

\vspace{0.8em}\hrule\vspace{0.8em}

\section{[隐患] 四、KV Cache 的代价：显存爆炸与 PagedAttention 的诞生}

天下没有免费的午餐。KV Cache 用空间换取了时间，但带来了一个极度致命的副作用：\textbf{显存空间被疯狂吞噬！}

\subsection{KV Cache 显存占用精确公式推导}

对于一个大模型：
\begin{itemize}
  \item 层数 $L$（每一层 Transformer 都有自己独立的 $K$ 和 $V$）
  \item 隐藏层维度 $H$
  \item 上下文总长度 $S$（输入的 Prompt 词数 + 已生成的词数）
  \item 并发请求数 $B$（Batch Size，同时服务的用户数）
  \item 浮点精度字节数 $b$（FP16/BF16 每个数字占 2 字节）
\end{itemize}

每个词需要存储一个 Key 向量（$H$ 个数字）和一个 Value 向量（$H$ 个数字），跨越全部 $L$ 层：

\[
\text{每个 Token 的 KV 显存} = 2 \times L \times H \times b \quad \text{(字节)}
\]

如果同时服务 $B$ 个用户，平均每个请求长 $S$ 个词：

\[
\text{总 KV Cache 显存} = B \times S \times 2 \times L \times H \times b
\]

\subsubsection{实例计算（一个 80 层、隐藏维度 8192、采用原始 MHA 的 70B 级模型）：}

\begin{noteBox}[核心导读与背景]
注：真实的 Llama-2-70B 已经使用了 GQA（64 个 Q 头、8 个 KV 头），下面的数字是"假如它仍用原始 MHA"的反事实计算——正是这个数字吓人，才逼出了 GQA。LLaMA-1-65B（80 层、8192 维、MHA）的量级与之相同。
\end{noteBox}

\begin{itemize}
  \item $L = 80$, $H = 8192$, $b = 2$（FP16）
  \item 单个 Token 的 KV Cache 显存：
\end{itemize}

\[
2 \times 80 \times 8192 \times 2 = 2,621,440 \text{ 字节} \approx 2.5 \text{ MB / token}
\]

\begin{itemize}
  \item 如果有 16 个用户并发（$B=16$），每个人对话长度为 4096 字（$S=4096$）：
\end{itemize}

\[
16 \times 4096 \times 2.5 \text{ MB} \approx 163,840 \text{ MB} = \mathbf{160 \text{ GB 显存}}!
\]

看到了吗？光是保存这一批对话的 KV Cache，就需要整整 \textbf{160 GB} 显存，这比 70B 模型本身的权重（140 GB）还要大！
而换成 GQA-8（Llama-2/3-70B 的真实配置）后，每 token 只需 $2 \times 80 \times (8 \times 128) \times 2 = 327{,}680$ 字节 $\approx 320$ KB，同样场景只要 \textbf{20 GB}——这就是 GQA 的价值。

\subsection{现代模型的架构救赎：GQA 与 PagedAttention}

为了解决这个显存无底洞，学术界与工业界爆发了两场革命：
\begin{enumerate}
  \item \textbf{模型结构革命：GQA（分组查询注意力）}
\end{enumerate}

\begin{itemize}
  \item 既然存储全部 64 个头太费显存，那就让 8 个 Query 头共享 1 组 Key/Value 头（KV 头数压缩为 1/8）。
  \item 现代 Llama-3、Qwen2.5 广泛采用 GQA，将 KV Cache 显存直接缩减为原来的 \textbf{1/4 到 1/8}！
\end{itemize}

\begin{enumerate}
  \item \textbf{显存管理革命：PagedAttention（vLLM 破圈之作）}
\end{enumerate}

\begin{itemize}
  \item 传统框架必须为每个请求提前申请连续的显存块（例如预留 4096 个词的空间），但大部分请求只聊了 100 句就退出了。vLLM 论文测到 60\%\textasciitilde{}80\% 的 KV 显存被浪费——其中既有按 max\_len 预留的闲置，也有内部碎片与外部碎片（见本节的下一段），三者要分开看。
  \item vLLM 借鉴了操作系统的虚拟分页技术，将 KV Cache 切成 16 个 token 一块的物理页，按需分配，零碎片浪费，让服务器并发能力暴增数倍！
\end{itemize}

\vspace{0.8em}\hrule\vspace{0.8em}

\section{[目标] 总结与下一章预告}

在本篇中，我们彻底理清了大模型推理的微观全貌：
\begin{enumerate}
  \item \textbf{推理与训练的区别}：推理是权重冻结的自回归单向接龙；
  \item \textbf{Prefill vs Decode}：Prefill 是首字并行算力受限（Compute-Bound），Decode 是逐字生成访存受限（Memory-Bound）；
  \item \textbf{KV Cache 代数展开}：证明了历史 Key 和 Value 的不变性，通过缓存把\textbf{逐 token 重复做投影/FFN} 的部分从 $O(T^2)$ 降到 $O(T)$（注意力本身仍是 $O(T^2d)$，长上下文下它反过来成为主导，见本节 §3 的补充）；
  \item \textbf{显存挑战}：KV Cache 随并发与长度线性膨胀，催生了 GQA 与 PagedAttention。
\end{enumerate}

在下一篇中，我们将探讨：
\textbf{模型输出 Logits 之后，最终的文字是怎么被选择出来的？温度系数（Temperature）在数学极限下到底证明了什么？Top-K 和 Top-P 是怎么把胡言乱语切除的？}

\begin{mathBox}[核心数学定理推导]
{}[重点] 本篇说"Decode 访存受限"，但没有给出定量推导。为什么 batch 开大几乎不增加每一步的耗时？GQA 除了省显存还能提速吗？这些在 \textbf{第 8 篇：硬件视角} 中推导，并在你的显卡上实测。
显存管理（PagedAttention）的完整机制与实现见 \textbf{第 12 篇} 与 Lab 12b。
\end{mathBox}

\vspace{0.8em}\hrule\vspace{0.8em}

\textbf{上一篇}：\textbf{第 5 篇：从原版 Transformer 到 Llama / DeepSeek：每一处改动的历史动机} ｜ \textbf{下一篇}：\textbf{第 7 篇：大模型的决策大脑：Logits、温度系数极限证明与采样算法} ｜ \textbf{总路线}：\textbf{LEARNING\_PATH.md}

\chapter{大模型的决策大脑：Logits、温度系数极限证明与采样算法}

\begin{mathBox}[核心数学定理推导]
\textbf{阅读前须知}：经过了前面四篇的学习，大模型的 Transformer 网络已经完成了前向传播，把前文所有的上下文浓缩进了一个最终的隐藏层向量中。在这最后一公里，这个向量是如何变成真实世界的文字的？为什么温度调高了模型会胡思乱想？调低了就会一本正经？本篇将为你带来严格的数学极限证明与算法解析。
\end{mathBox}

\vspace{0.8em}\hrule\vspace{0.8em}

\section{[目标] 一、最后的终点线：LM Head（语言模型头）}

在模型的最后一层，针对序列的最后一个位置，输出一个经过归一化的隐藏特征向量：

\[
h_{\text{last}} \in \mathbb{R}^{d_{\text{model}}} \quad \text{(例如 } 4096 \text{ 维)}
\]

这个向量就像是一团高度压缩的“思考精华”，里面包含了模型对“下一个字应该是什么”的所有语义判断。
但它还不是人类能看懂的字，我们需要把它映射回\textbf{整个词表（Vocabulary）}。

\subsection{词表矩阵乘法}

假设模型的词表里包含 $V$ 个词（现代大模型如 Llama-3 的词表大小 $V = 128,256$；Qwen2.5 小尺寸 151,936、7B 及以上 152,064）。
我们使用一个投影矩阵 $W_{\text{head}} \in \mathbb{R}^{d_{\text{model}} \times V}$ 与之相乘：

\[
z = h_{\text{last}} W_{\text{head}} \in \mathbb{R}^V
\]

这个输出向量 $z$ 的长度恰好等于词表大小 $V$。
向量里的每一个数值，被称为一个 \textbf{Logit（未归一化对数打分）}：

\[
z = [z_1, z_2, \dots, z_V]
\]

\begin{itemize}
  \item $z_{1024}$ 对应词表里的“的”，数值可能是 $14.5$（极高）；
  \item $z_{305}$ 对应词表里的“烤鸭”，数值可能是 $8.2$；
  \item $z_{9999}$ 对应词表里的“宇宙飞船”，数值可能是 $-3.1$（极低）。
\end{itemize}

现在，模型拿到了对全人类十几万个词汇的打分单。接下来的关键问题是：\textbf{怎么做最终的选词决策？}

\vspace{0.8em}\hrule\vspace{0.8em}

\section{[剖析] 二、决策方法 1：贪婪搜索（Greedy Search）}

最直观的方法，就是直接挑选打分最高的那一个词：

\[
\hat{y} = \arg\max_{i} (z_i)
\]

\begin{center}
\small
\begin{tabularx}{0.7\textwidth}{l c X}
\toprule
\textbf{候选 Token} & \textbf{模型输出 Logit 打分} & \textbf{选择策略判定} \\
\midrule
\texttt{"的"} & $14.5$ & \textbf{最高分！Greedy 策略直接选中} \\
\texttt{"是"} & $12.1$ & 次高分 \\
\texttt{"在"} & $9.3$ & 低分 \\
\texttt{"飞船"} & $-2.0$ & 极低概率 \\
\bottomrule
\end{tabularx}
\end{center}

\subsection{优缺点分析：}

\begin{itemize}
  \item \textbf{优点}：计算最快，结果 100\% 确定，非常适合\textbf{数学推理、编写代码、事实性问答}等有唯一确定答案的场景。
  \item \textbf{缺点}：极其死板，缺乏创造力。更致命的是，贪婪算法只看眼前单步最优，容易陷入局部最优陷阱，导致模型在长文本生成时陷入可怕的\textbf{无脑复读死循环}（“我认为\dots{}\dots{}我认为\dots{}\dots{}我认为\dots{}\dots{}”）。
\end{itemize}

为了打破这种死板，人类引入了\textbf{随机概率采样（Stochastic Sampling）}。

\vspace{0.8em}\hrule\vspace{0.8em}

\section{[温度] 三、【重磅数学推导】：温度系数（Temperature）的极限证明}

为了在采样中可控地调节模型的“严肃度”与“创造力”，科学家在 Softmax 中引入了一个正实数参数 $T$（Temperature，温度系数）：

\[
p_i(T) = \frac{\exp\left(\frac{z_i}{T}\right)}{\sum_{j=1}^V \exp\left(\frac{z_j}{T}\right)}
\]

你一定在各类大模型界面中调节过这个滑动条（通常在 $0.0 \sim 2.0$ 之间）。
\textbf{但你是否知道，在严谨的数学分析中，温度参数在极限状态下到底意味着什么？}

下面我们给出两个经典的\textbf{数学定理与极限推导}。

\vspace{0.8em}\hrule\vspace{0.8em}

\subsection{定理 1：若最大 logit 唯一，当温度 $T \to 0^+$ 时，概率分布严格退化为贪婪搜索（确定性决策）}

\begin{mathBox}[核心数学定理推导]
\textbf{【证明】}：
设 Logits 向量中存在唯一的最大值 $z_{\max} = \max_j (z_j)$。
对任意其他非最大项 $i \neq \max$，均有：
\end{mathBox}

\[
z_i - z_{\max} < 0
\]

\begin{noteBox}[核心导读与背景]
我们将 Softmax 公式分子与分母同时除以 $\exp\left(\frac{z_{\max}}{T}\right)$（利用平移不变性）：
\end{noteBox}

\[
p_i(T) = \frac{\exp\left(\frac{z_i - z_{\max}}{T}\right)}{\sum_{j=1}^V \exp\left(\frac{z_j - z_{\max}}{T}\right)}
\]

\begin{noteBox}[核心导读与背景]
考察当 $T \to 0^+$ 时的极限行为：

\textbf{情况 A：对于最大值项自身（$i = \max$）}：
\end{noteBox}

\[
z_{\max} - z_{\max} = 0 \implies \exp\left(\frac{0}{T}\right) = e^0 = 1
\]

\begin{noteBox}[核心导读与背景]
\textbf{情况 B：对于任意非最大值项（$i \neq \max$）}：
因为 $z_i - z_{\max} < 0$（为一个严格负常数 $-c$），当 $T \to 0^+$ 时：
\end{noteBox}

\[
\lim_{T \to 0^+} \frac{z_i - z_{\max}}{T} = -\infty
\]

\begin{noteBox}[核心导读与背景]
根据自然指数极限：
\end{noteBox}

\[
\lim_{T \to 0^+} \exp\left(\frac{z_i - z_{\max}}{T}\right) = \lim_{x \to -\infty} e^x = \mathbf{0}
\]

\begin{noteBox}[核心导读与背景]
现在把分子和分母的极限合并：
- 分母的总和：
\end{noteBox}

\[
\lim_{T \to 0^+} \sum_{j=1}^V \exp\left(\frac{z_j - z_{\max}}{T}\right) = 1 + \underbrace{0 + 0 + \dots + 0}_{V-1 \text{ 个}} = \mathbf{1}
\]

\begin{noteBox}[核心导读与背景]
- 最大值项的概率：
\end{noteBox}

\[
\lim_{T \to 0^+} p_{\max}(T) = \frac{1}{1} = \mathbf{1.0} \quad (100\%)
\]

\begin{noteBox}[核心导读与背景]
- 所有非最大值项的概率：
\end{noteBox}

\[
\lim_{T \to 0^+} p_{i \neq \max}(T) = \frac{0}{1} = \mathbf{0.0} \quad (0\%)
\]

\begin{noteBox}[核心导读与背景]
\textbf{Q.E.D. 证毕！}
\end{noteBox}

\subsubsection{物理直觉：}

\begin{warningBox}[避坑指南与核心警示]
{}[注意] \textbf{"最大值唯一"这个前提不能省}：如果有 $k$ 个词的 logit 并列最大，同样的推导给出极限是\textbf{这 $k$ 个词各占 $1/k$ 的均匀分布}，而不是确定性地选其中一个。真实模型里 logit 恰好相等很少见，但 BF16 精度下前几名"数值上相等"并不罕见，这也是低温采样有时仍有随机性的原因之一。
\end{warningBox}

\textbf{温度越低（$T \to 0$），最大值项的优势被无限放大，概率瞬间形成尖锐的独峰（全部概率集中在一个词上的"单点分布"，即 one-hot），模型 100\% 必然挑选中打分最高的词！这就是为什么调低温度能让回答严谨、稳定、绝不胡说八道！}

\vspace{0.8em}\hrule\vspace{0.8em}

\subsection{定理 2：当温度 $T \to +\infty$ 时，概率分布严格退化为均匀分布（纯随机乱码）}

\begin{mathBox}[核心数学定理推导]
\textbf{【证明】}：
考察指数内部项在 $T \to +\infty$ 时的极限：
对于任意有限实数 $z_i$：
\end{mathBox}

\[
\lim_{T \to +\infty} \frac{z_i}{T} = 0
\]

\begin{noteBox}[核心导读与背景]
因此：
\end{noteBox}

\[
\lim_{T \to +\infty} \exp\left(\frac{z_i}{T}\right) = e^0 = 1
\]

\begin{mathBox}[核心数学定理推导]
代入 Softmax 公式：
\end{mathBox}

\[
\lim_{T \to +\infty} p_i(T) = \frac{1}{\sum_{j=1}^V 1} = \frac{1}{V}
\]

\begin{noteBox}[核心导读与背景]
\textbf{Q.E.D. 证毕！}
\end{noteBox}

\subsubsection{物理直觉：}

\textbf{温度无穷大（$T \to \infty$）时，词表中每一个词（不论是常见字还是生僻乱码）被选中的概率全部变成了毫无差别的 $1/V$！模型丧失了一切前向推理学习到的知识，完全退化为一个均匀扔骰子的随机数发生器！}

\vspace{0.8em}\hrule\vspace{0.8em}

\section{[目标] 四、工业级截断双子星：Top-K 与 Top-P}

单纯靠 Temperature 调节依然有风险：即使把温度设为正常的 $0.7$，长尾里那些概率只有 $0.0001\%$ 的荒谬生僻字，依然有极其微小的概率被“不幸抽中”，导致模型偶尔冒出一句胡话。

为了彻底扫清这些低质噪声，工业界标配了两种截断算法：

\subsection{Top-K 采样}

\begin{itemize}
  \item \textbf{逻辑}：不管词表有多大，\textbf{永远只保留概率最高的前 $K$ 个候选词}（例如 $K=40$ 或 $K=50$）。
  \item 将第 $K+1$ 到第 $V$ 个词的 Logit 强行赋值为 $-\infty$（概率归零）。
  \item 在保留的这 $K$ 个词中重新归一化 Softmax。
\end{itemize}

\begin{center}
\small
\begin{tabularx}{\textwidth}{X c X}
\toprule
\textbf{高概率候选词（保留采样池）} & \textbf{截断阈值} & \textbf{长尾低概率词（强制置 $-\infty$）} \\
\midrule
前 $K$ 个最高概率词: \texttt{[词 1, 词 2, $\dots$, 词 $K$]} & $\longleftrightarrow$ & 剩余词: \texttt{[词 $K+1$, $\dots$, 词 $V$]} 全部过滤 \\
\bottomrule
\end{tabularx}
\end{center}

\vspace{0.8em}\hrule\vspace{0.8em}

\subsection{Top-P 采样（核采样 Nucleus Sampling，更智能的动态截断）}

Top-K 有一个天然的死穴：\textbf{固定 $K$ 值无法适应不同语境的灵活性！}
\begin{itemize}
  \item \textbf{场景 A}：用户输入“人工智能的英文缩写是\_\_\_\_”。
  \item 这种问题几乎只有唯一正确答案“AI”。此时理想的候选池大小应该是 $K=1$。如果固定 $K=50$，模型反而可能把后 49 个错误答案放进抽奖箱。
  \item \textbf{场景 B}：用户输入“请描述一个美丽的黄昏景色：\_\_\_\_”。
  \item 此时有成百上千个优美的辞藻（“霞光”、“落日”、“晚风”、“云彩”）都是合理的，此时 $K=50$ 又显得太少，扼杀了多样性。
\end{itemize}

\textbf{Top-P 的破局之道（Holtzman et al., ICLR 2020）}：
\begin{noteBox}[核心导读与背景]
\textbf{不要固定候选词的个数，而是固定候选词的「累积置信度阈值 $P$」（例如 $P = 0.9$ 或 $90\%$）！}
\end{noteBox}

\begin{center}
\small
\begin{tabularx}{\textwidth}{p{3.5cm} c X}
\toprule
\textbf{语境场景} & \textbf{候选 Token 概率累积} & \textbf{动态候选池判定} \\
\midrule
\textbf{极确定语境} (如代码/事实) & \texttt{"AI"}: 92\% $\ge 90\%$ & \textbf{立即截断！候选池仅保留 1 个词} \\
\midrule
\textbf{发散创意语境} (如诗歌/创作) & \texttt{"落日"}: 25\%, \texttt{"晚霞"}: 20\%, $\dots$ & \textbf{累积至 90.5\% 截断！动态保留 15 个候选词} \\
\bottomrule
\end{tabularx}
\end{center}

\begin{itemize}
  \item \textbf{语境确定时}：头部词概率极高，累积迅速破 90\%，候选池自动收缩为 1\textasciitilde{}2 个词，极其精准；
  \item \textbf{语境发散时}：头部词概率分散，需要累积几十个词才破 90\%，候选池自动扩大，充分保障多样性！
\end{itemize}

\vspace{0.8em}\hrule\vspace{0.8em}

\section{五、防复读机制：重复惩罚（Repetition Penalty）}

你一定见过模型陷入“复读机”状态。怎么在数学上给它一记当头棒喝？

在生成第 $t$ 步时，算法会检查前文中\textbf{已经出现过的所有词}的索引集合 $\mathcal{S}_{\text{history}}$。
对于集合中的词，对其打分 $z_i$ 施加一个惩罚系数 $\theta > 1.0$（通常取 $1.1 \sim 1.3$）：

\[
z_i' = \begin{cases}
z_i / \theta & \text{若 } z_i > 0 \quad (\text{主动压低正向打分}) \\
z_i \times \theta & \text{若 } z_i \le 0 \quad (\text{加大负向惩罚力度})
\end{cases}
\]

\textbf{通过这一轻巧的代数缩放，曾经说过的词在下一步重新争夺第一名的难度被大幅提高，模型被迫转向寻找新的表达方式，显著缓解复读机问题。}

\begin{warningBox}[避坑指南与核心警示]
{}[注意] 但它是一把钝器：它惩罚的是"出现过的所有 token"，不区分该不该重复。人名、专有名词、标点、代码里的变量名本来就需要反复出现，惩罚过重会让模型改用同义词、写错变量名，甚至破坏格式。而且它连提示词里的 token 也一起惩罚（Lab 07b 实验 5 里 HF 的实现就是这样）。所以实际服务中惩罚系数通常只取 1.05\textasciitilde{}1.2，另有 frequency / presence penalty 等更细的变体。
\end{warningBox}

\vspace{0.8em}\hrule\vspace{0.8em}

\section{[检验] 前七篇小结}

到这里，你已经\textbf{从零起点、不依赖任何黑盒}，贯通了"推理在算什么"的整条链路：

\begin{enumerate}
  \item 第 1 篇：计算机如何从规则、统计到被 Transformer 彻底征服
  \item 第 2 篇：向量空间几何、点积的几何意义与 Safe Softmax 溢出防护
  \item 第 3 篇：导数、反向传播、交叉熵，Softmax 雅可比与 $p - y$ 梯度
  \item 第 4 篇：严格证明 Attention 为什么必须除以 $\sqrt{d_k}$，推导 RoPE 相对位置内积恒等性
  \item 第 5 篇：从原版 Transformer 到 Llama / DeepSeek，每一处改动的历史动机
  \item 第 6 篇：推理与训练的本质区别，代数展开 KV Cache 消除冗余并剖析显存挑战
  \item 第 7 篇：温度参数在 $T \to 0$ 和 $T \to \infty$ 下的数学极限，Top-K、Top-P 与重复惩罚
\end{enumerate}

但要真正理解推理\textbf{为什么快、为什么慢、怎么优化}，还需要继续：
{}\textbf{第 8 篇：硬件视角} $\rightarrow$ \textbf{第 9 篇：FlashAttention} $\rightarrow$ \textbf{第 10 篇：投机解码} $\rightarrow$ \textbf{第 11 篇：量化} $\rightarrow$ \textbf{第 12 篇：推理服务系统} $\rightarrow$ \textbf{第 13 篇：多卡并行}

完整顺序（含实验与外部资源）见 \textbf{LEARNING\_PATH.md}。

\vspace{0.8em}\hrule\vspace{0.8em}

\textbf{上一篇}：\textbf{第 6 篇：什么是推理？自回归生成全生命周期与 KV Cache 代数证明} ｜ \textbf{下一篇}：\textbf{第 8 篇：硬件视角：GPU 存储层级、FLOPs、算术强度与浮点格式} ｜ \textbf{总路线}：\textbf{LEARNING\_PATH.md}

\chapter{硬件视角：GPU 存储层级、FLOPs、算术强度与浮点格式}

\begin{mathBox}[核心数学定理推导]
\textbf{阅读前须知}：第 6 篇和 Lab 08 告诉了你一个结论——"Prefill 算力受限，Decode 访存受限"。但为什么？受限到什么程度？batch 开多大才能把 GPU 喂饱？GQA 除了省显存还有没有别的好处？量化为什么能加速？这些问题都需要\textbf{把硬件放进公式里}才能回答。本篇是从"懂模型"走向"懂推理系统"的分水岭。

\textbf{前置}：第 6 篇（Prefill/Decode、KV Cache）、第 3 篇（$2N$ FLOPs/token）、第 5 篇（GQA/MLA）。

\textbf{配套实验}：\textbf{\texttt{lab08\_memory\_and\_roofline/measure\_gpu\_roofline.py}}——在你自己的 RTX 5060 Ti 上实测带宽、峰值算力、Decode 形状矩阵乘随 batch 的变化曲线，以及 BF16/FP16 的精度陷阱。
\end{mathBox}

\vspace{0.8em}\hrule\vspace{0.8em}

\section{[硬件] 一、GPU 的基本构造：计算很便宜，搬数据很贵}

\subsection{计算单元}

\begin{itemize}
  \item GPU 由几十到上百个 \textbf{SM（Streaming Multiprocessor，流式多处理器）} 组成。RTX 5060 Ti 有 36 个 SM，H100 有 132 个。
  \item 每个 SM 里有两类计算单元：
  \item \textbf{CUDA Core}：做普通的标量浮点运算（FP32 加乘）；
  \item \textbf{Tensor Core}：专门做\textbf{小矩阵乘累加}（例如一次完成 $16 \times 16$ 的块乘），低精度（BF16/FP16/FP8/FP4）吞吐是 CUDA Core 的十几倍以上。\textbf{大模型的矩阵乘几乎全部跑在 Tensor Core 上。}
  \item 线程以 \textbf{32 个为一组（warp）} 同步执行同一条指令。
\end{itemize}

\subsection{存储层级（越往下越大、越慢）}

\begin{center}
\small
\begin{tabularx}{\textwidth}{p{2.5cm} p{3.5cm} X X}
\toprule
\textbf{层级} & \textbf{容量量级} & \textbf{带宽量级} & \textbf{谁能访问} \\
\midrule
寄存器 & 每 SM 256 KB & 最快 & 单个线程 \\
共享内存 / L1（SRAM） & 每 SM 约 100\textasciitilde{}228 KB & 全卡合计数十 TB/s & 同一个 SM 内的线程块 \\
L2 缓存 & 全卡数十 MB & 数 TB/s & 全卡 \\
\textbf{显存（HBM / GDDR）} & 16 \textasciitilde{} 192 GB & \textbf{0.4 \textasciitilde{} 8 TB/s} & 全卡 \\
\bottomrule
\end{tabularx}
\end{center}

RTX 5060 Ti 的显存是 16GB GDDR7、128 位总线，标称带宽 \textbf{448 GB/s}。作为对比，H100 SXM 的 HBM3 带宽是 3.35 TB/s。

\begin{tipBox}[核心直觉与思考]
{}[思考] \textbf{一个工厂比喻}（来自 Horace He 的 \textit{Making Deep Learning Go Brrrr}）：Tensor Core 是一座效率极高的工厂，显存是远处的仓库，显存带宽是连接两者的公路。工厂再快，原料运不过来也只能停工。\textbf{大模型 Decode 的处境就是：工厂 90\% 以上的时间在等货车。}
\end{tipBox}

\vspace{0.8em}\hrule\vspace{0.8em}

\section{[数字] 二、数清楚两件事：算了多少（FLOPs）、搬了多少（Bytes）}

\subsection{矩阵乘的 FLOPs}

$[M, K] \times [K, N]$：输出有 $MN$ 个元素，每个是 $K$ 次乘法 + $K$ 次加法：

\[
\text{FLOPs} = 2MKN
\]

一个参数量为 $N_{\text{param}}$ 的模型，处理一个 token 约 $2N_{\text{param}}$ FLOPs（第 3 篇）。

\subsection{矩阵乘的访存量}

每个元素 $b$ 字节（BF16 为 2），至少要读两个输入、写一个输出：

\[
\text{Bytes} = b\,(MK + KN + MN)
\]

\subsection{算术强度（Arithmetic Intensity）}

\[
\text{AI} = \frac{\text{FLOPs}}{\text{Bytes}} \quad (\text{单位：FLOP/Byte，每搬 1 字节做多少次运算})
\]

\vspace{0.8em}\hrule\vspace{0.8em}

\section{[几何] 三、Roofline 模型：一张图判断你被谁卡住}

硬件有两个上限：峰值算力 $\pi$（FLOP/s）和显存带宽 $\beta$（Byte/s）。一个算术强度为 AI 的算子，能达到的性能是：

\[
\text{可达性能} = \min(\pi, \ \text{AI} \times \beta)
\]

\begin{figure}[htbp]
\centering
\begin{tikzpicture}[scale=1.0, >=stealth]
  \draw[->, thick, darktext] (0,0) -- (8.5,0) node[below right, font=\small\sffamily] {算术强度 $I$ (FLOP/Byte)};
  \draw[->, thick, darktext] (0,0) -- (0,4.8) node[above left, font=\small\sffamily] {计算性能 $P$ (FLOP/s)};
  
  \coordinate (O) at (0,0);
  \coordinate (Ridge) at (3.5,3.6);
  \coordinate (MaxPeak) at (8.0,3.6);
  
  \draw[very thick, primary] (O) -- (Ridge) node[midway, above left, sloped, font=\small\sffamily\bfseries, color=primary] {访存受限区 (斜率 = 显存带宽 $\beta$)};
  \draw[very thick, warnred] (Ridge) -- (MaxPeak) node[midway, above, font=\small\sffamily\bfseries, color=warnred] {算力受限区 (峰值算力 $\pi$)};
  
  \draw[dashed, codegray] (3.5,0) node[below, font=\small\sffamily] {$I^* = \frac{\pi}{\beta}$ (拐点 Ridge Point)} -- (Ridge);
  \draw[dashed, codegray] (0,3.6) node[left, font=\small\sffamily] {$\pi$} -- (Ridge);
  
  \fill[warnred] (Ridge) circle (2.5pt);
\end{tikzpicture}
\caption{GPU Roofline 模型与瓶颈判定边界}
\end{figure}

\textbf{拐点} $\text{AI}^* = \pi / \beta$：算子的 AI 低于它，就是访存受限；高于它，就是算力受限。

\begin{center}
\small
\begin{tabularx}{\textwidth}{p{2.5cm} p{3.5cm} X X}
\toprule
\textbf{硬件} & \textbf{BF16 稠密峰值 $\pi$} & \textbf{带宽 $\beta$} & \textbf{拐点 $\text{AI}^*$} \\
\midrule
A100 SXM & 312 TFLOP/s & 2.04 TB/s & $\approx$ 153 \\
H100 SXM & 989 TFLOP/s & 3.35 TB/s & $\approx$ 295 \\
RTX 5060 Ti & \textbf{用配套实验实测} & 448 GB/s（标称） & 实测后自己算 \\
\bottomrule
\end{tabularx}
\end{center}

\begin{warningBox}[避坑指南与核心警示]
{}[注意] 厂商宣传的 Tensor 算力常常是"稀疏"或"FP16 累加"口径，GeForce 卡在 FP32 累加（训练/推理实际用的模式）下常常只有宣传值的一半。\textbf{永远以实测为准}——这也是配套实验存在的意义。
\end{warningBox}

\vspace{0.8em}\hrule\vspace{0.8em}

\section{[目标] 四、核心推导：Decode 的算术强度 \texorpdfstring{$\approx$}{≈} Batch Size}

\subsection{线性层（占 Decode 权重读取的绝大部分）}

Decode 时 batch 里有 $B$ 个请求，每个请求只有 1 个新 token，所以线性层是 $[B, K] \times [K, N]$，而权重 $KN$ 远大于激活 $BK$、$BN$（$B \ll K, N$）：

\[
\text{AI} = \frac{2BKN}{b(BK + KN + BN)} \approx \frac{2BKN}{b \cdot KN} = \frac{2B}{b} \overset{\text{BF16}}{=} \boxed{\ B\ }
\]

\textbf{Decode 时线性层的算术强度约等于 batch size。} 这一个结论解释了一大串现象：
\begin{itemize}
  \item batch = 1 时 AI $\approx$ 1，而拐点在一两百 $\rightarrow$ GPU 算力利用率只有 \textbf{1\% 量级}；
  \item 把 batch 从 1 提到 64，每一步读的权重\textbf{几乎不变}，时间几乎不变，吞吐却提升接近 64 倍——这就是\textbf{批处理的全部意义}，也是 Lab 12a 连续批处理仿真里"每一步耗时视为常数"这个假设的来源；
  \item 只有当 $B$ 接近拐点（A100 上约 150）后，再加 batch 才会让每一步变慢。
\end{itemize}

\subsection{Prefill}

Prefill 时 $M$ = prompt 长度。例如 $M = 2000$、$K = N = 4096$，代入完整公式：

\[
\text{AI} = \frac{2MKN}{b(MK + KN + MN)} = \frac{2 \cdot 2000 \cdot 4096^2}{2(2000 \cdot 4096 + 4096^2 + 2000 \cdot 4096)} \approx 1012
\]

（注意这里 $M$ 已经不再远小于 $K, N$，上面 Decode 用的近似 $2M/b$ 只是 $M \to 0$ 时的上界，会高估一倍。）即便如此，1012 也远超拐点 $\rightarrow$ \textbf{算力受限}。
所以 Prefill 的耗时估算是：

\[
T_{\text{prefill}} \approx \frac{2 N_{\text{param}} \cdot S_{\text{prompt}}}{\pi \cdot \text{MFU}}
\]

MFU（Model FLOPs Utilization）是实际达到峰值的比例，优秀实现通常 40\%\textasciitilde{}70\%。

\subsection{Decode 中的注意力：batch 救不了它}

线性层的权重被 batch 里所有请求\textbf{共享}，所以 batch 能摊薄权重读取。但\textbf{每个请求的 KV Cache 是它自己的}，只有它自己的 query 会用。
单个请求、单层、上下文长度 $S$：Q 头数 $n_h$、KV 头数 $n_{kv}$、头维度 $d_h$：

\[
\text{FLOPs} = \underbrace{2 n_h S d_h}_{QK^T} + \underbrace{2 n_h S d_h}_{PV} = 4 n_h S d_h, \qquad \text{Bytes} = \underbrace{2 S n_{kv} d_h}_{K \text{ 和 } V} \cdot b
\]

\[
\text{AI}_{\text{attn}} = \frac{4 n_h S d_h}{2 S n_{kv} d_h b} = \frac{2}{b} \cdot \frac{n_h}{n_{kv}} \overset{\text{BF16}}{=} \boxed{\ \frac{n_h}{n_{kv}}\ }
\]

\textbf{Decode 注意力的算术强度 = GQA 的分组大小，与 batch 无关、与上下文长度无关。}
\begin{itemize}
  \item MHA：AI = 1，永远被带宽卡死；
  \item Llama-3-8B（32 Q 头 / 8 KV 头）：AI = 4；
  \item MLA 让大量 Q 头共享同一个压缩潜向量，AI 更高。
\end{itemize}

\begin{warningBox}[避坑指南与核心警示]
{}[目标] \textbf{所以 GQA/MLA 的价值有两层}：第 5 篇讲的是\textbf{省显存}（能放更多请求、更长上下文），这里推出的是\textbf{提速}（同样的 KV 字节被更多次计算复用）。长上下文时注意力的 KV 读取会超过权重读取，成为 Decode 的主要耗时——这时 GQA 分组大小直接决定速度。
\end{warningBox}

\subsection{把两部分合起来：一步 Decode 的耗时模型}

\[
T_{\text{step}}(B) \approx \max\left( \frac{W_{\text{bytes}} + B \cdot S \cdot \text{KV}_{\text{bytes/token}}}{\beta}, \ \frac{2 N_{\text{param}} B}{\pi} \right)
\]

系统吞吐 $= B / T_{\text{step}}(B)$，单用户速度 $= 1 / T_{\text{step}}(B)$。
\textbf{加大 batch：总吞吐上升，单用户速度下降}——这就是推理服务永恒的"吞吐-延迟权衡"，Lab 08 的 \texttt{memory\_calculator.py} 已经按这个公式输出了结果。

\vspace{0.8em}\hrule\vspace{0.8em}

\section{[链接] 五、小算子与"算子融合"}

RMSNorm、残差相加、SiLU、逐元素乘\dots{}\dots{}这些算子每个元素只做几次运算，AI 远小于 1。如果每个都单独启动一个 kernel：

\begin{center}
\small
\begin{tabularx}{\textwidth}{p{2.8cm} X}
\toprule
\textbf{执行模式} & \textbf{显存（HBM）与片上缓存（SRAM）数据流动过程} \\
\midrule
\textbf{朴素未融合算子} & 读 $x$ (HBM) $\to$ RMSNorm $\to$ 写回 HBM $\to$ 读回 $\to$ 乘 $\gamma$ $\to$ 写回 $\to$ 读回 $\to$ 残差加法 $\to$ 写回 HBM（\textbf{访存往返 6 次！}） \\
\midrule
\textbf{Kernel Fusion 融合} & 读 $x$ (HBM) $\to$ \textbf{片上 SRAM / 寄存器一次性完成计算} $\to$ 写回 HBM（\textbf{访存仅 1 次，打满带宽！}） \\
\bottomrule
\end{tabularx}
\end{center}

大部分时间花在显存来回搬运上。\textbf{算子融合（Kernel Fusion）}把多个逐元素操作合成一个 kernel：数据只从显存读一次，在寄存器里做完所有运算，再写回一次。
\begin{itemize}
  \item \texttt{torch.compile}、Triton、TensorRT 做的核心优化之一就是自动融合；
  \item \textbf{FlashAttention 本质上就是把"QK$^T$ $\rightarrow$ 掩码 $\rightarrow$ Softmax $\rightarrow$ 乘 V"整个融合成一个 kernel}，而且巧妙地让 $S \times S$ 的注意力矩阵根本不落到显存上——这需要一个数学技巧（Online Softmax），见第 9 篇。
\end{itemize}

\vspace{0.8em}\hrule\vspace{0.8em}

\section{[实验] 六、浮点格式：量化与混合精度的前提}

\subsection{浮点数的结构}

一个浮点数由三部分组成：符号位 $s$、指数位 $e$（决定\textbf{范围}）、尾数位 $m$（决定\textbf{精度}）：

\[
x = (-1)^s \times 2^{e - \text{bias}} \times (1.m)_2
\]

\begin{center}
\small
\begin{tabularx}{\textwidth}{l X X X X }
\toprule
\textbf{格式} & \textbf{符号/指数/尾数} & \textbf{最大值} & \textbf{相对精度（机器 $\varepsilon$）} & \textbf{主要用途} \\
\midrule
FP32 & 1 / 8 / 23 & $3.4 \times 10^{38}$ & $2^{-23} \approx 1.2 \times 10^{-7}$ & 累加器、优化器状态 \\
FP16 & 1 / 5 / 10 & \textbf{65504} & $2^{-10} \approx 9.8 \times 10^{-4}$ & 早期混合精度 \\
\textbf{BF16} & 1 / \textbf{8} / 7 & $3.4 \times 10^{38}$ & $2^{-7} \approx 7.8 \times 10^{-3}$ & 现代训练与推理的默认格式 \\
FP8 E4M3 & 1 / 4 / 3 & 448 & $2^{-3}$ & 前向的权重与激活 \\
FP8 E5M2 & 1 / 5 / 2 & 57344 & $2^{-2}$ & 梯度（需要更大范围） \\
FP4 E2M1 & 1 / 2 / 1 & 6 & — & Blackwell（含你的 5060 Ti）原生支持，需配合分块缩放 \\
\bottomrule
\end{tabularx}
\end{center}

\subsection{为什么 BF16 打败了 FP16？}

BF16 就是把 FP32 的尾数\textbf{直接砍掉 16 位}：指数位和 FP32 一样多，\textbf{范围与 FP32 相同}，只是精度低。
FP16 精度更高，但最大只能表示 65504，训练时激活或梯度很容易\textbf{溢出成 inf}，需要"损失缩放（loss scaling）"等额外技巧。深度学习对\textbf{范围}比对\textbf{精度}敏感得多，所以 BF16 胜出。

\subsection{低精度的陷阱：累加必须用高精度}

BF16 只有 8 位有效精度（7 位尾数 + 隐含的 1）。在 256 附近，相邻两个可表示数之间的间隔是 2：

\[
256 + 1 \overset{\text{BF16}}{=} 256
\]

如果你把 4096 个数在 BF16 里逐个累加，后面加的小数会被\textbf{全部吞掉}。所以：
\begin{itemize}
  \item Tensor Core 做 BF16 矩阵乘时，\textbf{内部用 FP32 累加}；
  \item Softmax、RMSNorm 的求和通常在 FP32 下完成（Lab 07b 的实现里你会看到 \texttt{.float()}）。
\end{itemize}

配套实验会让你亲眼看到这个现象。

\subsection{整数格式与浮点格式的区别}

INT8/INT4 是\textbf{均匀网格}（相邻值间隔固定为缩放因子 $S$），浮点是\textbf{对数网格}（越靠近 0 越密）。
神经网络的权重大多集中在 0 附近，所以同样 8 位，FP8 对小值更友好、INT8 对均匀分布更友好。这是第 11 篇量化的出发点。

\vspace{0.8em}\hrule\vspace{0.8em}

\section{[OK] 本篇自测}

\begin{enumerate}
  \item $[1, 4096] \times [4096, 4096]$ 的 BF16 矩阵乘，FLOPs 和访存字节各是多少？AI 是多少？
  \item 为什么说 Decode 时"batch 从 1 提到 32，每一步耗时几乎不变"？什么时候这个结论不再成立？
  \item 推导 Decode 注意力的 AI，并说明为什么增大 batch 对它没有帮助。
  \item 在一张带宽 448 GB/s 的卡上，BF16 的 Qwen2.5-7B（约 15.2 GB 权重）batch=1 Decode 的理论上限是多少 token/s？换成 INT4 权重呢？
  \item BF16 和 FP16 都是 16 位，为什么训练更偏爱 BF16？
  \item 为什么 Softmax 的求和要在 FP32 里做？
\end{enumerate}

\vspace{0.8em}\hrule\vspace{0.8em}

\section{[资料] 外部资源}

\begin{itemize}
  \item Horace He：\href{https://horace.io/brrr\_intro.html}{Making Deep Learning Go Brrrr From First Principles} ——算力/带宽/开销三分法，\textbf{必读}。
  \item NVIDIA：\href{https://docs.nvidia.com/deeplearning/performance/dl-performance-gpu-background/index.html}{GPU Performance Background User's Guide} ——官方的算术强度与 Roofline 讲解。
  \item kipply：\href{https://kipp.ly/transformer-inference-arithmetic/}{Transformer Inference Arithmetic} ——用本篇的方法把推理的每一项算清楚，\textbf{必读}。
  \item Google DeepMind：\href{https://jax-ml.github.io/scaling-book/}{How to Scale Your Model} ——第 1 章 Rooflines、第 4 章 Transformer Math、第 7 章 Inference，是本篇最好的进阶读物。
  \item {}\href{https://github.com/gpu-mode/lectures}{GPU MODE 讲座} ——从 CUDA 入门到 FlashAttention 的社区课程；配合教材 \textit{Programming Massively Parallel Processors}（PMPP，第 4 版）。
  \item {}\href{https://docs.nvidia.com/cuda/cuda-c-programming-guide/}{NVIDIA CUDA C++ Programming Guide} 第 1\textasciitilde{}5 章（编程模型、存储层级）。
  \item Goldberg：\href{https://docs.oracle.com/cd/E19957-01/806-3568/ncg\_goldberg.html}{What Every Computer Scientist Should Know About Floating-Point Arithmetic}（经典，选读）。
  \item Williams, Waterman, Patterson 2009：\textit{Roofline: An Insightful Visual Performance Model for Multicore Architectures}（CACM，Roofline 的原始论文）。
\end{itemize}

\vspace{0.8em}\hrule\vspace{0.8em}

\textbf{上一篇}：\textbf{第 7 篇：大模型的决策大脑：Logits、温度系数极限证明与采样算法} ｜ \textbf{下一篇}：\textbf{第 9 篇：FlashAttention：Online Softmax 的严格推导与 IO 复杂度} ｜ \textbf{总路线}：\textbf{LEARNING\_PATH.md}

\chapter{FlashAttention：Online Softmax 的严格推导与 IO 复杂度}

\begin{warningBox}[避坑指南与核心警示]
\textbf{阅读前须知}：FlashAttention（2022）可能是近几年影响最大的一个推理/训练优化——它\textbf{不改变任何数学结果}，却让注意力快了数倍、显存从 $O(N^2)$ 降到 $O(N)$，从而让长上下文成为可能。它的核心是一个非常漂亮的代数技巧：\textbf{Online Softmax}。本篇把它从头推导一遍。

\textbf{前置}：第 2 篇（Softmax 平移不变性）、第 4 篇（注意力公式）、第 8 篇（存储层级、算术强度、算子融合）。

\textbf{配套实验}：\textbf{\texttt{lab09\_flash\_attention/flash\_attention\_numpy.py}}——用 NumPy 实现 Online Softmax、分块 FlashAttention 与 FlashDecoding 的 split-K 合并，逐一验证与朴素注意力完全相等；并在 GPU 上实测朴素注意力与 PyTorch SDPA 的显存差距。
\textbf{真机版}：\textbf{\texttt{lab09\_flash\_attention/cuda\_softmax/}}——把本节和下一节的 Online Softmax 写成\textbf{真正的 CUDA C++ kernel}，在你自己显卡上量出 DRAM 带宽、L2 命中与占用率。不需要装 CUDA toolkit。
\end{warningBox}

\vspace{0.8em}\hrule\vspace{0.8em}

\section{[低速] 一、朴素注意力慢在哪里？}

对序列长度 $N$、头维度 $d$，朴素实现分三步，每一步都是一个独立的 kernel：

\begin{center}
\small
\begin{tabularx}{\textwidth}{l X X}
\toprule
\textbf{标准 Attention 步骤} & \textbf{数学算子公式} & \textbf{朴素显存读写量（HBM IO）} \\
\midrule
1. 计算打分矩阵 & $S = Q K^T / \sqrt{d}$ & 读 $Q, K$ ($2Nd$) $\longrightarrow$ 写 $S$ 到显存 ($N^2$) \\
2. 概率归一化 & $P = \mathrm{softmax}(S)$ & 读 $S$ ($N^2$) $\longrightarrow$ 写 $P$ 到显存 ($N^2$) \\
3. 加权求和输出 & $O = P V$ & 读 $P$ ($N^2$), $V$ ($Nd$) $\longrightarrow$ 写 $O$ 到显存 ($Nd$) \\
\midrule
\textbf{总访存量} & \multicolumn{2}{l}{\textbf{高达 $4Nd + 2N^2$ 字节（受 $O(N^2)$ 显存访问严重瓶颈束缚）}} \\
\bottomrule
\end{tabularx}
\end{center}

\begin{itemize}
  \item \textbf{显存}：必须存下 $N \times N$ 的 $S$ 和 $P$。$N = 32\text{K}$ 时单个头就是 $10^9$ 个元素，BF16 下 2GB——而一层有几十个头。
  \item \textbf{访存}：总共读写 $\Theta(Nd + N^2)$ 个元素。而 $d$ 通常只有 64\textasciitilde{}128，$N^2$ 项完全主导。Softmax 这一步每个元素只做几次运算，算术强度极低（第 8 篇），\textbf{GPU 大部分时间在搬运 $N^2$ 的中间矩阵}。
\end{itemize}

\textbf{自然的想法}：像第 8 篇说的那样做算子融合——把 Q、K、V 分块读进片上 SRAM，算完直接输出 O，$S$ 和 $P$ 永远不落到显存。
\textbf{障碍}：Softmax 需要\textbf{整行}的最大值和总和，才能归一化任何一个元素。只看到一个块时，你不知道这一行后面还有没有更大的值。

\vspace{0.8em}\hrule\vspace{0.8em}

\section{[推导] 二、Online Softmax：一遍扫描算出 Softmax}

\subsection{传统的"安全 Softmax"要扫三遍}

对一行 $x_1, \dots, x_N$（第 2 篇）：
\begin{enumerate}
  \item 第一遍求 $m = \max_i x_i$；
  \item 第二遍求 $\ell = \sum_i e^{x_i - m}$；
  \item 第三遍输出 $p_i = e^{x_i - m} / \ell$。
\end{enumerate}

\subsection{把前两遍合成一遍}

维护两个"运行中"的量：到目前为止看过的最大值 $m_j$ 和对应的和 $\ell_j$：

\[
m_j = \max(m_{j-1}, x_j), \qquad \ell_j = \ell_{j-1} \cdot e^{m_{j-1} - m_j} + e^{x_j - m_j}
\]

初始 $m_0 = -\infty$，$\ell_0 = 0$。

\textbf{定理}：对所有 $j$，$\ell_j = \sum_{i=1}^{j} e^{x_i - m_j}$。

\textbf{证明（归纳法）}：
\begin{itemize}
  \item $j = 1$：$m_1 = x_1$，$\ell_1 = 0 + e^{x_1 - x_1} = 1 = \sum_{i=1}^{1} e^{x_i - m_1}$ [OK]
  \item 假设 $\ell_{j-1} = \sum_{i=1}^{j-1} e^{x_i - m_{j-1}}$，则
\end{itemize}

\[
\ell_j = \left(\sum_{i=1}^{j-1} e^{x_i - m_{j-1}}\right) e^{m_{j-1} - m_j} + e^{x_j - m_j} = \sum_{i=1}^{j-1} e^{x_i - m_j} + e^{x_j - m_j} = \sum_{i=1}^{j} e^{x_i - m_j} \quad \blacksquare
\]

关键就是那个\textbf{修正因子} $e^{m_{j-1} - m_j}$：当最大值变大时，之前累加的和是以旧的最大值为基准的，乘上这个因子就把它们"换算"到新基准下。因为 $m_j \ge m_{j-1}$，这个因子 $\le 1$，\textbf{永远不会溢出}。

上面是逐个元素更新；换成逐块更新（一次处理一个块 $B$）形式完全一样：

\[
m_{\text{new}} = \max\left(m_{\text{old}}, \max_{i \in B} x_i\right), \qquad \ell_{\text{new}} = \ell_{\text{old}} \, e^{m_{\text{old}} - m_{\text{new}}} + \sum_{i \in B} e^{x_i - m_{\text{new}}}
\]

\vspace{0.8em}\hrule\vspace{0.8em}

\section{[加速] 三、把输出也"在线"算出来：FlashAttention 的核心}

注意力的一行输出是 $o = \sum_i p_i v_i = \frac{1}{\ell} \sum_i e^{x_i - m} v_i$（$x_i$ 是这一行第 $i$ 个分数 $q \cdot k_i / \sqrt{d}$）。
我们连最后的归一化都不急着做，只维护一个\textbf{未归一化的累加器}：

\[
\text{acc}_j = \text{acc}_{j-1} \cdot e^{m_{j-1} - m_j} + \sum_{i \in B_j} e^{x_i - m_j} \, v_i
\]

\textbf{定理}：处理完第 $j$ 块后，$\text{acc}_j = \sum_{i \in B_1 \cup \dots \cup B_j} e^{x_i - m_j} v_i$。证明与上一节完全相同（把标量 1 换成向量 $v_i$）。
于是全部块处理完后：

\[
o = \frac{\text{acc}_{\text{final}}}{\ell_{\text{final}}}
\]

\textbf{这和一次性算完整的 Softmax 再乘 V 在数学上严格相等}，不是近似。

\subsection{FlashAttention 算法（前向，单个头）}

\begin{lstlisting}[language=Python]
把 Q 按行切成若干块 Q_1..Q_Tr（每块 Br 行），K、V 切成 K_1..K_Tc（每块 Bc 行）
for 每个 Q 块 Q_i（不同的 Q 块互相独立 → 分给不同的 SM 并行）:
    把 Q_i 读进 SRAM；初始化 m = -∞, l = 0, acc = 0（都在 SRAM / 寄存器里）
    for 每个 K/V 块 j:
        把 K_j、V_j 读进 SRAM
        S_ij = Q_i K_jᵀ / √d                    # Br × Bc 的小矩阵，只存在片上
        (因果掩码：若 K_j 整块都在 Q_i 之后，直接跳过；对角块内部再逐元素掩码)
        m_new = max(m, rowmax(S_ij))
        P_ij  = exp(S_ij - m_new)
        l     = l · exp(m - m_new) + rowsum(P_ij)
        acc   = acc · exp(m - m_new) + P_ij V_j
        m     = m_new
    O_i = acc / l，写回显存；另外存下 logsumexp = m + log l（训练的反向传播要用）
\end{lstlisting}

（这是 FlashAttention-2 的循环顺序：外层循环 Q，内层循环 K/V。FlashAttention-1 的外层是 K/V。）

\subsection{它省下了什么？}

\begin{itemize}
  \item \textbf{显存}：$S$、$P$ 从不写入显存，额外空间只有每行的 $m$、$\ell$：$O(N)$。
  \item \textbf{访存}：Q、O 各读写一次；K、V 对每个 Q 块读一次，共 $N / B_r$ 次。设片上 SRAM 大小为 $M$（按元素计），块大小 $B_r \sim M / d$，论文证明 FlashAttention 的 HBM 访问量为 $\Theta\left(\frac{N^2 d^2}{M}\right)$，而朴素实现为 $\Theta(Nd + N^2)$。
\end{itemize}

  两者之比的量级是 $\frac{M}{d^2}$——\textbf{这是一个与 $N$ 无关的常数}。
  按配套实验 3 的简化模型（约 100 KB SRAM 要同时放下 Q、K、V、O 四个块），$d = 64$ 约省 6 倍，$d = 128$ 约省 1.5 倍。
\begin{noteBox}[核心导读与背景]
注：实验 3 用的是\textbf{元素个数}口径，它数出的是 $2B_r/d = M/(2d^2)$，和论文的 $M/d^2$ 差一个 2。这不是矛盾——实验 3 的朴素项数了 4 趟 $N^2$（写 $S$、读 $S$、写 $P$、读 $P$），而论文的 $\Theta(Nd + N^2)$ 只按元素数算一趟。\textbf{量级和"与 $N$ 无关"这个结论都不受影响}，但引用具体倍数时要说清用的是哪套账。
另外，小 $N$ 时这个常数还到不了（尾块占比大）：$d=64$ 时要到 $N \gtrsim 16\text{K}$ 才收敛到 6 倍，$N=1024$ 时只有 4.9 倍。真实 kernel 的块形状与 L2 缓存命中也会让具体数字不同，但"常数倍、与 $N$ 无关"这一点不变。
\end{noteBox}

\begin{itemize}
  \item \textbf{计算量}：FLOPs 与朴素完全相同（甚至略多一点点 exp 修正）。
  \item \textbf{随 $N$ 增长的真正收益是显存}：朴素实现需要 $O(N^2)$ 的额外显存，FlashAttention 只要 $O(N)$。再加上多个 kernel 融合成一个、不再产生 FP32 的 $N^2$ 中间张量，实测加速往往比访存模型预测的还大（配套实验 5 在 5060 Ti 上 $N = 4096$ 时约 16 倍）。
\end{itemize}

\begin{tipBox}[核心直觉与思考]
{}[思考] \textbf{一个反直觉的教训}：FlashAttention 没有减少任何计算，只是\textbf{减少了数据搬运}，就换来了数倍加速。这正是第 8 篇 Roofline 的核心观点——在访存受限的区域，优化的对象是字节而不是 FLOPs。它的反向传播更进一步：\textbf{宁可重新计算 $S$ 和 $P$，也不去读存下来的 $N^2$ 矩阵}，因为重算比读显存更便宜。
\end{tipBox}

\textbf{想亲手验证这句话？} \textbf{\texttt{lab09\_flash\_attention/cuda\_softmax/}} 把本节的 Online Softmax 写成真正的 CUDA C++ kernel，在你自己显卡上量出：三个实现\textbf{都贴着 380 GB/s 的 Roofline}（也就是都"优化到位了"），差别只在搬运了 268 MB 还是 134 MB —— 约 1.9x。里面还有一个反直觉的结果：把整行缓存进 shared memory 的版本在 $N = 8192$ 时\textbf{反而掉到 74\%}，因为 shared 占用把 SM 的常驻 block 数从 6 压到 2。\textbf{这正是上面说的"分块"为什么必须存在。}

\vspace{0.8em}\hrule\vspace{0.8em}

\section{[组件] 四、FlashDecoding：把同一个技巧用在 Decode 上}

\subsection{问题}

Decode 时每个请求的 Q 只有 \textbf{1 行}。FlashAttention 靠"不同的 Q 块分给不同的 SM"来并行，现在只剩 batch $\times$ 头数 这么多个独立任务。batch = 1、32 个头、上下文 64K 时，只有 32 个任务，而 H100 有 132 个 SM——大部分 SM 在闲着，每个任务却要顺序扫完 64K 个 key。

\subsection{解法：沿 K/V 的序列维度切分（split-K）}

把长度为 $S$ 的 KV Cache 切成 $s$ 段，\textbf{每段独立}算出自己的 $(m^{(r)}, \ell^{(r)}, \text{acc}^{(r)})$，分给不同的 SM 并行。最后用 Online Softmax 的规则\textbf{合并}：

\[
m = \max_r m^{(r)}, \qquad \ell = \sum_r \ell^{(r)} e^{m^{(r)} - m}, \qquad \text{acc} = \sum_r \text{acc}^{(r)} e^{m^{(r)} - m}, \qquad o = \text{acc} / \ell
\]

\textbf{为什么能这样合并？} 因为第三节的定理说明 $\text{acc}^{(r)} = \sum_{i \in \text{段}r} e^{x_i - m^{(r)}} v_i$，乘以 $e^{m^{(r)} - m}$ 后正好变成 $\sum_{i \in \text{段}r} e^{x_i - m} v_i$，所有段相加就是完整的和。
这个合并操作\textbf{满足结合律}，所以切成几段、按什么顺序合并都不影响结果——这也是为什么它能用在分布式场景（例如把超长上下文的 KV 分到多张卡上，Ring Attention）。

\subsection{与 PagedAttention 的关系}

PagedAttention（第 12 篇）把 KV Cache 分成固定大小的物理块，块在显存里\textbf{不连续}。它的 kernel 本质上就是"一块一块读 KV、用 Online Softmax 累加"——\textbf{Online Softmax 让注意力天然可以分块计算，这是分页 KV Cache 能高效工作的数学前提}。

\vspace{0.8em}\hrule\vspace{0.8em}

\section{五、版本演进一览}

\begin{center}
\small
\begin{tabularx}{\textwidth}{p{3cm} p{4.5cm} X}
\toprule
\textbf{版本} & \textbf{年份} & \textbf{关键改进} \\
\midrule
Online Softmax（Milakov \& Gimelshein） & 2018 & 一遍扫描求 Softmax 的归一化因子 \\
FlashAttention & 2022 & 分块 + Online Softmax + 反向重计算，IO 感知 \\
FlashAttention-2 & 2023 & 交换内外循环、减少非矩阵乘运算、更好的并行划分，约 2 倍提速 \\
FlashDecoding & 2023 & Decode 时沿 KV 序列切分并行 \\
FlashAttention-3 & 2024 & 针对 Hopper：异步数据搬运（TMA）、warp 专门化、FP8 \\
\bottomrule
\end{tabularx}
\end{center}

在 PyTorch 里，\texttt{F.scaled\_dot\_product\_attention} 会自动选择 FlashAttention 等融合 kernel；vLLM、SGLang 则使用 FlashAttention / FlashInfer 等专门的推理注意力库。

\vspace{0.8em}\hrule\vspace{0.8em}

\section{[OK] 本篇自测}

\begin{enumerate}
  \item 朴素注意力的显存和 HBM 访存量各是多少量级？为什么 $d$ 很小时问题特别严重？
  \item 写出 Online Softmax 的更新公式，并用归纳法证明 $\ell_j = \sum_{i \le j} e^{x_i - m_j}$。
  \item 修正因子 $e^{m_{\text{old}} - m_{\text{new}}}$ 为什么永远不会溢出？
  \item FlashAttention 的 FLOPs 比朴素实现少吗？那它为什么快？
  \item FlashDecoding 解决的是什么场景下的什么问题？写出两段部分结果的合并公式。
  \item 为什么说 Online Softmax 是 PagedAttention 的数学前提？
\end{enumerate}

\vspace{0.8em}\hrule\vspace{0.8em}

\section{[资料] 外部资源}

\begin{itemize}
  \item Zihao Ye：\href{https://courses.cs.washington.edu/courses/cse599m/23sp/notes/flashattn.pdf}{From Online Softmax to FlashAttention}（华盛顿大学 CSE 599M 讲义，推导最清晰的短文，\textbf{必读}）
  \item Milakov \& Gimelshein 2018：\href{https://arxiv.org/abs/1805.02867}{Online normalizer calculation for softmax}
  \item Dao et al. 2022：\href{https://arxiv.org/abs/2205.14135}{FlashAttention}；Dao 2023：\href{https://arxiv.org/abs/2307.08691}{FlashAttention-2}；Shah et al. 2024：\href{https://arxiv.org/abs/2407.08608}{FlashAttention-3}
  \item Stanford CRFM 博客：\href{https://crfm.stanford.edu/2023/10/12/flashdecoding.html}{Flash-Decoding for long-context inference}
  \item Triton 官方教程：\href{https://triton-lang.org/main/getting-started/tutorials/06-fused-attention.html}{Fused Attention}（用 Python 写出真正在 GPU 上跑的 FlashAttention；先做完前面的矩阵乘教程）
  \item 代码：\href{https://github.com/Dao-AILab/flash-attention}{Dao-AILab/flash-attention}；\href{https://github.com/gpu-mode/lectures}{GPU MODE 讲座} 中关于 FlashAttention 的几讲
\end{itemize}

\vspace{0.8em}\hrule\vspace{0.8em}

\textbf{上一篇}：\textbf{第 8 篇：硬件视角：GPU 存储层级、FLOPs、算术强度与浮点格式} ｜ \textbf{下一篇}：\textbf{第 10 篇：投机解码：无偏性证明与期望加速比} ｜ \textbf{总路线}：\textbf{LEARNING\_PATH.md}

\chapter{投机解码：无偏性证明与期望加速比}

\begin{noteBox}[核心导读与背景]
\textbf{阅读前须知}：Lab 10 给出了投机采样的算法与仿真，但最关键的"为什么"没有讲：它凭什么说"\textbf{输出分布与大模型完全一致}"？加速比能不能事先算出来？什么时候它反而不划算？

\textbf{前置}：第 7 篇（采样）、第 8 篇（访存受限、拐点）。

\textbf{配套实验}：\textbf{\texttt{lab10\_speculative\_decoding/verify\_unbiasedness.py}}——蒙特卡洛验证无偏性定理与期望加速比公式，并演示两种常见错误写法如何产生偏差。
\end{noteBox}

\vspace{0.8em}\hrule\vspace{0.8em}

\section{[加速] 一、为什么"猜 K 个再一起验证"会更快？}

第 8 篇推出：Decode 时线性层的算术强度约等于\textbf{这一步处理的 token 数}。batch=1 时一次只处理 1 个 token，GPU 算力利用率只有百分之一左右。
如果让大模型\textbf{一次前向验证 K 个 token}，读的权重还是那一份，算术强度变成 K，只要 K 远小于拐点（5060 Ti 上约 130），\textbf{耗时几乎不变}。

于是问题变成：\textbf{从哪里弄来 K 个"大概率正确"的候选 token？}——让一个小得多的草稿模型先猜。关键是：猜错了不能影响结果的正确性。

\section{[几何] 二、单个 token 的无偏性定理}

设在同一个前缀下，目标（大）模型给出分布 $p(x)$，草稿（小）模型给出分布 $q(x)$。算法：
\begin{enumerate}
  \item 从 $q$ 采样一个 token $x$；
  \item 以概率 $\min\left(1, \frac{p(x)}{q(x)}\right)$ \textbf{接受} $x$；
  \item 若拒绝，则从\textbf{残差分布} $r(x) = \frac{\max(0, \ p(x) - q(x))}{Z}$ 重新采样，其中 $Z = \sum_{x'} \max(0, \ p(x') - q(x'))$。
\end{enumerate}

（约定：若 $q(x) = 0$，草稿根本不会提出 $x$，第 2 步的比值不会被计算，所以不会出现除以 0。定理对 $q$ 唯一的要求是它是一个合法的概率分布。）

\textbf{定理}：最终输出的 token 服从分布 $p$，与草稿模型 $q$ 的好坏无关。

\textbf{证明}：输出 $x$ 有两种途径——草稿恰好是 $x$ 且被接受；或者草稿被拒绝、然后从残差分布抽中了 $x$。

(1) 草稿是 $x$ 且被接受的概率：

\[
q(x) \cdot \min\left(1, \frac{p(x)}{q(x)}\right) = \min\big(q(x), \ p(x)\big)
\]

(2) 发生拒绝的总概率（对所有可能的草稿 token 求和）：

\[
P(\text{拒绝}) = 1 - \sum_{x'} \min(q(x'), p(x')) = \sum_{x'} \Big( p(x') - \min(q(x'), p(x')) \Big) = \sum_{x'} \max\big(0, \ p(x') - q(x')\big) = Z
\]

（第二个等号用了 $\sum_{x'} p(x') = 1$；第三个等号是因为 $p - \min(p, q) = \max(0, p - q)$。）

(3) 拒绝后从残差分布抽中 $x$ 的概率：

\[
P(\text{拒绝}) \cdot r(x) = Z \cdot \frac{\max(0, \ p(x) - q(x))}{Z} = \max\big(0, \ p(x) - q(x)\big)
\]

两条途径相加：

\[
P(\text{输出} = x) = \min\big(q(x), p(x)\big) + \max\big(0, \ p(x) - q(x)\big) = p(x) \quad \blacksquare
\]

（最后一步分两种情况验证：若 $p(x) \le q(x)$，是 $p(x) + 0$；若 $p(x) > q(x)$，是 $q(x) + p(x) - q(x)$。若 $Z = 0$，说明 $p = q$，永远不会拒绝，结论同样成立。）

\textbf{直觉}：草稿"猜多了"的地方（$q > p$）按比例拒掉，"猜少了"的地方（$p > q$）由残差分布精确补上。

\section{[链接] 三、推广到 K 个 token}

草稿模型自回归地猜 $x_1, \dots, x_K$；大模型\textbf{一次前向}同时得到 $p(\cdot \mid \text{前缀})$、$p(\cdot \mid \text{前缀}, x_1)$、\dots{}\dots{}、$p(\cdot \mid \text{前缀}, x_1..x_K)$ 共 $K+1$ 个分布。然后\textbf{从左到右}逐个用第二节的规则检验：
\begin{itemize}
  \item 第 $i$ 个被接受 $\rightarrow$ 继续检验第 $i+1$ 个；
  \item 第 $i$ 个被拒绝 $\rightarrow$ 从残差分布采样一个 token 替代它，\textbf{丢弃后面所有草稿}，本轮结束；
  \item $K$ 个全部被接受 $\rightarrow$ 用已经算好的第 $K+1$ 个分布\textbf{免费再采样一个} token。
\end{itemize}

由于每个位置都在"正确的前缀"下应用了单 token 定理，按位置归纳可得：\textbf{整段输出的联合分布与大模型逐个 token 自回归采样完全一致}。

\begin{noteBox}[核心导读与背景]
贪婪解码（温度 0）是特例：$p$ 是 one-hot，接受规则退化为"草稿 token 等于大模型的 argmax 才接受"。
\end{noteBox}

\section{[增长] 四、期望加速比公式}

设每个草稿 token 的\textbf{接受率}为 $\alpha$（假设各位置独立同分布）。由第二节，

\[
\alpha = \sum_x \min(p(x), q(x)) = 1 - \text{TV}(p, q)
\]

其中 TV 是两个分布的\textbf{总变差距离}——草稿越像目标，接受率越高。

一轮能产出多少 token？"至少前 $i$ 个草稿都被接受"的概率是 $\alpha^i$，每轮还一定会多出 1 个 token（要么是拒绝后的替代 token，要么是全接受后的免费 token）：

\[
\mathbb{E}[\text{每轮产出}] = 1 + \sum_{i=1}^{K} \alpha^i = \frac{1 - \alpha^{K+1}}{1 - \alpha}
\]

设草稿模型单步耗时与大模型单步耗时之比为 $c$（$c < 1$，例如 $c = 0.1$ 即十分之一），一轮耗时 $= K c \, T + T$（K 次草稿 + 1 次验证）。相对于大模型逐个生成（每 token 耗时 $T$）的加速比：

\[
\text{加速比} = \frac{1 - \alpha^{K+1}}{(1 - \alpha)(Kc + 1)}
\]

例：$\alpha = 0.8$、$K = 4$、$c = 0.1$：期望每轮产出 $\frac{1 - 0.8^5}{0.2} \approx 3.36$，加速比 $\frac{3.36}{1.4} \approx 2.4$ 倍。

\textbf{两个重要推论}：
\begin{enumerate}
  \item $K$ 不是越大越好：$\alpha^i$ 衰减很快，而草稿成本随 $K$ 线性增加，存在最优 $K$（配套实验会扫描出来）。
  \item \textbf{batch 越大，投机解码越不划算}：它的前提是"验证 K 个 token 和验证 1 个一样快"，而 batch 很大时 Decode 本身已接近拐点（第 8 篇），多验证的 token 不再免费。所以投机解码主要用于\textbf{低延迟、小 batch} 的场景。
\end{enumerate}

\section{[启蒙] 五、演进}

\begin{center}
\small
\begin{tabularx}{\textwidth}{p{3.5cm} X}
\toprule
\textbf{方法} & \textbf{草稿从哪来} \\
\midrule
Blockwise Parallel Decoding（2018） & 模型自己加多个预测头（最早的思想） \\
Speculative Decoding / Sampling（2022/2023） & 独立的小模型，并给出无偏性证明 \\
Medusa（2024） & 在大模型上加几个轻量解码头，一次猜多个位置，用树形注意力并行验证多条候选 \\
EAGLE（2024） & 在\textbf{特征层}（倒数第二层隐藏状态）上做自回归草稿，接受率更高 \\
Prompt Lookup / N-gram & 直接从提示词里找重复片段当草稿（代码编辑、摘要中效果很好），零成本 \\
\bottomrule
\end{tabularx}
\end{center}

\vspace{0.8em}\hrule\vspace{0.8em}

\section{[OK] 本篇自测}

\begin{enumerate}
  \item 写出投机采样单 token 无偏性的完整证明。关键的两步等式各用到了什么？
  \item 接受率 $\alpha$ 与总变差距离是什么关系？
  \item 推导每轮期望产出 token 数 $\frac{1-\alpha^{K+1}}{1-\alpha}$。为什么 batch 大时投机解码收益变小？
  \item "拒绝后直接从 $p$ 重新采样"为什么是错的？（提示：Lab 10 实验 3）
\end{enumerate}

\vspace{0.8em}\hrule\vspace{0.8em}

\section{[资料] 外部资源}

\textbf{投机采样}
\begin{itemize}
  \item Leviathan et al. 2022：\href{https://arxiv.org/abs/2211.17192}{Fast Inference from Transformers via Speculative Decoding}；Chen et al. 2023：\href{https://arxiv.org/abs/2302.01318}{Accelerating LLM Decoding with Speculative Sampling}（两篇独立提出，证明方式与本篇一致）
  \item Stern et al. 2018：\href{https://arxiv.org/abs/1811.03115}{Blockwise Parallel Decoding}；Cai et al. 2024：\href{https://arxiv.org/abs/2401.10774}{Medusa}；Li et al. 2024：\href{https://arxiv.org/abs/2401.15077}{EAGLE}
\end{itemize}

\vspace{0.8em}\hrule\vspace{0.8em}

\textbf{上一篇}：\textbf{第 9 篇：FlashAttention：Online Softmax 的严格推导与 IO 复杂度} ｜ \textbf{下一篇}：\textbf{第 11 篇：量化：从每位 6 dB 到 SmoothQuant、AWQ 与 GPTQ} ｜ \textbf{总路线}：\textbf{LEARNING\_PATH.md}

\chapter{量化：从每位 6 dB 到 SmoothQuant、AWQ 与 GPTQ}

\begin{mathBox}[核心数学定理推导]
\textbf{阅读前须知}：Lab 11 给出了量化公式与分组量化的仿真，但没有讲：大模型为什么\textbf{特别难}量化？什么时候该量化权重、什么时候该量化激活？GPTQ、AWQ、SmoothQuant 分别用什么数学手段解决问题？

\textbf{前置}：第 3 篇（矩阵乘的梯度与形状）、第 8 篇（访存受限、拐点、浮点格式）、Lab 11 的 \texttt{quant\_basics.py}。

\textbf{配套实验}：\textbf{\texttt{lab11\_quantization/advanced\_quant\_numpy.py}}——量化误差的 6 dB/bit 规律、激活离群值、SmoothQuant、AWQ 缩放、GPTQ 误差补偿，全部用 NumPy 从零实现。
\end{mathBox}

\vspace{0.8em}\hrule\vspace{0.8em}

\section{一、量化误差的基本规律：每多 1 位，信噪比 +6 dB}

均匀量化（Lab 11 的公式）把一个区间切成 $2^b$ 格，格宽为 $S$。舍入误差近似在 $[-S/2, \ S/2]$ 上均匀分布，方差为：

\[
\sigma_{\text{err}}^2 = \frac{S^2}{12}
\]

（均匀分布 $U(-a, a)$ 的方差是 $a^2/3$，代入 $a = S/2$。）
位数每加 1，格宽 $S$ 减半，误差方差变为 $1/4$，信噪比提高 $10 \log_{10} 4 \approx$ \textbf{6.02 dB}。INT8 到 INT4 理论上损失约 24 dB——这就是为什么 INT4 必须配合分组量化和下面的算法，而 INT8 往往直接可用。

\begin{warningBox}[避坑指南与核心警示]
{}[注意] 这条"6 dB 规律"依赖三个前提：(1) 量化区间固定（这里由最大绝对值决定，不随位数变）；(2) \textbf{高分辨率近似}——格子足够细，每个格子里数据的分布近似平坦，误差才近似均匀分布；(3) 信号功率不变。
位数很低时前提 (2) 不再成立：高斯分布的数据集中在 0 附近，几个粗格子根本覆盖不了分布形状，误差不再是均匀的。所以 Lab 11 实测（高斯数据、最大绝对值定标）里，高位宽时每位约 +6.1 dB，INT3$\rightarrow$INT4 却有 +7.35 dB，INT8$\rightarrow$INT4 共 25.2 dB 而非 24 dB。
\end{warningBox}

误差方差正比于 $S^2$，而 $S$ 由该组数据的\textbf{最大绝对值}决定：\textbf{一个离群值会拉大整组的 $S$}，让组里其它正常值的精度全部变差。这是理解后面所有算法的钥匙。

\section{[瓶颈] 二、大模型为什么难量化：激活离群值}

Dettmers 等（LLM.int8(), 2022）发现：模型规模超过约 6.7B 后，\textbf{激活值}里会出现"涌现的离群特征"——少数几个\textbf{固定的通道（隐藏维度）}，数值比其它通道大几十倍，而且几乎在每个 token 上都出现。
\begin{itemize}
  \item 如果对激活做 per-tensor 量化，整个张量的 $S$ 被这几个通道决定，其它通道几乎全被舍入成 0；
  \item 而\textbf{权重}的分布通常比较规整，容易量化。
\end{itemize}

LLM.int8() 的解法是\textbf{混合精度分解}：把离群通道拿出来用 FP16 算，其余用 INT8 算。效果好但实现复杂、速度不理想，于是有了下面的方法。

\subsection{先想清楚：你要量化什么？}

第 8 篇的 Roofline 决定了答案：

\begin{center}
\small
\begin{tabularx}{\textwidth}{p{2.5cm} p{3.5cm} X X}
\toprule
\textbf{方案} & \textbf{量化对象} & \textbf{加速了什么} & \textbf{代表} \\
\midrule
\textbf{W4A16 / W8A16}（仅权重） & 权重存成低位，计算前反量化成 BF16 & \textbf{Decode}（访存受限，少读字节 = 加速） & GPTQ、AWQ \\
\textbf{W8A8 / FP8}（权重 + 激活） & 两者都低精度，直接用 INT8 / FP8 Tensor Core & \textbf{Prefill} 与大 batch（算力受限，低精度算力翻倍） & SmoothQuant、FP8 \\
\textbf{KV Cache 量化} & K、V 存成 FP8 / INT8 / INT4 & 长上下文 Decode（读 KV 的字节减少），能放更多请求 & FP8 KV、KIVI \\
\bottomrule
\end{tabularx}
\end{center}

\section{[基础] 三、SmoothQuant：把激活的难度"转移"给权重}

对线性层 $Y = XW$（$X$：$[T, C_{in}]$，$W$：$[C_{in}, C_{out}]$），对每个输入通道 $j$ 选一个缩放因子 $s_j > 0$：

\[
Y = XW = \big(X \, \text{diag}(s)^{-1}\big) \big(\text{diag}(s) \, W\big) = \hat{X} \hat{W}
\]

\textbf{这是严格的数学恒等式}。选择

\[
s_j = \frac{\max |X_{:,j}|^{\alpha}}{\max |W_{j,:}|^{1-\alpha}}, \qquad \alpha \approx 0.5
\]

离群通道的激活被除以一个大数、变得温和；对应的权重行被乘以同一个数、变大一些。两边的"量化难度"被拉平，两者都能用 INT8 表示。

\textbf{为什么没有额外开销？} $X$ 来自前面的 RMSNorm，$\text{diag}(s)^{-1}$ 可以直接\textbf{合并进 RMSNorm 的权重 $\gamma$}；$\text{diag}(s)$ 合并进 $W$。推理时的计算图完全没有变化。

\section{[目标] 四、AWQ：保护 1\% 的关键权重}

AWQ 的观察：权重里只有很少一部分（约 1\%）是"关键"的——它们对应的\textbf{输入激活很大}，舍入误差会被放大。把这 1\% 保留为 FP16 效果很好，但混合精度对硬件不友好。
\textbf{替代方案：把关键通道的权重放大 $s$ 倍再量化，同时把对应激活缩小 $s$ 倍}（与 SmoothQuant 同样的恒等变换）：

\[
Q(w \cdot s) \cdot \frac{x}{s}
\]

量化的绝对误差上界由格宽 $S$ 决定，大致不变（只要这个通道不是组内最大值，$S$ 就不变），但它乘的激活缩小到了 $1/s$，所以\textbf{这个通道贡献的输出误差约缩小为 $1/s$}。
AWQ 用少量校准数据，在 $s = (\text{激活平均幅度})^{\beta}$ 这一族里网格搜索 $\beta \in [0, 1]$，选输出误差最小的那个。
它不需要反向传播、不依赖校准数据的具体内容，泛化性好，是目前 W4A16 最常用的方法之一。

\section{[工具] 五、GPTQ：用二阶信息补偿舍入误差}

\subsection{先看一个只有两个权重的例子}

设某个输出 $y = w_1 x_1 + w_2 x_2$，权重 $w_1 = 0.34,\ w_2 = 0.10$，量化格宽 $S = 0.25$（只能取 0、$\pm$0.25、$\pm$0.5\dots{}\dots{}）。再假设两个输入\textbf{总是相等}：$x_1 = x_2 = x$（极端的"相关"）。真实输出 $y = 0.44x$。
\begin{itemize}
  \item \textbf{RTN（各自四舍五入）}：$w_1 \to 0.25$，$w_2 \to 0$，输出 $0.25x$，误差 $0.19x$。
  \item \textbf{GPTQ 的做法}：先量化 $w_1 \to 0.25$，误差 $0.09$；因为 $x_1 = x_2$，把这 $0.09$ 加到还没量化的 $w_2$ 上，$w_2 = 0.19$；再量化 $w_2 \to 0.25$，输出 $0.5x$，误差 $0.06x$。
  \item 单看权重，GPTQ 把 $w_2$ 从 0.10 改成了 0.25，改得更多；但输出误差从 $0.19x$ 降到 $0.06x$。下面的公式就是把"该往哪列补、补多少"在一般相关结构下算到最优。
\end{itemize}

\subsection{目标}

逐层求解：让量化后的层在校准数据 $X$ 上输出尽可能不变（下面把权重写成 $[C_{out}, C_{in}]$，$X$ 为 $[C_{in}, n]$）：

\[
\min_{\hat{W}} \ \| W X - \hat{W} X \|_F^2
\]

这是一个关于 $\hat{W}$ 的二次函数，它的 Hessian 是 $H = 2 X X^T$（$[C_{in}, C_{in}]$，每一行输出共享同一个 $H$）。

\subsection{核心思想（来自 1993 年的 Optimal Brain Surgeon）}

\textbf{一个一个地量化输入通道（列）}。量化第 $q$ 列时产生误差 $w_q - \text{quant}(w_q)$；因为输入通道之间是\textbf{相关的}，可以调整\textbf{还没量化的}其它列来补偿这个误差。最优补偿量为：

\[
\delta = -\frac{w_q - \text{quant}(w_q)}{[H^{-1}]_{qq}} \cdot [H^{-1}]_{q,:}
\]

然后把第 $q$ 列从问题中移除（对 $H^{-1}$ 做一步高斯消元），继续下一列。

\textbf{直觉}：如果第 3 个输入通道和第 7 个输入通道高度相关，第 3 列舍入多了一点，就让第 7 列少一点，输出几乎不变。\textbf{RTN（直接四舍五入）每一列各自为战；GPTQ 让后面的列替前面的列"还债"。}

\subsection{让它能用在 175B 模型上的工程技巧}

GPTQ 论文的贡献在于把 OBS 做得足够快：所有行共享同一个 $H^{-1}$、按固定顺序量化、批量延迟更新、用 Cholesky 分解保证数值稳定。175B 模型在单卡上约 4 小时量化完。

\section{[OK] 六、小结：选哪个？}

\begin{center}
\small
\begin{tabularx}{\textwidth}{p{3.5cm} X}
\toprule
\textbf{场景} & \textbf{推荐} \\
\midrule
单卡本地部署、batch 小（例如你的 5060 Ti） & W4A16：AWQ / GPTQ / llama.cpp 的 GGUF Q4\_K 等 \\
数据中心、大 batch、H100 / Blackwell & FP8（W8A8），可加 FP8 KV Cache \\
超长上下文 & 在权重量化基础上加 KV Cache 量化 \\
\bottomrule
\end{tabularx}
\end{center}

评估量化好坏：先看 PPL 变化（第 3 篇），再看下游任务准确率——PPL 对小损失不够敏感。

\vspace{0.8em}\hrule\vspace{0.8em}

\section{[OK] 本篇自测}

\begin{enumerate}
  \item 为什么每多 1 位信噪比提高约 6 dB？
  \item Decode 为什么适合"只量化权重"，而 Prefill 需要"权重 + 激活都量化"才能加速？
  \item 为什么说 SmoothQuant 的缩放"没有推理开销"？
  \item AWQ 放大关键通道的权重，为什么反而减小了误差？
  \item GPTQ 相比 RTN，利用了输入通道之间的什么性质？为什么它的权重误差更大、输出误差却更小？
\end{enumerate}

\vspace{0.8em}\hrule\vspace{0.8em}

\section{[资料] 外部资源}

\begin{itemize}
  \item Maarten Grootendorst：\href{https://newsletter.maartengrootendorst.com/p/a-visual-guide-to-quantization}{A Visual Guide to Quantization}（图解入门，先看这个）
  \item Dettmers et al. 2022：\href{https://arxiv.org/abs/2208.07339}{LLM.int8()}（离群值的发现）
  \item Xiao et al. 2022：\href{https://arxiv.org/abs/2211.10438}{SmoothQuant}
  \item Lin et al. 2023：\href{https://arxiv.org/abs/2306.00978}{AWQ}
  \item Frantar et al. 2022：\href{https://arxiv.org/abs/2210.17323}{GPTQ}；以及它的前身 \href{https://arxiv.org/abs/2208.11580}{Optimal Brain Compression}
  \item Liu et al. 2024：\href{https://arxiv.org/abs/2402.02750}{KIVI (KV Cache 2-bit 量化)}
\end{itemize}

\vspace{0.8em}\hrule\vspace{0.8em}

\textbf{上一篇}：\textbf{第 10 篇：投机解码：无偏性证明与期望加速比} ｜ \textbf{下一篇}：\textbf{第 12 篇：推理服务系统：PagedAttention、前缀缓存与调度} ｜ \textbf{总路线}：\textbf{LEARNING\_PATH.md}

\chapter{推理服务系统：PagedAttention、前缀缓存与调度}

\begin{noteBox}[核心导读与背景]
\textbf{阅读前须知}：前面九篇都在讲"\textbf{一个请求}怎么算得又对又快"。真实的推理服务同时面对成百上千个请求：它们长短不一、随时到达、随时结束，还常常共享同样的系统提示词。本篇讲推理引擎（vLLM、SGLang、TensorRT-LLM）在\textbf{系统层面}解决的问题：怎么衡量好坏、显存怎么管、请求怎么排。

\textbf{前置}：第 6 篇（KV Cache 显存公式）、Lab 12a（连续批处理仿真）、第 8 篇（Decode 耗时模型 $T_{\text{step}}(B)$）、第 9 篇（Online Softmax 可分块计算）。

\textbf{配套实验}：\textbf{\texttt{lab12b\_paged\_attention/paged\_kv\_cache\_sim.py}}——实现块分配器、块表、分页注意力、写时复制、前缀缓存与 Chunked Prefill 的仿真。
\end{noteBox}

\vspace{0.8em}\hrule\vspace{0.8em}

\section{一、先定义"好"：推理服务的指标体系}

\begin{center}
\small
\begin{tabularx}{\textwidth}{p{2.5cm} p{3.5cm} X X}
\toprule
\textbf{指标} & \textbf{含义} & \textbf{用户感受} & \textbf{主要受谁影响} \\
\midrule
\textbf{TTFT}（Time To First Token） & 从请求到达，到第一个 token 返回 & "它开始回答了吗" & 排队时间 + Prefill 时间 \\
\textbf{TPOT / ITL}（Time Per Output Token / Inter-Token Latency） & 相邻两个输出 token 的间隔 & "打字速度够不够流畅" & 每一步 Decode 的耗时 \\
\textbf{端到端延迟} & $\text{TTFT} + \text{TPOT} \times (n_{\text{out}} - 1)$ & 总等待时间 & 两者之和 \\
\textbf{吞吐（Throughput）} & 每秒生成的总 token 数 / 完成的请求数 & 老板关心的成本 & batch 大小、显存利用率 \\
\textbf{Goodput} & 满足延迟目标（SLO，如 TTFT < 1s 且 TPOT < 50ms）的请求吞吐 & 真正有意义的吞吐 & 以上全部 \\
\bottomrule
\end{tabularx}
\end{center}

工业界通常看 \textbf{P50 / P99 分位数}，而不是平均值——一次偶发的 2 秒卡顿会被用户记住。

\subsection{Little 定律：并发数从哪来}

排队论的 Little 定律：\textbf{系统中平均请求数 = 到达率 $\times$ 平均停留时间}，$L = \lambda W$。
例：每秒到达 10 个请求，每个请求平均要生成 10 秒 $\rightarrow$ 系统里同时约有 \textbf{100 个}请求在跑。每个请求上下文 2K token、每 token KV 128 KB（Llama-3-8B），光 KV Cache 就要 $100 \times 2048 \times 128\text{KB} \approx 25$ GB。
\textbf{并发数由业务决定，显存必须装得下它}——这就是为什么 KV Cache 的显存管理是推理系统的核心问题。

\vspace{0.8em}\hrule\vspace{0.8em}

\section{[反向] 二、迭代级调度回顾：为什么连续批处理成立}

Lab 12a 已经仿真过：静态批处理要等整批最长的请求结束，连续批处理在\textbf{每一步}都可以让结束的请求离开、新请求加入（Orca, 2022）。
它能成立的物理基础是第 8 篇的推导：\textbf{Decode 线性层的算术强度 $\approx$ batch size}，在 batch 远小于拐点时，一步的耗时几乎与 batch 无关。所以"多塞一个请求进这一步"几乎是免费的。

但连续批处理立刻带来两个新问题：
\begin{enumerate}
  \item 请求随时进出，\textbf{KV Cache 的显存怎么分配和回收？}（第三\textasciitilde{}五节）
  \item 新请求的 Prefill 很重，插进正在 Decode 的批次里会怎样？（第六节）
\end{enumerate}

\vspace{0.8em}\hrule\vspace{0.8em}

\section{[模块] 三、KV Cache 的显存碎片：PagedAttention 要解决的问题}

\subsection{传统做法：按最大长度预留连续显存}

请求到达时，并不知道它最终会生成多少 token，于是只能按 \texttt{max\_len}（例如 2048）预留一整块\textbf{连续}显存：

\begin{figure}[htbp]
\centering
\begin{tikzpicture}[scale=0.95]
  \node[anchor=west, font=\sffamily\small\bfseries] at (0, 2.5) {传统连续预留 (按 max\_len=2048 预分配，严重浪费)：};
  
  \node[anchor=east, font=\sffamily\footnotesize] at (2.2, 1.8) {请求 A (实际用 300):};
  \draw[fill=primary!70, draw=primary!90] (2.4, 1.5) rectangle (4.2, 2.1) node[midway, white, font=\tiny\bfseries] {已用 300};
  \draw[fill=codegray!15, draw=codegray!60, dashed] (4.2, 1.5) rectangle (11.0, 2.1) node[midway, codegray, font=\footnotesize] {内部碎片 (未用浪费 1748)};
  
  \node[anchor=east, font=\sffamily\footnotesize] at (2.2, 0.8) {请求 B (实际用 1200):};
  \draw[fill=primary!70, draw=primary!90] (2.4, 0.5) rectangle (7.8, 1.1) node[midway, white, font=\tiny\bfseries] {已用 1200};
  \draw[fill=codegray!15, draw=codegray!60, dashed] (7.8, 0.5) rectangle (11.0, 1.1) node[midway, codegray, font=\footnotesize] {内部碎片 (848)};
  
  \draw[<->, thick, darktext] (2.4, 0.1) -- (11.0, 0.1) node[midway, below, font=\scriptsize\sffamily] {连续预留长度 max\_len = 2048 tokens};
\end{tikzpicture}
\caption{传统连续显存分配的内部碎片问题}
\end{figure}

\begin{itemize}
  \item \textbf{内部碎片}：预留了却没用上的部分；
  \item \textbf{外部碎片}：不同大小的连续块反复分配释放后，空闲显存被切成零碎的小段，总量够但凑不出一整块。
\end{itemize}

vLLM 论文测量发现，当时的系统里 KV Cache 显存\textbf{只有 20\%\textasciitilde{}38\% 真正存了 token}。

\subsection{PagedAttention：把操作系统的虚拟内存搬过来}

操作系统早在几十年前就解决过同样的问题——\textbf{分页}：

\begin{center}
\small
\begin{tabularx}{\textwidth}{p{3.5cm} X}
\toprule
\textbf{操作系统} & \textbf{PagedAttention} \\
\midrule
进程 & 请求（序列） \\
虚拟页 & 逻辑块（第 0 块 = 第 0\textasciitilde{}15 个 token） \\
物理页框 & 显存里的物理块（每块存 16 个 token 的 K、V） \\
页表 & \textbf{块表（Block Table）}：逻辑块号 $\rightarrow$ 物理块号 \\
缺页时分配 & 写满一块时才分配下一块 \\
\bottomrule
\end{tabularx}
\end{center}

\begin{center}
\small
\begin{tabularx}{\textwidth}{l X}
\toprule
\textbf{请求 ID} & \textbf{逻辑块到物理块映射表 (Block Table)} \\
\midrule
\textbf{请求 A} & 逻辑块 0 $\to$ 物理块 7, \quad 逻辑块 1 $\to$ 物理块 2, \quad 逻辑块 2 $\to$ 物理块 9 \\
\textbf{请求 B} & 逻辑块 0 $\to$ 物理块 4, \quad 逻辑块 1 $\to$ 物理块 0, \quad 逻辑块 2 $\to$ 物理块 5 \\
\midrule
\textbf{核心优势} & \textbf{按需分配固定大小的物理块 (如 16 tokens)，物理显存无需连续，内部碎片近乎为零！} \\
\bottomrule
\end{tabularx}
\end{center}

\begin{itemize}
  \item 每个请求最多只浪费\textbf{最后一个块里没写满的部分}（< 16 个 token）。注意这只是\textbf{上界}：$\frac{15}{S} < 4\%$ 要求序列长度 $S \gtrsim 400$ token；短对话（$S=100$ 时最坏 15\%）远达不到这个数；
  \item 没有外部碎片：所有块一样大，任何空闲块都能用；
  \item 结果：同样的显存能容纳 \textbf{2\textasciitilde{}4 倍} 的并发请求，吞吐相应提升。
\end{itemize}

\subsection{代价：注意力 kernel 要"查表读 KV"}

KV 不再连续，注意力计算时要按块表逐块读取——而第 9 篇已经证明，Online Softmax 让注意力可以\textbf{逐块累加}，所以分页几乎不影响计算效率。这是 PagedAttention 在数学上可行的原因。

\subsection{显存不够时怎么办：抢占}

请求在运行中不断变长，块可能被分完。此时调度器\textbf{抢占}一部分请求：
\begin{itemize}
  \item \textbf{换出（Swap）}：把它的 KV 块拷到 CPU 内存，之后再换回；
  \item \textbf{重算（Recompute）}：直接丢掉它的 KV，之后重新 Prefill（Prefill 是算力受限的，对短序列往往比经 PCIe 拷贝更快）。
\end{itemize}

\vspace{0.8em}\hrule\vspace{0.8em}

\section{四、共享：引用计数与写时复制}

同一个提示词需要生成多个回答时（并行采样 $n > 1$、束搜索），它们的提示词 KV \textbf{完全相同}。有了块表，共享变得非常自然：
\begin{itemize}
  \item 多个序列的块表指向\textbf{同一组物理块}，每个物理块维护一个\textbf{引用计数}；
  \item 当某个序列要往一个\textbf{被共享且未写满}的块里追加 token 时，先\textbf{复制一份}给自己再写（\textbf{写时复制，Copy-on-Write}）——这同样是操作系统 \texttt{fork()} 的经典技术。
\end{itemize}

并行采样 4 个回答时，提示词部分的 KV 只存 1 份而不是 4 份。

\vspace{0.8em}\hrule\vspace{0.8em}

\section{五、前缀缓存：跨请求复用 KV}

很多请求的\textbf{开头}是一样的：
\begin{itemize}
  \item 所有请求共享同一段几千 token 的\textbf{系统提示词}；
  \item 多轮对话中，第 $k$ 轮的输入 = 前 $k-1$ 轮的全部内容 + 新问题；
  \item Few-shot 示例、RAG 中被反复检索到的同一篇文档、Agent 的工具描述。
\end{itemize}

这些前缀的 KV 完全相同（同样的 token、同样的位置 $\rightarrow$ 同样的 K、V，RoPE 也一样）。\textbf{前缀缓存}让请求结束后 KV 块不立即释放，而是留着等下一个相同前缀的请求复用：
\begin{itemize}
  \item \textbf{vLLM 的自动前缀缓存}：对每个\textbf{写满的块}，用"它自己的 token + 它之前所有 token"算一个哈希值作为键；新请求逐块查哈希，命中就直接复用物理块；
  \item \textbf{SGLang 的 RadixAttention}：用\textbf{基数树（Radix Tree）}组织所有缓存的前缀，按 LRU 淘汰。
\end{itemize}

\textbf{收益}：命中的部分\textbf{完全跳过 Prefill}。系统提示词 4000 token、用户问题 50 token 时，TTFT 可能降低一个数量级。

\begin{warningBox}[避坑指南与核心警示]
{}[注意] 前缀必须\textbf{逐 token 完全相同}且从第一个 token 开始。把会变化的内容（例如当前时间）放在系统提示词开头，会让整个缓存失效——这是写提示词时要注意的工程细节。
\end{warningBox}

\vspace{0.8em}\hrule\vspace{0.8em}

\section{六、Prefill 与 Decode 的冲突：Chunked Prefill 与 P/D 分离}

\subsection{问题：一个长 Prefill 会让所有人卡住}

连续批处理把新请求的 Prefill 和其他请求的 Decode 放在同一步里。第 8 篇告诉我们：
\begin{itemize}
  \item Decode 一步耗时 $\approx$ 读一遍权重的时间（例如 10ms）；
  \item 一个 8000 token 的 Prefill 是算力受限的，可能要几百毫秒。
\end{itemize}

它们在同一步里，\textbf{这一步的耗时就由 Prefill 决定}，正在 Decode 的几十个用户会同时看到一次几百毫秒的"卡顿"（TPOT 尖刺）。

\subsection{Chunked Prefill（Sarathi-Serve, 2024）}

把长 Prefill \textbf{切成若干块}（例如每块 512 token），每一步只做一块，和其他请求的 Decode 拼在一起。每一步设一个 \textbf{token 预算}：
\begin{itemize}
  \item 每一步的耗时有了上界，Decode 用户不再卡顿；
  \item 而且 Decode 本来算力闲置（访存受限），顺便塞进一段 Prefill，正好把闲置的算力用上——\textbf{两种瓶颈互补}。
\end{itemize}

\subsection{P/D 分离（DistServe、Mooncake，2024）}

更彻底的做法：Prefill 和 Decode \textbf{放到不同的 GPU 上}，Prefill 算完后把 KV Cache 通过高速网络传给 Decode 节点。
\begin{itemize}
  \item 理由正是第 8 篇：两者的瓶颈完全相反（算力 vs 带宽），硬件配置和并行策略可以分别优化；
  \item 代价：KV Cache 需要跨机传输，对网络带宽要求高。
\end{itemize}

\vspace{0.8em}\hrule\vspace{0.8em}

\section{七、一个现代推理引擎里还有什么}

\begin{center}
\small
\begin{tabularx}{\textwidth}{p{3cm} p{4.5cm} X}
\toprule
\textbf{技术} & \textbf{解决的问题} & \textbf{本仓库对应} \\
\midrule
CUDA Graph & Decode 一步要启动几百个小 kernel，每个都有微秒级开销（Lab 07b 实测只跑到带宽下界的一成左右，约 7\%\textasciitilde{}14\%） & Lab 07b 练习 5 \\
融合的注意力 kernel（FlashAttention / FlashInfer） & 第 9 篇 & Lab 09 \\
量化（权重 / KV Cache） & 显存与带宽 & Lab 11、第 11 篇 \\
投机解码 & Decode 串行、算力闲置 & Lab 10、第 10 篇 \\
张量并行 / 专家并行 & 单卡放不下 & 第 13 篇、Lab 13 \\
结构化输出（JSON Schema 约束解码） & 让输出满足格式 & 第 7 篇的延伸：对 logits 做掩码 \\
多 LoRA 服务 & 一个基座模型服务多个微调版本 & 选读 S-LoRA \\
\bottomrule
\end{tabularx}
\end{center}

\vspace{0.8em}\hrule\vspace{0.8em}

\section{[OK] 本篇自测}

\begin{enumerate}
  \item TTFT 和 TPOT 分别由什么决定？为什么服务要看 P99 而不只看平均？
  \item 用 Little 定律估算：每秒 5 个请求、每个请求平均耗时 20 秒、平均上下文 4K、每 token KV 128 KB，需要多少 KV Cache 显存？
  \item 预留连续显存为什么会浪费？PagedAttention 每个请求最多浪费多少？
  \item 为什么说 Online Softmax（第 9 篇）是 PagedAttention 的数学前提？
  \item 写时复制在什么时候触发？
  \item 为什么把"当前时间"写在系统提示词开头会破坏前缀缓存？
  \item 一个长 Prefill 为什么会让所有正在 Decode 的请求卡顿？Chunked Prefill 如何解决？为什么说它让"两种瓶颈互补"？
\end{enumerate}

\vspace{0.8em}\hrule\vspace{0.8em}

\section{[资料] 外部资源}

\begin{itemize}
  \item Kwon et al. 2023：\href{https://arxiv.org/abs/2309.06180}{Efficient Memory Management for LLM Serving with PagedAttention}（vLLM 论文，\textbf{必读}），以及 \href{https://blog.vllm.ai/2023/06/20/vllm.html}{vLLM 博客}
  \item Yu et al. 2022：\href{https://www.usenix.org/conference/osdi22/presentation/yu}{Orca: A Distributed Serving System for Transformer-Based Generative Models}（OSDI'22，迭代级调度）
  \item Anyscale 博客：\href{https://www.anyscale.com/blog/continuous-batching-llm-inference}{How continuous batching enables 23x throughput in LLM inference}
  \item Zheng et al. 2023：\href{https://arxiv.org/abs/2312.07104}{SGLang (RadixAttention)}
  \item Agrawal et al. 2023/2024：\href{https://arxiv.org/abs/2308.16369}{Sarathi}、\href{https://arxiv.org/abs/2403.02310}{Sarathi-Serve (Chunked Prefill)}
  \item Zhong et al. 2024：\href{https://arxiv.org/abs/2401.09670}{DistServe (P/D 分离)}；Qin et al. 2024：\href{https://arxiv.org/abs/2407.00079}{Mooncake}
  \item vLLM 官方文档：\href{https://docs.vllm.ai/}{docs.vllm.ai} 中的 Design / Architecture 部分
  \item Lilian Weng：\href{https://lilianweng.github.io/posts/2023-01-10-inference-optimization/}{Large Transformer Model Inference Optimization}（推理优化综述）
\end{itemize}

\vspace{0.8em}\hrule\vspace{0.8em}

\textbf{上一篇}：\textbf{第 11 篇：量化：从每位 6 dB 到 SmoothQuant、AWQ 与 GPTQ} ｜ \textbf{下一篇}：\textbf{第 13 篇：多卡并行推理：张量并行、流水线并行与专家并行} ｜ \textbf{总路线}：\textbf{LEARNING\_PATH.md}

\chapter{多卡并行推理：张量并行、流水线并行与专家并行}

\begin{noteBox}[核心导读与背景]
\textbf{阅读前须知}：一张 16GB 的卡能跑 7B（量化后），但 70B 的 BF16 权重就要 140GB，DeepSeek-V3 更是 671B。模型放不下一张卡时，必须把它\textbf{切开}放到多张卡上。怎么切、切完之后卡与卡之间要传什么、传多少——这些决定了多卡推理的速度。

\textbf{前置}：第 3 篇（矩阵乘的形状）、第 5 篇（MoE）、第 8 篇（带宽与耗时模型）、第 9 篇（Online Softmax 合并）。

\textbf{配套实验}：\textbf{\texttt{lab13\_parallelism/tensor\_parallel\_sim.py}}——用 NumPy 模拟多卡：张量并行的 MLP 与注意力、错误切法的反例、Ring All-Reduce、多卡 Decode 延迟模型。
\end{noteBox}

\vspace{0.8em}\hrule\vspace{0.8em}

\section{一、先认识"通信原语"和它们的代价}

多卡之间交换数据的几种标准操作（NCCL 库实现）：

\begin{center}
\small
\begin{tabularx}{\textwidth}{p{3cm} p{4.5cm} X}
\toprule
\textbf{原语} & \textbf{做什么} & \textbf{典型用途} \\
\midrule
\textbf{All-Reduce} & 每张卡有一个同形状张量，结果是\textbf{所有卡的和}，每张卡都拿到 & 张量并行合并部分和 \\
All-Gather & 每张卡有一块，结果是\textbf{拼接}后的完整张量 & 拼回被切开的激活 \\
Reduce-Scatter & 先求和，再把结果切块，每张卡拿一块 & All-Reduce 的前半段 \\
\textbf{All-to-All} & 每张卡把自己的数据切成 $n$ 份分别发给每张卡 & MoE 把 token 发给对应专家 \\
Send / Recv & 点对点 & 流水线并行传激活 \\
\bottomrule
\end{tabularx}
\end{center}

\textbf{Ring All-Reduce 的代价}：$n$ 张卡连成环，先做 Reduce-Scatter 再做 All-Gather，每张卡发送的数据量为

\[
2 \cdot \frac{n-1}{n} \cdot \text{数据大小}
\]

几乎与卡数无关——这是它被广泛采用的原因。但每一次通信还有一个\textbf{固定延迟}（几微秒到几十微秒），数据很小时，延迟占主导。

\textbf{互联带宽差距巨大}：

\begin{center}
\small
\begin{tabularx}{\textwidth}{p{3.5cm} X}
\toprule
\textbf{互联} & \textbf{单向带宽量级} \\
\midrule
NVLink（H100，同机 8 卡） & 约 450 GB/s \\
PCIe 4.0 x16 / 5.0 x16 & 约 25 / 50 GB/s \\
InfiniBand 400 Gb/s（跨机） & 约 50 GB/s \\
\bottomrule
\end{tabularx}
\end{center}

所以"切法"必须和"连接方式"匹配：通信频繁的切法只能用在 NVLink 连接的同一台机器内。

\vspace{0.8em}\hrule\vspace{0.8em}

\section{二、张量并行（Tensor Parallelism, TP）：把每一层的矩阵切开}

Megatron-LM（2019）给出了 Transformer 最经典的切法。以 FFN 为例：$Y = f(X A) B$，$f$ 是逐元素的非线性（GeLU / SiLU）。

\subsection{第一个矩阵按列切}

把 $A$ 按列切成 $[A_1, A_2]$，卡 $i$ 持有 $A_i$：

\[
X A = [X A_1, \ X A_2] \quad \Rightarrow \quad f(XA) = [f(X A_1), \ f(X A_2)]
\]

\textbf{因为 $f$ 是逐元素的}，每张卡可以独立地对自己那半算非线性，\textbf{不需要任何通信}。

\subsection{第二个矩阵按行切}

把 $B$ 按行切成 $\begin{bmatrix} B_1 \\ B_2 \end{bmatrix}$，卡 $i$ 持有 $B_i$：

\[
Y = [f(XA_1), \ f(XA_2)] \begin{bmatrix} B_1 \\ B_2 \end{bmatrix} = f(XA_1) B_1 + f(XA_2) B_2
\]

每张卡算出一个\textbf{部分和}，最后做\textbf{一次 All-Reduce} 求和。

\subsection{为什么不能反过来？}

如果 $A$ 按\textbf{行}切（即切分输入维度），每张卡得到的是 $X_i A_i$ 这样的\textbf{部分和}，而

\[
f(X_1 A_1 + X_2 A_2) \neq f(X_1 A_1) + f(X_2 A_2)
\]

非线性函数对加法不分配，必须先 All-Reduce 求和才能做 $f$——多了一次通信。\textbf{"先列后行"让两次矩阵乘之间的非线性完全在本地完成，每个子层只通信一次。}

SwiGLU 同理：$W_{\text{gate}}$、$W_{\text{up}}$ 按列切，$W_{\text{down}}$ 按行切（$\text{SiLU}(\cdot) \odot (\cdot)$ 是逐元素的）。

\subsection{注意力：按"头"切}

多个注意力头本来就互相独立，这恰好就是"按列切"：$W_q, W_k, W_v$ 按头分给不同的卡（列切），每张卡独立算自己那几个头的注意力，$W_o$ 按行切，最后一次 All-Reduce。
\textbf{附带好处}：每张卡只存自己那些头的 KV Cache，KV 显存也被均分。（GQA 的 KV 头数少于 TP 卡数时，KV 头需要复制到多张卡上。）

\subsection{代价与收益}

\begin{itemize}
  \item \textbf{每层 2 次 All-Reduce}（注意力后一次、FFN 后一次），数据量是 $[\text{token 数}, d_{\text{model}}]$。
  \item \textbf{Decode 收益}：每张卡只读 $1/n$ 的权重，第 8 篇的访存时间直接除以 $n$——\textbf{TP 能降低单请求延迟}，这是它在推理中最重要的价值。
  \item \textbf{Decode 代价}：每步 $2L$ 次 All-Reduce，每次数据量很小（batch $\times$ 8192 $\times$ 2 字节），\textbf{延迟主导}。在 PCIe 上，通信时间可能超过省下的计算时间，TP 加卡反而变慢（配套实验会算给你看）。
\end{itemize}

\vspace{0.8em}\hrule\vspace{0.8em}

\section{[硬件] 三、流水线并行（Pipeline Parallelism, PP）：按层切}

把 80 层分成 4 段，每张卡放 20 层，激活依次在卡之间传递。
\begin{itemize}
  \item \textbf{通信少}：每段边界只传一次激活（点对点），适合跨机器；
  \item \textbf{不降低单请求延迟}：一个 token 仍要依次穿过所有层，第 1 张卡算完第 2 张卡才能开始；
  \item \textbf{气泡}：如果只有一个请求，任一时刻只有一张卡在工作。推理服务里有大量请求时，可以把不同请求（micro-batch）错开塞进流水线，让每张卡都忙起来——\textbf{PP 提升吞吐，不提升延迟}。
\end{itemize}

llama.cpp 在多张消费级显卡之间"按层分配"，本质上就是流水线式切分。

\vspace{0.8em}\hrule\vspace{0.8em}

\section{[人员][实验] 四、专家并行（Expert Parallelism, EP）：MoE 的切法}

MoE 层有几十上百个专家（第 5 篇），把不同的专家放在不同的卡上：
\begin{enumerate}
  \item 路由器算出每个 token 要去哪 $k$ 个专家；
  \item \textbf{All-to-All}：把 token 发到对应专家所在的卡；
  \item 各卡计算自己的专家；
  \item 再一次 \textbf{All-to-All} 把结果发回原卡，按路由权重加权求和。
\end{enumerate}

难点：
\begin{itemize}
  \item \textbf{负载均衡}：热门专家所在的卡会成为瓶颈，训练时要加负载均衡损失或偏置调整（DeepSeek-V3 使用无辅助损失的偏置调整）；
  \item \textbf{All-to-All 通信量大}，需要高带宽网络；DeepSeek 为此专门开发了通信库 DeepEP。
\end{itemize}

\vspace{0.8em}\hrule\vspace{0.8em}

\section{[并行] 五、其它并行方式}

\begin{itemize}
  \item \textbf{数据并行（DP）}：每张卡放一份完整模型，各自服务不同请求。最简单，吞吐线性扩展，但每张卡都要放得下整个模型。
  \item \textbf{序列 / 上下文并行}：超长上下文（百万 token）时，把一个序列的 KV 切到多张卡上。\textbf{Ring Attention} 让 KV 块在卡之间环形传递，每张卡用 \textbf{Online Softmax 的合并公式（第 9 篇 §4）} 累加部分结果——又一次用到了那个满足结合律的合并操作。
\end{itemize}

\section{[指南] 六、怎么组合？一些经验法则}

\begin{center}
\small
\begin{tabularx}{\textwidth}{p{3.5cm} X}
\toprule
\textbf{场景} & \textbf{常见方案} \\
\midrule
单机多卡（NVLink），模型放不下一张卡 & 机内 TP \\
跨多台机器的超大模型 & 机内 TP + 机间 PP \\
MoE 大模型（DeepSeek-V3 等） & 专家用 EP，注意力用 DP 或 TP \\
消费级多卡（PCIe，无 NVLink） & 优先 PP / 按层切分（通信少），或者量化后单卡 \\
吞吐优先、单卡放得下 & 多副本 DP \\
\bottomrule
\end{tabularx}
\end{center}

\vspace{0.8em}\hrule\vspace{0.8em}

\section{[OK] 本篇自测}

\begin{enumerate}
  \item Ring All-Reduce 每张卡发送多少数据？为什么说它"几乎与卡数无关"？
  \item 为什么 FFN 的第一个矩阵按列切、第二个按行切？反过来会怎样？
  \item 张量并行时每层需要几次 All-Reduce？数据量多大？
  \item 为什么 TP 能降低 Decode 延迟而 PP 不能？
  \item 在两张通过 PCIe 连接的 5060 Ti 上跑 TP=2，Decode 一定比单卡快吗？用第 8 篇的耗时模型加上通信项分析。
  \item Ring Attention 和第 9 篇的哪个公式有关？
\end{enumerate}

\vspace{0.8em}\hrule\vspace{0.8em}

\section{[资料] 外部资源}

\begin{itemize}
  \item Shoeybi et al. 2019：\href{https://arxiv.org/abs/1909.08053}{Megatron-LM}（张量并行的原始论文，第 3 节图 3 就是本篇第二节）
  \item Pope et al. 2022：\href{https://arxiv.org/abs/2211.05102}{Efficiently Scaling Transformer Inference}（专门讲\textbf{推理}的多卡切分与延迟/成本权衡，\textbf{强烈推荐}）
  \item Google DeepMind：\href{https://jax-ml.github.io/scaling-book/}{How to Scale Your Model} 第 3 章（Sharding）与第 7\textasciitilde{}8 章（推理）
  \item Hugging Face：\href{https://huggingface.co/spaces/nanotron/ultrascale-playbook}{The Ultra-Scale Playbook}（训练视角，但通信原语与并行方式讲得非常直观）
  \item Huang et al. 2018：\href{https://arxiv.org/abs/1811.06965}{GPipe}（流水线并行）；Liu et al. 2023：\href{https://arxiv.org/abs/2310.01889}{Ring Attention}
  \item DeepSeek-AI 2024：\href{https://arxiv.org/abs/2412.19437}{DeepSeek-V3 技术报告}（EP、负载均衡、部署方案）
  \item {}\href{https://docs.nvidia.com/deeplearning/nccl/user-guide/docs/usage/collectives.html}{NCCL 文档}：各通信原语的定义
\end{itemize}

\vspace{0.8em}\hrule\vspace{0.8em}

\textbf{上一篇}：\textbf{第 12 篇：推理服务系统：PagedAttention、前缀缓存与调度} ｜ \textbf{总路线}：\textbf{LEARNING\_PATH.md}


\part{工业实战与实验源码精解篇 (Practical Labs \& Implementations)}
\chapter{从零手写自动微分与反向传播引擎 (Micrograd)}

\section{Lab 03：从零手写自动求导（Autograd）}

\begin{noteBox}[核心导读与背景]
对应理论：\textbf{第 3 篇：神经网络如何学习：导数、链式法则、反向传播与交叉熵}
前置：读完第 3 篇第一 \textasciitilde{} 六节。
依赖：纯 Python 标准库。
\end{noteBox}

\section{为什么要自己写一遍？}

PyTorch 的 \texttt{loss.backward()} 是整个深度学习最大的"黑盒"。只要你亲手写过一次，就会明白：
\begin{itemize}
  \item 反向传播 = \textbf{链式法则 + 拓扑排序 + 梯度累加}，没有任何魔法；
  \item 为什么每一步训练前必须 \texttt{zero\_grad()}（梯度是 \texttt{+=} 累加的）；
  \item 为什么训练要保存激活值（\texttt{\_backward} 闭包里用到了前向时的 \texttt{self.data}）；
  \item 为什么推理时可以用 \texttt{torch.no\_grad()} / \texttt{torch.inference\_mode()} 省掉大量显存（不建图、不存激活）。
\end{itemize}

\section{运行}

\begin{lstlisting}[language=bash]
python3 lab03_autograd_from_scratch/micrograd_from_scratch.py
\end{lstlisting}

\begin{center}
\small
\begin{tabularx}{\textwidth}{p{3cm} p{4.5cm} X}
\toprule
\textbf{实验} & \textbf{验证的结论} & \textbf{对应章节} \\
\midrule
1 & 手算表格 dL/dw = −12, dL/db = −6 & 第 3 篇 §2.3 \\
2 & 自动求导 = 数值差分（梯度检验） & 第 3 篇 §3 \\
3 & Softmax 雅可比 $p_i(\delta_{ij}-p_j)$；交叉熵梯度 $p - y$；logits 放大后梯度坍缩 & 第 3 篇 §6、第 4 篇 §3 \\
4 & $\partial L/\partial W = X^T G$，$\partial L/\partial X = G W^T$ & 第 3 篇 §4 \\
5 & Sigmoid 连乘梯度消失 vs 残差畅通 & 第 3 篇 §7、第 4 篇 §5 \\
6 & 单层学不会 XOR，两层可以 & 第 1 篇 §4.2、第 0 篇族谱 \\
\bottomrule
\end{tabularx}
\end{center}

\section{动手练习（做完才算过关）}

\begin{enumerate}
  \item 给 \texttt{Value} 加一个 \texttt{silu()} 算子（$x\sigma(x)$），自己推导导数，并用实验 2 的方法做梯度检验。
  \item 把实验 6 的激活函数从 \texttt{tanh} 换成 \texttt{sigmoid}，学习率不变，观察收敛速度变化，并用第 3 篇 §7 的导数上界解释。
  \item 把实验 5 的深度改成 30，看纯 Sigmoid 的梯度变成多少。
  \item （进阶）跟着 \href{https://karpathy.ai/zero-to-hero.html}{Karpathy 的 micrograd 视频} 实现 \texttt{Neuron / Layer / MLP} 类，并训练一个二分类的"月牙"数据集。
\end{enumerate}

\section{学完之后}

你已经可以放心使用 \texttt{torch.autograd} 了——它做的事和这 350 行完全一样，只是算子换成了张量、实现换成了 C++/CUDA。
下一步：\textbf{第 5 篇} $\rightarrow$ \textbf{Lab 05：亲手训练一个迷你 GPT}。


\section{实验核心源码精选与解析}

本实验配套有完整可运行的工业级验证代码，重点实现细节如下：


\subsection{源码实现：\texttt{micrograd\_from\_scratch.py}}

\lstinputlisting[language=Python, caption={\texttt{micrograd\_from\_scratch.py}}]{code/lab03_autograd_from_scratch_micrograd_from_scratch.py}


\chapter{Transformer 推理显微镜：逐层张量变化与 RoPE 实测}

\section{Lab 04：大模型推理的微观全貌（零起点极细拆解）}

如果你对「大模型推理到底在算什么」、「张量每一步怎么变」感到模糊，本模块就是为你量身打造的\textbf{显微镜}。

无论模型是 0.5B、7B 还是 70B，无论叫 Llama-3、Qwen2.5 还是 DeepSeek，其自回归推理的\textbf{微观数学本质完全一致}。

\vspace{0.8em}\hrule\vspace{0.8em}

\section{[指南] 1. 宏观全貌：从一句话到下一个字的 8 个步骤}

假设你对模型说：\texttt{"人工智能"}，模型要推理出下一个字 \texttt{"的"}。整个过程分为严格的 8 步：

\begin{figure}[htbp]
\centering
\begin{tikzpicture}[
  stage/.style={draw=primary!80, fill=primary!5, rounded corners=1.5mm, align=center, font=\small\sffamily, inner sep=4.5pt, line width=0.8pt, minimum width=11cm},
  arr/.style={-{Stealth[scale=0.9]}, draw=primary, line width=1pt}
]
\node[stage] (s0) {原始输入文本：\texttt{"人工智能"}};
\node[stage, below=0.4cm of s0] (s1) {1. Tokenizer 分词 $\rightarrow$ Token IDs: \texttt{[102, 305]} \quad (Shape: $[2]$)};
\node[stage, below=0.4cm of s1] (s2) {2. Embedding 查表 $\rightarrow$ 稠密向量矩阵 \quad (Shape: $[2, d_{\mathrm{model}}]$)};
\node[stage, below=0.4cm of s2] (s3) {3. RMSNorm 归一化 $\rightarrow$ 消除激活值方差偏置 \quad (Shape: $[2, d_{\mathrm{model}}]$)};
\node[stage, below=0.4cm of s3] (s4) {4. 多头自注意力 (Self-Attention + RoPE 旋转位置编码) \quad (Shape: $[2, d_{\mathrm{model}}]$)};
\node[stage, below=0.4cm of s4] (s5) {5. 残差连接 + SwiGLU / MLP 非线性知识映射 \quad (Shape: $[2, d_{\mathrm{model}}]$)};
\node[stage, below=0.4cm of s5, fill=secondary!10, draw=secondary] (s6) {6. 循环堆叠 $L$ 个 Transformer 块 (重复步骤 3 $\sim$ 5)};
\node[stage, below=0.4cm of s6] (s7) {7. Final RMSNorm + LM Head 词表投影 $\rightarrow$ Logits \quad (Shape: $[2, V]$)};
\node[stage, below=0.4cm of s7, fill=accent!10, draw=accent] (s8) {8. 采样算法 (Top-P / Temperature) $\rightarrow$ 生成下一个 Token ID: 520 (\texttt{"的"})};

\draw[arr] (s0) -- (s1);
\draw[arr] (s1) -- (s2);
\draw[arr] (s2) -- (s3);
\draw[arr] (s3) -- (s4);
\draw[arr] (s4) -- (s5);
\draw[arr] (s5) -- (s6);
\draw[arr] (s6) -- (s7);
\draw[arr] (s7) -- (s8);
\end{tikzpicture}
\caption{Transformer 单步推理张量生命周期微观全貌}
\end{figure}

\vspace{0.8em}\hrule\vspace{0.8em}

\section{[观察] 2. 张量形状流（Tensor Shape Flow Tracker）}

理解大模型推理的关键，是死死盯住\textbf{每一步张量的形状（Shape）}。
下表用一个\textbf{假设的}微型配置来说明形状如何流动（注意：脚本 \texttt{01\_tensor\_trace\_step\_by\_step.py} 为了能把每个数字打印出来，用的是更小的配置：$d_{\text{model}} = 16$、2 个头、词表 11 个字、只有 1 层；下面的流程图描述的是真实模型的 $L$ 层结构）。

假设微型模型配置：
\begin{itemize}
  \item 序列长度 $T = 3$（例如输入 3 个词）
  \item 隐藏层维度 $d_{\text{model}} = 64$
  \item 注意力头数 $n_{\text{heads}} = 4$，每个头的维度 $d_{\text{head}} = 64 / 4 = 16$
  \item 词表大小 $\text{vocab\_size} = 1000$
\end{itemize}

\begin{center}
\small
\begin{tabularx}{\textwidth}{l X X X X }
\toprule
\textbf{处理阶段} & \textbf{算子名称} & \textbf{输入形状 (Input Shape)} & \textbf{输出形状 (Output Shape)} & \textbf{物理含义解释} \\
\midrule
\textbf{分词} & Tokenizer & 字符串文本 & \texttt{[3]} & 文本转整数 ID \\
\textbf{词嵌入} & Embedding Lookup & \texttt{[3]} & \texttt{[3, 64]} & 每个词映射成 64 维空间的一个点 \\
\textbf{归一化} & RMSNorm & \texttt{[3, 64]} & \texttt{[3, 64]} & 均方根缩放，保持方差稳定 \\
\textbf{QKV投影} & Linear ($W_q, W_k, W_v$) & \texttt{[3, 64]} & \texttt{[3, 64]} (各 1 个) & 生成查询、键、值向量 \\
\textbf{多头拆分} & Reshape \& Transpose & \texttt{[3, 64]} & \texttt{[4, 3, 16]} & 4 个注意力头并行关注不同维度的特征 \\
\textbf{位置编码} & RoPE (旋转编码) & \texttt{[4, 3, 16]} & \texttt{[4, 3, 16]} & 对 Q 和 K 按位置角频率旋转，注入顺序感 \\
\textbf{注意力打分} & $Q \cdot K^T / \sqrt{d_k}$ & Q:\texttt{[4, 3, 16]}, K:\texttt{[4, 3, 16]} & \texttt{[4, 3, 3]} & 第 $i$ 个词与第 $j$ 个词的关联度矩阵 \\
\textbf{因果掩码} & Causal Mask & \texttt{[4, 3, 3]} & \texttt{[4, 3, 3]} & 将右上角置为 $-\infty$（禁止看到未来的词） \\
\textbf{注意力归一} & Softmax & \texttt{[4, 3, 3]} & \texttt{[4, 3, 3]} & 转换成概率百分比分布（每一行和为 1） \\
\textbf{值加权} & Score $\cdot V$ & Scores:\texttt{[4, 3, 3]}, V:\texttt{[4, 3, 16]} & \texttt{[4, 3, 16]} & 提取汇聚上下文信息 \\
\textbf{多头合并} & Transpose \& Reshape & \texttt{[4, 3, 16]} & \texttt{[3, 64]} & 拼接 4 个头的特征 \\
\textbf{输出投影} & Linear ($W_o$) & \texttt{[3, 64]} & \texttt{[3, 64]} & 特征融合 \\
\textbf{残差连接} & $X + \text{Attn}(X)$ & \texttt{[3, 64]} + \texttt{[3, 64]} & \texttt{[3, 64]} & 保留原始输入，防止梯度/特征退化 \\
\textbf{前馈网络} & SwiGLU / MLP & \texttt{[3, 64]} & \texttt{[3, 64]} & 升维到例如 172 维激活后再降维回 64 维 \\
\textbf{残差连接} & $X + \text{FFN}(X)$ & \texttt{[3, 64]} + \texttt{[3, 64]} & \texttt{[3, 64]} & 第二层残差 \\
\textbf{词表投影} & LM Head ($W_{\text{head}}$) & \texttt{[3, 64]} & \texttt{[3, 1000]} & 将每个位置映射为词表中每个词的得分 \\
\textbf{预测下个词} & Slice Last Position & \texttt{[3, 1000]} & \texttt{[1000]} & \textbf{我们只取最后一个位置（第 3 个词）的输出} \\
\textbf{采样策略} & Greedy / Top-P & \texttt{[1000]} & 整数（标量） & 采样出第 4 个词的 Token ID！ \\
\bottomrule
\end{tabularx}
\end{center}

\vspace{0.8em}\hrule\vspace{0.8em}

\section{[实验] 3. 逐个核心算子的底层细节与数学公式}

\subsection{(1) 为什么注意力机制必须有「因果掩码（Causal Mask）」？}

在文本生成中，我们不能让前面的词“偷看”后面的词。
对于序列长度为 3 的输入：

\[
\text{Mask} = \begin{bmatrix}
0 & -\infty & -\infty \\
0 & 0 & -\infty \\
0 & 0 & 0
\end{bmatrix}
\]

当注意力分数矩阵与该 Mask 相加后，在执行 $\text{Softmax}$ 时：

\[
e^{-\infty} = 0
\]

被掩码的位置权重精确变成 \textbf{0}！这样第 1 个词只能看第 1 个词，第 2 个词只能看第 1 和第 2 个词，保证严格的单向自回归因果性。

\subsection{(2) RMSNorm（均方根归一化）到底算什么？}

现代大模型（LLaMA/Qwen/Mistral）舍弃了传统的 LayerNorm，改用更高效的 RMSNorm：
\begin{itemize}
  \item 不需要计算均值 $\mu$，不需要减均值；
  \item 仅计算向量元素的均方根（Root Mean Square）：
\end{itemize}

\[
\text{RMS}(x) = \sqrt{\frac{1}{d} \sum_{i=1}^{d} x_i^2 + \epsilon}
\]

\begin{itemize}
  \item 归一化公式（$\gamma$ 是可学习的缩放向量）：
\end{itemize}

\[
\bar{x}_i = \frac{x_i}{\text{RMS}(x)} \odot \gamma_i
\]

计算量少了一半，硬件执行速度更快，数值更平稳。

\subsection{(3) SwiGLU（门控前馈网络）算什么？}

现代模型基本不再使用经典的 \texttt{ReLU(W1 x) W2}，而是采用 \textbf{SwiGLU} 结构：
它包含三个权重矩阵：$W_{\text{gate}}, W_{\text{up}}, W_{\text{down}}$：

\[
\text{SwiGLU}(x) = \left( \text{SiLU}(x W_{\text{gate}}) \otimes (x W_{\text{up}}) \right) W_{\text{down}}
\]

\begin{itemize}
  \item $x W_{\text{gate}}$ 通过 $\text{SiLU}$ 激活函数充当“开关门（Gate）”，控制哪些特征信息可以通过。
  \item 与 $x W_{\text{up}}$ 做逐元素乘法（Hadamard product $\otimes$）。
  \item 最终经过 $W_{\text{down}}$ 降维回隐藏层维度。
\end{itemize}

\subsection{(4) 采样：从 Logits 到文字}

最后一步"从 Logits 选出下一个 token"单独放在 \textbf{Lab 07a}，配合第 7 篇学习。

\vspace{0.8em}\hrule\vspace{0.8em}

\section{[执行] 运行显微镜级追踪脚本}

本目录下提供了可以直接运行的 Python 代码：

\begin{enumerate}
  \item \textbf{\textbf{\texttt{01\_tensor\_trace\_step\_by\_step.py}}}：
\end{enumerate}

   \textbf{单步前向传播显微镜}。完整实现了一个微型 Transformer，单步逐行打印所有输入/输出张量 Shape、注意力矩阵数值、Logits 排行榜以及生成循环。
\begin{enumerate}
  \item \textbf{\textbf{\texttt{02\_rope\_and\_attention\_microscope.py}}}：
\end{enumerate}

   \textbf{位置编码与因果掩码显微镜}。无位置编码时的置换等变性、RoPE 的相对位置性质、长程衰减的正确表述（只在 q、k 相关时的期望上成立）、因果掩码。

对应理论：\textbf{第 4 篇}。依赖：纯 Python 标准库。

\vspace{0.8em}\hrule\vspace{0.8em}

\section{[链接] 说明与下一步}

\begin{itemize}
  \item 脚本 01 的 "Decode" 步骤为了看清张量，故意把整段序列重算一遍（没有 KV Cache）；权重是随机的，所以输出没有语义。
  \item 脚本 02 已修正：旧版把"点积交换律"当成了"没有语序感"的证明，并声称单一频率的 RoPE 有长程衰减（实际上单一频率下点积是 $\cos(\Delta\theta)$，是周期函数）。新版演示了真正的\textbf{置换等变性}；长程衰减实验对比了 q、k 完全相关 / 部分相关 / 独立随机三种情况——只有相关时点积的期望才随距离衰减，独立随机时完全不衰减。
  \item 想看一个\textbf{真正训练过}的模型，请继续做 \textbf{Lab 05} 与 \textbf{Lab 07b}。完整顺序见 \textbf{LEARNING\_PATH.md}。
\end{itemize}



\section{实验核心源码精选与解析}

本实验配套有完整可运行的工业级验证代码，重点实现细节如下：


\subsection{源码实现：\texttt{01\_tensor\_trace\_step\_by\_step.py}}

\lstinputlisting[language=Python, caption={\texttt{01\_tensor\_trace\_step\_by\_step.py}}]{code/lab04_transformer_microscope_01_tensor_trace_step_by_step.py}


\subsection{源码实现：\texttt{02\_rope\_and\_attention\_microscope.py}}

\lstinputlisting[language=Python, caption={\texttt{02\_rope\_and\_attention\_microscope.py}}]{code/lab04_transformer_microscope_02_rope_and_attention_microscope.py}


\chapter{从零训练与推理最小完整 Llama 模型}

\section{Lab 05：亲手训练一个迷你 Llama，再在它身上做推理}

\begin{warningBox}[避坑指南与核心警示]
对应理论：第 4 篇（注意力公式）、第 6 篇（KV Cache）、第 7 篇（采样）、第 3 篇（交叉熵 / 困惑度 / Adam）、第 5 篇（RMSNorm、RoPE、SwiGLU、GQA、Pre-LN）
前置：完成 Lab 03（你已经亲手写过反向传播，所以这里可以放心用 \texttt{loss.backward()}）。
依赖：PyTorch（见根目录 \texttt{requirements.txt}，RTX 50 系必须装 cu128 版本）。
\end{warningBox}

\section{为什么需要这个模块？}

\texttt{lab04\_transformer\_microscope} 和 \texttt{lab06\_kv\_cache} 用的都是\textbf{随机权重}，输出没有任何意义。
这里你会得到一个\textbf{真正学过东西}的模型（虽然很小），然后在它身上验证推理理论——这样"KV Cache 输出完全一致""温度越高越胡说"这些结论就不再是抽象描述。

\section{文件}

\begin{center}
\small
\begin{tabularx}{\textwidth}{p{3.5cm} X}
\toprule
\textbf{文件} & \textbf{内容} \\
\midrule
{}\textbf{\texttt{model.py}} & 约 180 行的 Llama 架构：RMSNorm、RoPE（半分配对，与 HF 一致）、GQA、因果掩码、SwiGLU、Pre-LN、KV Cache。每个组件的注释都标注了对应的教程章节 \\
{}\textbf{\texttt{train.py}} & 用本仓库的教程 Markdown 做语料，字符级训练；包含"初始损失 $\approx$ ln V"检查、验证集 PPL、early stopping \\
{}\textbf{\texttt{generate.py}} & 朴素生成 vs KV Cache 逐 token 对拍；用真实缓存张量验证显存公式；逐步计时；采样策略对比 \\
\bottomrule
\end{tabularx}
\end{center}

\section{运行}

\begin{lstlisting}[language=bash]
.venv/bin/python lab05_train_tiny_llama/train.py      # GPU 约 0.5~1.5 分钟（取决于显卡时钟状态）
.venv/bin/python lab05_train_tiny_llama/generate.py
\end{lstlisting}

\section{你会看到什么}

以下结论都由脚本自己\textbf{断言}（不满足会直接报错），具体数值以运行输出为准：

\begin{itemize}
  \item 初始损失 $\approx$ ln(V) ——初始化正确。脚本会打印实测值和理论值，并断言两者相差 < 0.5。
\end{itemize}

  {}[注意] \textbf{V（字表大小）和下面所有 PPL 数字都会随仓库文档的增删而变}，因为语料就是本仓库的文档。
  具体数值以你运行时打印的为准，这里只描述\textbf{形状}。
\begin{itemize}
  \item 验证集 PPL：从约 $V$（均匀瞎猜 = ln V）一路降到十几，随后\textbf{过拟合}：训练损失继续下降
\end{itemize}

  （默认 1000 步结束时约 0.4\textasciitilde{}0.5），验证损失却回升。3.5M 参数 vs 十几万字符，数据量只有
  Chinchilla 比例的千分之一，模型开始背书——这正是第 5 篇讲规模定律时说的"参数和数据要匹配"。
\begin{noteBox}[核心导读与背景]
记住这条曲线的\textbf{形状}（先降后升、训练/验证分叉），不要记具体数字：
评审实测中，同一份代码在语料增加 120 行后，最优 PPL 从 11.1 变成 14.3 —— 不是代码变了，
是验证集的切分变了。\textbf{这也是"评测集必须冻结"这条工程常识的活例子。}
\end{noteBox}

\begin{itemize}
  \item 朴素生成与 KV Cache 生成 400 个 token \textbf{完全一致}，logits 最大差 \textasciitilde{}1e-5
  \item KV Cache 真实字节数 = 公式 $2 \times L \times n_{kv} \times d_{head} \times b \times T$，\textbf{逐字节相等}
  \item 每步耗时（本机一次运行）：
\end{itemize}

\begin{noteBox}[核心导读与背景]
{}[计时] 以下计时类数字是本机某一次运行的\textbf{量级示例}，前提是显卡空闲、已预热；GPU 降频或有其它负载时可能慢 40\% 以上。
\textbf{CPU 列更不可靠}：CPU 首步有首次内存分配/缺页开销，且没有预热路径，实测可能与表里差 5 倍。
请看\textbf{趋势和比例}（GPU 上 Cache 基本持平 vs 朴素线性增长），以你自己运行的输出为准。
\end{noteBox}

\begin{center}
\small
\begin{tabularx}{\textwidth}{l X X X X }
\toprule
\textbf{上下文长度} & \textbf{朴素 (GPU)} & \textbf{Cache (GPU)} & \textbf{朴素 (CPU)} & \textbf{Cache (CPU)} \\
\midrule
11 & 3.1 ms & 3.3 ms & 3.4 ms & 1.9 ms \\
210 & 4.1 ms & 3.1 ms & 10.6 ms & 3.4 ms \\
409 & 3.9 ms & 4.7 ms & 37.2 ms & 2.5 ms \\
\bottomrule
\end{tabularx}
\end{center}

CPU 上完美呈现"朴素线性变慢、缓存持平"；GPU 上两者差不多——因为模型太小，\textbf{耗时被 kernel 启动等固定开销主导}（第 8 篇的 overhead 区间）。这是真实世界的现象，不是实验失败：它解释了为什么推理引擎要用 CUDA Graph，以及为什么 Lab 07b 要换成真实的大模型再测。

\section{动手练习}

\begin{enumerate}
  \item 把 \texttt{n\_kv\_head} 改成 8（即 MHA）重新训练，比较参数量、KV Cache 大小和验证 PPL。
  \item 在 \texttt{model.py} 里把 RoPE 去掉（\texttt{apply\_rope} 直接返回 \texttt{x}）重新训练，看验证 PPL 变差多少——对应 \texttt{lab04\_transformer\_microscope/02} 里"没有位置编码就是词袋"的实验。
  \item 把 \texttt{SwiGLU} 换成 \texttt{ReLU(x W1) W2}，保持参数量相同（隐藏维度要改成多少？见第 5 篇 §5.3）。
  \item 生成长度超过训练时的 \texttt{block=128} 后，文本质量有没有明显下降？这和第 5 篇 RoPE 外推的讨论有什么关系？
  \item （进阶）把 \texttt{Attention.forward} 里的 \texttt{torch.cat} 缓存改成\textbf{预分配}一块 \texttt{[B, n\_kv, max\_len, hd]} 的张量、按位置写入，比较速度。这就是从"动态拼接"走向 PagedAttention（Lab 12b）的第一步。
\end{enumerate}

\section{下一步}

{}\textbf{Lab 07b：不用 transformers 的模型类，亲手实现 Qwen2.5 的前向传播并与官方逐位对拍}


\section{实验核心源码精选与解析}

本实验配套有完整可运行的工业级验证代码，重点实现细节如下：


\subsection{源码实现：\texttt{model.py}}

\lstinputlisting[language=Python, caption={\texttt{model.py}}]{code/lab05_train_tiny_llama_model.py}


\subsection{源码实现：\texttt{train.py}}

\lstinputlisting[language=Python, caption={\texttt{train.py}}]{code/lab05_train_tiny_llama_train.py}


\subsection{源码实现：\texttt{generate.py}}

\lstinputlisting[language=Python, caption={\texttt{generate.py}}]{code/lab05_train_tiny_llama_generate.py}


\chapter{KV Cache 性能基准测试与吞吐加速比实测}

\section{Lab 06：KV Cache 原理与自回归解码基准测试}

\section{[思考] 为什么需要 KV Cache？}

在大语言模型（如 GPT、Llama、Qwen 等 Decoder-only 架构）生成文本时，是\textbf{逐字（Token-by-Token）生成}的：

\begin{lstlisting}[numbers=none]
Prompt: "人工智能的未来是"
Step 1 -> "星"
Step 2 -> "辰"
Step 3 -> "大"
Step 4 -> "海"
\end{lstlisting}

在标准的 Self-Attention 计算中：

\[
\text{Attention}(Q, K, V) = \text{softmax}\left(\frac{Q K^T}{\sqrt{d_k}}\right) V
\]

\subsection{朴素实现的问题（无 KV Cache）：}

生成第 $t$ 个 token 时，把前 $t-1$ 个 token 连同最新 token 一起重新送入模型：
\begin{itemize}
  \item 重新计算前 $t-1$ 个 token 的 $Q, K, V$ 投影矩阵。
  \item 但事实上，\textbf{历史 token 的 Key 和 Value 向量是固定不变的}！
  \item 每次生成都会重复计算历史所有 token 的 $K$ 和 $V$，导致总计算复杂度高达 $O(N^2)$，生成越到后面越慢。
\end{itemize}

\subsection{KV Cache 的优化思路：}

\begin{itemize}
  \item 在 \textbf{Prefill 阶段}（首字），计算出 Prompt 中所有 token 的 $K$ 和 $V$，并保存在显存中（KV Cache）。
  \item 在 \textbf{Decode 阶段}（后续生成），每次\textbf{只输入最新的 1 个 token}：
  \item 只对最新 token 计算 $q_{\text{new}}, k_{\text{new}}, v_{\text{new}}$。
  \item 把 $k_{\text{new}}, v_{\text{new}}$ 拼接到缓存列表中：$K_{\text{all}} = [K_{\text{cache}}, k_{\text{new}}]$。
  \item 用单个查询向量 $q_{\text{new}}$ 与整个 $K_{\text{all}}$ 做点积得到注意力分数，加权聚合 $V_{\text{all}}$。
  \item 这样，单步生成的计算复杂度从 $O(t)$ 矩阵乘降为向量-矩阵乘，避免了大量冗余算力消耗！
\end{itemize}

\vspace{0.8em}\hrule\vspace{0.8em}

\section{[执行] 运行测试}

本目录下提供了 \textbf{\texttt{kv\_cache\_benchmark.py}}，演示：
\begin{enumerate}
  \item 朴素生成 vs KV Cache 生成的数学等价性（验证输出完全一致）。
  \item 每步生成耗时随序列长度增加的变化趋势对比。
\end{enumerate}

\begin{lstlisting}[language=bash]
python3 lab06_kv_cache/kv_cache_benchmark.py
\end{lstlisting}

\textit{(本脚本自适应纯 Python 或 NumPy，即使暂未安装任何第三方库也能直接运行！)}

\vspace{0.8em}\hrule\vspace{0.8em}

\section{[链接] 继续深入}

本模块用的是随机权重的单层注意力。想在\textbf{真正训练过的模型}上验证同样的结论（输出逐 token 一致、缓存字节数与公式逐字节相等、每步耗时对比），请做：
\begin{itemize}
  \item {}\textbf{Lab 05}：自己训练的迷你 Llama（\texttt{generate.py}）；
  \item {}\textbf{Lab 07b}：真实的 Qwen2.5-0.5B，上下文越长 KV Cache 的优势越大（2048 时约快一个数量级）。
\end{itemize}

注意：KV Cache 消除的是对历史 token 重复做投影 / FFN 的冗余，\textbf{注意力本身每步仍要读完整个缓存}，长上下文时这会成为新的瓶颈（第 6 篇 §3 的补充说明、第 8 篇 §4.3）。


\section{实验核心源码精选与解析}

本实验配套有完整可运行的工业级验证代码，重点实现细节如下：


\subsection{源码实现：\texttt{kv\_cache\_benchmark.py}}

\lstinputlisting[language=Python, caption={\texttt{kv\_cache\_benchmark.py}}]{code/lab06_kv_cache_kv_cache_benchmark.py}


\chapter{工业级文本采样策略实操 (Top-K / Top-P / Min-P)}

\section{Lab 07a：采样策略显微镜}

\begin{mathBox}[核心数学定理推导]
对应理论：\textbf{第 7 篇：Logits、温度系数极限证明与采样算法}
依赖：纯 Python 标准库。
\end{mathBox}

\section{运行}

\begin{lstlisting}[language=bash]
python3 lab07a_sampling/sampling_strategies.py
\end{lstlisting}

用一组固定的 Logits，逐一对比：原始分布、低温（T=0.2）、高温（T=2.0）、Top-K、Top-P、重复惩罚，打印每种策略处理后的概率分布条形图。

\section{速查}

模型 LM Head 输出的不是概率，而是一堆未经归一化的实数值，称为 \textbf{Logits}（例如 \texttt{[2.1, -0.5, 8.4, 0.1]}）。
\begin{itemize}
  \item \textbf{贪婪搜索（Greedy Search）}：
\end{itemize}

\[
\text{next\_token} = \arg\max(\text{logits})
\]

  直接选得分最高的那一个（确定性强，但容易重复或死板）。
\begin{itemize}
  \item \textbf{温度系数（Temperature $T$）}：
\end{itemize}

\[
\text{probs} = \text{softmax}\left(\frac{\text{logits}}{T}\right)
\]

\begin{itemize}
  \item $T \to 0$：分布极度尖锐，退化为贪婪搜索；
  \item $T > 1$：分布平缓，低分词也有机会被选中，回答更具创意但更容易胡说八道；
  \item \textbf{Top-P（核采样 Nucleus Sampling）}：
\end{itemize}

  按概率从高到低排序累加，只保留累计概率达到 $P$（如 $0.9$）的头部候选词，截断长尾无意义低概率词，在多样性与合理性之间取得最佳平衡。

\section{下一步}

\begin{itemize}
  \item 在\textbf{真正训练过的模型}上看采样效果：\textbf{Lab 05} 的 \texttt{generate.py} 实验 4；
  \item 手写重复惩罚并与 Hugging Face 官方实现逐 token 对拍：\textbf{Lab 07b} 实验 5。
\end{itemize}



\section{实验核心源码精选与解析}

本实验配套有完整可运行的工业级验证代码，重点实现细节如下：


\subsection{源码实现：\texttt{sampling\_strategies.py}}

\lstinputlisting[language=Python, caption={\texttt{sampling\_strategies.py}}]{code/lab07a_sampling_sampling_strategies.py}


\chapter{纯 Python 零依赖 Qwen 模型真实权重推理}

\section{Lab 07b：亲手实现 Qwen2.5 的推理，并与官方逐位对拍}

\begin{noteBox}[核心导读与背景]
这是整条学习路线里\textbf{"从用库的人变成懂原理的人"}的关键一步。

对应理论：第 4、5、6、7、8 篇全部。
前置：完成 Lab 05（你已经在自己训练的模型上写过同样的结构）。
依赖：PyTorch、transformers、safetensors、huggingface\_hub（见根目录 \texttt{requirements.txt}）。第一次运行会下载约 1GB 权重，国内网络可先 \texttt{export HF\_ENDPOINT=https://hf-mirror.com}。
\end{noteBox}

\section{规则}

\texttt{transformers} 只允许用来做两件事：\textbf{分词}，以及\textbf{作为标准答案和我们对拍}。
模型的前向传播、RoPE、GQA、KV Cache、重复惩罚全部自己写，直接读取 \texttt{model.safetensors} 里的原始张量。核心的 \texttt{QwenFromScratch.forward}（含 RMSNorm / RoPE / GQA 注意力 / SwiGLU / KV Cache）约 80 行。

\section{运行}

\begin{lstlisting}[language=bash]
.venv/bin/python lab07b_qwen_from_scratch/qwen_from_scratch.py
\end{lstlisting}

\section{本机（RTX 5060 Ti）一次运行的关键结果}

实验 1\textasciitilde{}5 的数值是确定性的，并且由脚本断言；实验 6 是计时。
\begin{noteBox}[核心导读与背景]
{}[计时] 实验 6 的计时数字是本机某一次运行的\textbf{量级示例}，前提是显卡空闲、已预热；GPU 降频或有其它负载时可能慢 40\% 以上。请看\textbf{趋势和比例}，以你自己运行的输出为准。
\end{noteBox}

\begin{center}
\small
\begin{tabularx}{\textwidth}{p{3.5cm} X}
\toprule
\textbf{实验} & \textbf{结果} \\
\midrule
1. 读 config.json & GQA 分组 7（14 个 Q 头 / 2 个 KV 头），每 token KV 仅 \textbf{12 KB} \\
2. 参数分布 & 词嵌入 \textbf{28\%}、FFN 64\%、注意力只有 9\%——小模型的词表开销惊人，所以共享 LM Head \\
3. BPE 分词 & "大模型推理的瓶颈是显存带宽。" $\rightarrow$ 11 个 token：\texttt{大/模型/推理/的/瓶颈/是/显/存/带/宽/。} \\
4. logits 对拍（FP32） & 最大误差 \textbf{1.5e-5}，每个位置 top-1 全部一致 \\
5. 生成对拍 & 纯贪婪逐 token 一致；并且发现 HF 的 \texttt{generate(do\_sample=False)} \textbf{偷偷使用了} \texttt{generation\_config.json} 里的 \texttt{repetition\_penalty=1.1}，我们按第 7 篇公式手写惩罚后再次逐 token 一致 \\
6. 性能（BF16） & Prefill 1024 token $\approx$ 80 ms（约 1.25 万 token/s）；Decode 约 20 ms/token，只达到带宽下界 2.6 ms 的\textbf{一成左右}（本机多次运行在 7\%\textasciitilde{}14\% 之间波动，取决于 GPU 频率与负载） \\
6. 朴素 vs Cache & 上下文越长差距越大：128 时两者接近（约 0.7\textasciitilde{}1.7x，此时单步耗时被 kernel 启动等固定开销主导，谁快取决于当时的噪声），2048 时约 \textbf{9\textasciitilde{}11x} \\
\bottomrule
\end{tabularx}
\end{center}

\subsection{为什么 Decode 只达到理论上限的一成左右？}

权重只有 1GB，按带宽算每 token 只要 2.6ms，但我们的实现每一层要启动十几个小 kernel（24 层就是几百个），每个 kernel 启动都有微秒级开销，逐元素算子还在反复读写显存，GQA 的 \texttt{repeat\_interleave} 甚至把 KV 复制了 7 份。
\textbf{这九成左右的差距，就是 vLLM / llama.cpp / TensorRT-LLM 这些推理引擎存在的全部理由}：算子融合、CUDA Graph、专用的 GEMV 与注意力 kernel（不复制 KV 的 GQA）。

\section{动手练习（按难度递增）}

\begin{enumerate}
  \item 把 \texttt{repeat\_interleave} 去掉：把 Q reshape 成 \texttt{[B, n\_kv, group, T, hd]} 直接和 \texttt{[B, n\_kv, 1, S, hd]} 的 K 做广播矩阵乘，测 Decode 耗时变化。（第 8 篇 §4.3 说明了为什么这很重要）
  \item 用 \texttt{torch.nn.functional.scaled\_dot\_product\_attention} 替换手写的注意力，确认结果不变，比较 Prefill 耗时（它内部会调用 FlashAttention，见第 9 篇）。
  \item 加入第 7 篇的 Temperature + Top-P 采样，并用 \texttt{torch.manual\_seed} 与 HF 的 \texttt{generate(do\_sample=True, ...)} 对比分布是否一致（提示：逐 token 对拍需要完全相同的随机数使用顺序，比较难；改为比较多次采样的统计分布）。
  \item 支持 batch > 1：多个不同长度的请求一起 Decode，需要 padding 与注意力掩码。体会为什么这件事很麻烦——这就是 Lab 12a（连续批处理）和 Lab 12b（PagedAttention）要解决的问题。
  \item （进阶）参考 \href{https://github.com/pytorch-labs/gpt-fast}{gpt-fast} 用 \texttt{torch.compile(mode="reduce-overhead")}（内部使用 CUDA Graph）加速 Decode，看能逼近带宽下界的多少。
  \item （进阶）把权重换成 Lab 11 里的分组 INT4 量化再反量化回 BF16，测量 PPL 与生成质量的变化（第 11 篇）。
\end{enumerate}

\section{学完之后}

你已经能读懂任何 Llama 系模型的 \texttt{modeling\_*.py}，也知道性能差距在哪。下一步：
\begin{itemize}
  \item {}\textbf{第 9 篇：FlashAttention} + \textbf{Lab 09}
  \item {}\textbf{第 12 篇：推理服务系统} + \textbf{Lab 12b}
\end{itemize}



\section{实验核心源码精选与解析}

本实验配套有完整可运行的工业级验证代码，重点实现细节如下：


\subsection{源码实现：\texttt{qwen\_from\_scratch.py}}

\lstinputlisting[language=Python, caption={\texttt{qwen\_from\_scratch.py}}]{code/lab07b_qwen_from_scratch_qwen_from_scratch.py}


\chapter{显存占用严谨核算与 GPU 算力 Roofline 实测}

\section{Lab 08：显存开销与 Roofline 性能模型}

在进行大模型推理系统设计或显卡选型时，有两大黄金法则：
\begin{enumerate}
  \item \textbf{显存容量法则}：决定了你「能不能跑得下」（模型参数 + KV Cache + 运行时开销 $\le$ 显存容量）。
  \item \textbf{Roofline 性能模型与访存带宽法则}：决定了你「能跑多快」（自回归解码受限于 GPU 显存带宽）。
\end{enumerate}

\vspace{0.8em}\hrule\vspace{0.8em}

\section{[推导] 1. 显存开销组成详解}

大模型在推理运行时的显存占用主要由以下四部分组成：

\[
\text{Total Memory} = \text{Mem}_{\text{weights}} + \text{Mem}_{\text{KV\_Cache}} + \text{Mem}_{\text{activation}} + \text{Mem}_{\text{CUDA\_overhead}}
\]

\subsection{(1) 模型权重 (Model Weights)}

\begin{itemize}
  \item \textbf{FP16 / BF16}: 每个参数占用 \textbf{2 字节 (Bytes)}。
  \item 例如 7B 模型需要约 $7 \times 2 = 14 \text{ GB}$。
  \item \textbf{INT8 / FP8}: 每个参数占用 \textbf{1 字节}（7B 约 7 GB）。
  \item \textbf{INT4 / AWQ / GPTQ}: 每个参数占用 \textbf{0.5 字节}（7B 约 3.5 GB）。
\end{itemize}

\subsection{(2) KV Cache 显存}

这是动态显存开销的最大头，随并发请求数和序列长度成正比膨胀：
对于采用分组查询注意力（GQA, Grouped-Query Attention）的模型：
\begin{itemize}
  \item 层数 $L$
  \item KV 头数 $n_{\text{kv\_heads}}$，单头维度 $d_{\text{head}}$
  \item 并发请求数 $B$（Batch Size）
  \item 序列总长度 $S$（Prompt 长度 + 已生成长度）
  \item 存储精度字节数 $b$（FP16 为 2）
\end{itemize}

\[
\text{KV Cache (Bytes)} = 2 \times L \times (n_{\text{kv\_heads}} \times d_{\text{head}}) \times S \times B \times b
\]

\begin{tipBox}[核心直觉与思考]
{}[思考] \textbf{MHA vs GQA 的显存革命}：
- 早期模型（如 LLaMA-1 65B）使用 Multi-Head Attention (MHA)，KV 头数等于 Query 头数（例如 64 头）。
- 现代模型（如 Llama-3、Qwen2.5）几乎全部采用 GQA（例如 8 个 KV 头），KV Cache 显存直接暴降为原来的 $1/4 \sim 1/8$！
\end{tipBox}

\subsection{(3) 激活值 (Activations)}

在 Prefill 阶段较大（因为输入序列长），但在 Decode 阶段因为一次只输入 1 个 token，激活值极小（通常仅几十 MB 到几百 MB）。

\subsection{(4) CUDA 上下文与运行时开销}

PyTorch / CUDA runtime、NCCL 通信缓冲区等，通常占用约 $1 \sim 2 \text{ GB}$。

\vspace{0.8em}\hrule\vspace{0.8em}

\section{[高速] 2. Roofline 模型与为什么 Decode 是访存受限？}

\subsection{计算密度（Arithmetic Intensity）}

\[
\text{Arithmetic Intensity (AI)} = \frac{\text{计算量 (FLOPs)}}{\text{访存数据量 (Bytes)}}
\]

GPU 的极限速度受制于：
\begin{itemize}
  \item \textbf{最大算力 (TFLOPs)}
  \item \textbf{显存带宽 (GB/s)}
\end{itemize}

两者的交界点为硬件的转折算术密度 $AI_{\text{roof}} = \frac{\text{Peak FLOPs}}{\text{Peak Bandwidth}}$。

\begin{center}
\small
\begin{tabularx}{\textwidth}{l X X X X X }
\toprule
\textbf{阶段} & \textbf{输入规模} & \textbf{计算量} & \textbf{访存量} & \textbf{算术强度 $AI$} & \textbf{性能瓶颈} \\
\midrule
\textbf{Prefill 阶段} & $N$ 个 token 并行 & $O(N \cdot P)$ 大矩阵乘法 & 读取一次模型权重 $P$ & 很高 ($> 100$) & \textbf{算力受限 (Compute-Bound)}，GPU 计算核心满载 \\
\textbf{Decode 阶段} & 每次仅 1 个 token & 仅对 1 个 token 做矩阵-向量乘 & 必须把全部模型权重和全部历史 KV Cache 从显存重新读一遍！ & 极低 ($\approx 1 \sim 2$) & \textbf{访存受限 (Memory-Bound)}，算力被闲置，卡在显存带宽搬运上 \\
\bottomrule
\end{tabularx}
\end{center}

\subsection{单 Token 解码耗时下界理论估算}

在 Batch Size = 1 时，生成 1 个 token 必须把整个模型权重读取一遍：

\[
T_{\text{decode\_min}} \approx \frac{\text{模型权重体积 (Bytes)}}{\text{显存有效读取带宽 (Bytes/s)}}
\]

\begin{noteBox}[核心导读与背景]
\textbf{例如}：在显存带宽为 448 GB/s 的显卡上运行 14 GB 的 7B FP16 模型：
每生成 1 个 token 至少需要 $14 \text{ GB} / 448 \text{ GB/s} \approx 31.25 \text{ ms}$（即吞吐上限约为 $32 \text{ tokens/s}$）。
如果通过 INT4 量化将权重压缩为 3.5 GB，理论下界降为 $3.5 / 448 \approx 7.8 \text{ ms}$（吞吐暴涨到 $128 \text{ tokens/s}$）！这就是量化在推理阶段能带来巨大速度提升的根本原因！

{}[注意] 上面的 448 GB/s 是\textbf{厂商标称值}，实际达不到（本机实测约 380 GB/s，见下面"一次实测结果"）。
用实测值算，同一个例子的下界是 $14/380 \approx 36.8 \text{ ms}$ 而不是 31.25 ms —— \textbf{标称值会乐观约 18\%}。
\texttt{memory\_calculator.py} 现在两种都打印，默认用实测值。
\end{noteBox}

\begin{noteBox}[核心导读与背景]
\textbf{单位口径}：本节算例用\textbf{十进制 GB}（7B $\times$ 2 字节 = $14 \times 10^9$ B = 14 GB），
因为要和 GB/s 相除。而 \texttt{memory\_calculator.py} 的权重/KV 打印用 \textbf{GiB}（除以 $1024^3$），
因为要和显卡容量比较。同一个 7B FP16 模型：14 GB = 13.04 GiB，\textbf{两个都对，差 7.4\%}。
混用是这一行最常见的错误，脚本里专门有一段注释说明。
\end{noteBox}

\vspace{0.8em}\hrule\vspace{0.8em}

\section{[工具] 显存与性能计算工具使用}

运行 \textbf{\texttt{memory\_calculator.py}}：

\begin{lstlisting}[language=bash]
python3 lab08_memory_and_roofline/memory_calculator.py
\end{lstlisting}

可评估主流模型（Llama-3-8B/70B、Qwen-2.5-7B/72B 等）在各类 GPU（如 RTX 5060 Ti、RTX 4090、A100）下的显存与并发上限。

\vspace{0.8em}\hrule\vspace{0.8em}

\section{[实验] 在你自己的显卡上实测 Roofline（需要 PyTorch）}

理论推导见 \textbf{第 8 篇：硬件视角}。运行：

\begin{lstlisting}[language=bash]
.venv/bin/python lab08_memory_and_roofline/measure_gpu_roofline.py
\end{lstlisting}

它会测出你这张卡的真实带宽 $\beta$、BF16 峰值 $\pi$、拐点 AI* = $\pi$/$\beta$，并验证两个关键推导：
\begin{enumerate}
  \item \textbf{Decode 线性层的算术强度 $\approx$ batch size}：batch 从 1 到 16，耗时几乎不变；
  \item \textbf{Decode 注意力的算术强度 $\approx$ GQA 分组大小}：与 batch 无关，所以 GQA/MLA 不仅省显存，还直接提速。
\end{enumerate}

本机（RTX 5060 Ti 16GB）的一次实测结果：

\begin{noteBox}[核心导读与背景]
{}[计时] 以下计时类数字是本机某一次运行的\textbf{量级示例}，前提是显卡空闲、已预热；GPU 降频或有其它负载时可能慢 40\% 以上。请看\textbf{趋势和比例}，以你自己运行的输出为准。 例如有评审在并发负载下测到过 217\textasciitilde{}234 GB/s 的带宽。
\end{noteBox}

\begin{center}
\small
\begin{tabularx}{\textwidth}{p{3cm} p{4.5cm} X}
\toprule
\textbf{指标} & \textbf{厂商标称} & \textbf{实测} \\
\midrule
显存带宽 & 448 GB/s & $\approx$ 380 GB/s（拷贝），$\approx$ 400 GB/s（GEMV 读权重） \\
BF16 Tensor 稠密算力 & 宣传口径更高 & \textbf{$\approx$ 48 TFLOP/s} \\
拐点 AI* & — & \textbf{$\approx$ 130 FLOP/Byte}：batch < 100 左右的 Decode 都是访存受限 \\
\bottomrule
\end{tabularx}
\end{center}

\begin{warningBox}[避坑指南与核心警示]
{}[注意] 实测中还会看到 batch = 32/64 的线性层比 batch = 128 还慢——这是 cuBLAS 对这些形状选中了低效 kernel，是真实存在的现象，不是噪声。它提醒你：\textbf{Roofline 是上界，实现质量决定你离上界有多远。}
\end{warningBox}



\section{实验核心源码精选与解析}

本实验配套有完整可运行的工业级验证代码，重点实现细节如下：


\subsection{源码实现：\texttt{memory\_calculator.py}}

\lstinputlisting[language=Python, caption={\texttt{memory\_calculator.py}}]{code/lab08_memory_and_roofline_memory_calculator.py}


\subsection{源码实现：\texttt{measure\_gpu\_roofline.py}}

\lstinputlisting[language=Python, caption={\texttt{measure\_gpu\_roofline.py}}]{code/lab08_memory_and_roofline_measure_gpu_roofline.py}


\chapter{FlashAttention 算法仿真与 CUDA Softmax 算子}

\section{Lab 09：FlashAttention 与 FlashDecoding}

\begin{mathBox}[核心数学定理推导]
对应理论：\textbf{第 9 篇：FlashAttention——Online Softmax 的严格推导与 IO 复杂度}
前置：第 8 篇（存储层级、算子融合）。
依赖：NumPy；实验 5 可选 PyTorch + GPU。
\end{mathBox}

\section{运行}

\begin{lstlisting}[language=bash]
.venv/bin/python lab09_flash_attention/flash_attention_numpy.py
\end{lstlisting}

\begin{center}
\small
\begin{tabularx}{\textwidth}{p{3.5cm} X}
\toprule
\textbf{实验} & \textbf{验证的结论} \\
\midrule
1 & Online Softmax 一遍扫描 = 三遍扫描；打印每一步的修正因子 $e^{m_{old}-m_{new}}$ \\
2 & 任意块大小的 FlashAttention（含因果掩码、整块跳过）与朴素注意力误差 \textasciitilde{}1e-15 \\
3 & 访存量模型（元素个数口径）：渐近节省倍数 = $2B_r/d = M/(2d^2)$，与论文的 $M/d^2$ 同量级，差的 2 是记账口径；小 N 时因尾块占比大还到不了这个常数；随 N 增长\textbf{真正}的收益是 $O(N^2) \to O(N)$ 的显存 \\
4 & FlashDecoding：KV 切 1/4/16/64 段分别计算再合并，结果一致；合并满足结合律 \\
5 & GPU 实测：朴素注意力 vs \texttt{F.scaled\_dot\_product\_attention} \\
\bottomrule
\end{tabularx}
\end{center}

\section{本机（RTX 5060 Ti）实验 5 结果}

\begin{noteBox}[核心导读与背景]
{}[计时] 以下计时类数字是本机某一次运行的\textbf{量级示例}，前提是显卡空闲、已预热；GPU 降频或有其它负载时可能慢 40\% 以上。请看\textbf{趋势和比例}，以你自己运行的输出为准。
\end{noteBox}

\begin{center}
\small
\begin{tabularx}{\textwidth}{l X X X X }
\toprule
\textbf{N} & \textbf{朴素} & \textbf{朴素额外显存} & \textbf{SDPA} & \textbf{SDPA 额外显存} \\
\midrule
1024 & 0.67 ms & 88 MB & 0.12 ms & 2 MB \\
4096 & 14.4 ms & 1.3 GB & 0.91 ms & 8 MB \\
8192 & 57.3 ms & 5.1 GB & 3.3 ms & 16 MB \\
16384 & \textbf{2309 ms} & \textbf{20 GB} & 12.6 ms & 33 MB \\
\bottomrule
\end{tabularx}
\end{center}

最后一行的朴素实现需要 20GB，超过了 16GB 显存却没有报 OOM，而是慢了 180 倍：Windows/WSL 的驱动把溢出部分放进了\textbf{系统内存}，通过 PCIe 访问。这是一个很好的提醒——"能跑"不等于"跑在显存里"。

\section{[附件] 附录：真机 CUDA 版（强烈建议做）}

上面的实验全部是 NumPy 仿真——它证明\textbf{数学}是对的，但访存量都是算出来的。\textbf{\textbf{\texttt{cuda\_softmax/}} 用真正的 CUDA C++ 把同一套 online softmax 写成 kernel}，在你的显卡上量出带宽：

\begin{center}
\small
\begin{tabularx}{\textwidth}{l X X X X }
\toprule
\textbf{实现} & \textbf{DRAM 搬运} & \textbf{耗时} & \textbf{占理论上界} & \textbf{相对} \\
\midrule
三段式（= 本实验 §1 的朴素版） & 4 次 & 0.694 ms & 51\% & 1.00x \\
单 kernel + online softmax & 2 次 & 0.359 ms & \textbf{97\textasciitilde{}98\%} & \textbf{1.9x} \\
同上 + 行缓存进 shared & 2 次 & 0.360 ms & \textbf{97\textasciitilde{}98\%} & \textbf{1.9x} \\
\bottomrule
\end{tabularx}
\end{center}

三个实现\textbf{都贴着 380 GB/s 的 Roofline}，差别只在搬了多少字节。里面还有一个反直觉结果：\textbf{行缓存进 shared 的版本在 N=8192 时反而崩到 74\%}——因为 34 KB shared/block 把每个 SM 的常驻 block 数从 6 压到 2。\textbf{这正是 FlashAttention 要做分块而不是缓存整行的原因。}

\begin{lstlisting}[language=bash]
.venv/bin/python lab09_flash_attention/cuda_softmax/run.py
\end{lstlisting}

不需要装 CUDA toolkit（用的是 PyTorch 自带的 NVRTC）。详见 \textbf{\texttt{cuda\_softmax/README.md}}。

\section{动手练习}

\begin{enumerate}
  \item 在 \texttt{flash\_attention} 里加上对 GQA 的支持：多个 Q 头共享同一组 K/V 块时，同一个 K/V 块读一次就能服务多个头（对照第 8 篇 §4.3）。
  \item 实现 FlashAttention-1 的循环顺序（外层 K/V、内层 Q），需要把每个 Q 块的 $m, \ell$, acc 写回"显存"再读出来；统计访存量，理解 FlashAttention-2 为什么交换循环。
  \item 把实验 4 的 \texttt{partial\_attention} 改成读取\textbf{不连续}的 KV 块（给一个 block table，见 Lab 12b），这就是 PagedAttention kernel 的核心逻辑。
  \item （进阶）在 \textbf{\texttt{cuda\_softmax/}} 里把行缓存版改成\textbf{真正的分块}（只缓存 \texttt{BLOCK} 个元素、循环推进），验证占用率恢复后带宽回到 97\textasciitilde{}98\%——做完这题你就理解了 tiling 是怎么来的。
  \item （进阶）跟着 \href{https://triton-lang.org/main/getting-started/tutorials/06-fused-attention.html}{Triton Fused Attention 教程} 写一个真正在 GPU 上运行的 FlashAttention 前向，与 SDPA 对比速度。
\end{enumerate}



\section{实验核心源码精选与解析}

本实验配套有完整可运行的工业级验证代码，重点实现细节如下：


\subsection{源码实现：\texttt{flash\_attention\_numpy.py}}

\lstinputlisting[language=Python, caption={\texttt{flash\_attention\_numpy.py}}]{code/lab09_flash_attention_flash_attention_numpy.py}


\chapter{投机解码仿真实现与无偏性实证校验}

\section{Lab 10：投机采样 (Speculative Decoding) 原理与算法仿真}

\section{为什么需要投机采样？}

在 Lab 08 中我们已经知道：
\begin{itemize}
  \item \textbf{自回归 Decode 阶段是访存带宽受限（Memory-Bound）的}。
  \item 大模型每生成 1 个 token，就必须把整个模型的权重（数十 GB）完整搬运一次，因此单步延迟很难降低。
  \item 但如果同时输入 $K$ 个 token（如 Prefill 阶段），GPU 的算力利用率大幅提高，耗时\textbf{几乎与输入 1 个 token 相同}！
\end{itemize}

\textbf{投机采样（Speculative Decoding）的洞察}：
\begin{noteBox}[核心导读与背景]
"与其每次花很长时间让昂贵的大模型只猜 1 个词，不如让一个廉价的小模型快速一口气猜 $K$ 个词，再让大模型做一次并行前向传播进行批量验证！"
\end{noteBox}

\vspace{0.8em}\hrule\vspace{0.8em}

\section{核心算法流程}

假设目标大模型为 $M_{\text{target}}$，轻量草稿模型为 $M_{\text{draft}}$（参数量通常仅为大模型的 $1/10 \sim 1/5$）：

\begin{figure}[htbp]
\centering
\begin{tikzpicture}[
  actor/.style={draw=primary, fill=primary!10, rounded corners=2mm, font=\bfseries\small\sffamily, align=center, inner sep=6pt, minimum width=4cm},
  msg/.style={-{Stealth[scale=0.9]}, draw=secondary, line width=1.1pt, font=\small\sffamily},
  note/.style={draw=gray!50, fill=gray!5, rounded corners=1.5mm, font=\footnotesize\sffamily, align=center, inner sep=5pt}
]
\node[actor] (draft) {草稿小模型 ($M_{\mathrm{draft}}$)\\小而快 (Draft Model)};
\node[actor, right=5cm of draft] (target) {目标大模型 ($M_{\mathrm{target}}$)\\大而准 (Target Model)};

\coordinate[below=6.2cm of draft] (draft_end);
\coordinate[below=6.2cm of target] (target_end);

\draw[line width=1pt, gray!60, dashed] (draft) -- (draft_end);
\draw[line width=1pt, gray!60, dashed] (target) -- (target_end);

\node[note, below=0.8cm of draft] (n1) {自回归快速生成 $K$ 个候选 Token\\$[x_1, x_2, x_3, x_4]$};
\draw[msg] ($(draft.south)+(0,-1.8)$) -- node[above, font=\small\sffamily\color{secondary}] {1. 提交草稿预测序列 $[x_1, \dots, x_K]$} ($(target.south)+(0,-1.8)$);
\node[note, below=2.2cm of target] (n2) {单次并行前向传播 (Parallel Forward)\\同时计算全部 $K$ 个位置真实概率分布 $p(x)$};
\node[note, below=3.6cm of target] (n3) {逐位执行投机判定：\\接受率 $r_i = \min\left(1, \frac{p(x_i)}{q(x_i)}\right)$};
\draw[msg] ($(target.south)+(0,-5.2)$) -- node[above, font=\small\sffamily\color{secondary}] {2. 返回接收的前缀 + 补采样调整词元} ($(draft.south)+(0,-5.2)$);

\end{tikzpicture}
\caption{投机采样（Speculative Decoding）草稿与目标模型协同序列交互}
\end{figure}

\subsection{严格的数学无偏性保证 (Unbiased Sampling)}

投机采样最迷人的特点是：\textbf{完全无损（Lossless）}。
通过巧妙的拒绝采样概率公式：

\[
\alpha = \min\left(1, \frac{P_{\text{target}}(x)}{P_{\text{draft}}(x)}\right)
\]

\begin{itemize}
  \item 若以概率 $\alpha$ 接受草稿 token $x$；
  \item 若拒绝，则根据校正后的残差分布采样：
\end{itemize}

\[
P_{\text{residual}}(x) = \frac{\max(0, P_{\text{target}}(x) - P_{\text{draft}}(x))}{\sum \max(0, P_{\text{target}}(x') - P_{\text{draft}}(x'))}
\]

数学上严格证明：\textbf{无论草稿模型多弱，最终生成的文本统计分布与单独用大模型生成 100\% 完全一致！}
（草稿模型越弱，只会让接受率变低、加速比变小，而不会改变输出分布。）

\begin{mathBox}[核心数学定理推导]
{}[几何] \textbf{证明在哪？} 完整的逐行证明见 \textbf{第 10 篇：投机解码}；
用蒙特卡洛实验"亲眼验证"这个定理，请运行 \textbf{\texttt{verify\_unbiasedness.py}}。

{}[注意] 注意：本目录的 \texttt{toy\_speculative\_decoding.py} 里，接收/拒绝只用了 \textbf{被采样 token 的一个概率值}，所以是正确的；
但 \texttt{p\_target / p\_draft} 必须用\textbf{同一个前缀}下的两个分布——这一点在真实实现里最容易写错。
\end{mathBox}

\vspace{0.8em}\hrule\vspace{0.8em}

\section{[执行] 运行仿真实验}

运行 \textbf{\texttt{toy\_speculative\_decoding.py}}：

\begin{lstlisting}[language=bash]
python3 lab10_speculative_decoding/toy_speculative_decoding.py
\end{lstlisting}

本脚本模拟了草稿模型与目标模型的联合采样过程，展示每轮投机接受长度、大模型调用次数节省比例及端到端加速比。

\begin{lstlisting}[language=bash]
python3 lab10_speculative_decoding/verify_unbiasedness.py
\end{lstlisting}

对同一个前缀重复做 20 万次"投机采样的一步"，统计输出 token 的经验分布，与目标模型分布逐项比较（总变差距离应趋近 0），并给出理论接受率 $\sum_x \min(p, q)$ 与期望加速比公式的实测验证。还会演示两种\textbf{常见错误写法}（拒绝后直接从 $p$ 或 $q$ 重采样）如何让分布产生 11\%\textasciitilde{}17\% 的总变差偏差。


\section{实验核心源码精选与解析}

本实验配套有完整可运行的工业级验证代码，重点实现细节如下：


\subsection{源码实现：\texttt{toy\_speculative\_decoding.py}}

\lstinputlisting[language=Python, caption={\texttt{toy\_speculative\_decoding.py}}]{code/lab10_speculative_decoding_toy_speculative_decoding.py}


\subsection{源码实现：\texttt{verify\_unbiasedness.py}}

\lstinputlisting[language=Python, caption={\texttt{verify\_unbiasedness.py}}]{code/lab10_speculative_decoding_verify_unbiasedness.py}


\chapter{工业量化机制（RTN / SmoothQuant / AWQ / GPTQ）实战}

\section{Lab 11：模型量化技术 (INT8 / INT4 / FP8, AWQ, GPTQ)}

\section{[目标] 什么是量化？为什么对推理如此重要？}

在深度学习推理中，\textbf{量化（Quantization）是将高精度浮点数（如 FP32 / FP16）映射到低位宽定点数或微缩浮点数（如 INT8 / INT4 / FP8）的技术}。

对于大模型推理，量化带来双重巨大收益：
\begin{enumerate}
  \item \textbf{显存容量降低}：70B 模型从 FP16 的 140GB 显存，经 INT4 压缩后仅需约 35GB，可以在单张或双张消费级显卡上加载。
  \item \textbf{生成速度翻倍}：由于 Decode 阶段瓶颈在于显存带宽，将权重体积压缩至 1/4 后，显存读取时间直接缩短至 1/4，理论生成速度可提升至接近 4 倍！
\end{enumerate}

\vspace{0.8em}\hrule\vspace{0.8em}

\section{[几何] 核心量化数学公式}

\subsection{对称量化 (Symmetric Quantization)}

将浮点零点映射为定点零点（即零偏置 $Z = 0$）：
\begin{itemize}
  \item 缩放因子（Scale）：
\end{itemize}

\[
S = \frac{\max(|X|)}{2^{b-1} - 1}
\]

\begin{itemize}
  \item 量化过程：
\end{itemize}

\[
q = \text{clip}\left(\text{round}\left(\frac{X}{S}\right), -2^{b-1}+1, 2^{b-1}-1\right)
\]

\begin{itemize}
  \item 反量化（Dequantization）：
\end{itemize}

\[
\hat{X} = q \times S
\]

\subsection{非对称量化 (Asymmetric / Affine Quantization)}

当数据分布具有明显单边偏置时，引入零点 Zero-point ($Z$)：
\begin{itemize}
  \item 缩放因子与零点：
\end{itemize}

\[
S = \frac{X_{\max} - X_{\min}}{2^b - 1}
\]

\[
Z = \text{round}\left(-\frac{X_{\min}}{S}\right)
\]

\begin{itemize}
  \item 量化与反量化：
\end{itemize}

\[
q = \text{clip}\left(\text{round}\left(\frac{X}{S}\right) + Z, 0, 2^b - 1\right)
\]

\[
\hat{X} = (q - Z) \times S
\]

\vspace{0.8em}\hrule\vspace{0.8em}

\section{[模块] 量化粒度 (Granularity)}

\begin{center}
\small
\begin{tabularx}{\textwidth}{p{2.5cm} p{3.5cm} X X}
\toprule
\textbf{粒度模式} & \textbf{描述} & \textbf{硬件友好度} & \textbf{精度保持度} \\
\midrule
\textbf{Per-Tensor} & 整个张量共享一个 Scale & 最高（计算开销最小） & 最差（易受离群异常值干扰） \\
\textbf{Per-Channel} & 矩阵的每个输出通道拥有独立的 Scale & 高 & 较好 \\
\textbf{Per-Group (分组量化)} & 连续每 32 / 64 / 128 个权重组成一组独立量化（AWQ/GPTQ 标配；本仓库脚本用 \texttt{group\_size=32} 演示） & 中等 & \textbf{最好}（极大减小离群值污染范围） \\
\bottomrule
\end{tabularx}
\end{center}

\vspace{0.8em}\hrule\vspace{0.8em}

\section{[思考] 工业级算法速览}

\begin{enumerate}
  \item \textbf{W4A16 (Weight-Only Quantization)}:
\end{enumerate}

\begin{itemize}
  \item 权重存为 INT4，计算时在寄存器/SRAM 实时反量化为 FP16 与激活值做运算。
  \item 代表：\textbf{AWQ (Activation-aware Weight Quantization)}（保护 1\% 关键特征通道）和 \textbf{GPTQ}（二阶 Hessian 误差补偿）。原理与实现见下方"进阶"。
\end{itemize}

\begin{enumerate}
  \item \textbf{W8A8 (Weight \& Activation Quantization)}:
\end{enumerate}

\begin{itemize}
  \item 权重与激活值均量化为 INT8 或 FP8，直接使用硬件 INT8/FP8 Tensor Core 进行矩阵乘法，Prefill 和 Decode 均获加速。
  \item 代表：\textbf{SmoothQuant}（通过数学等价变换将激活值的离群峰值转移给权重）。
\end{itemize}

\begin{enumerate}
  \item \textbf{FP8 量化（现代 GPU 标配）}:
\end{enumerate}

\begin{itemize}
  \item Hopper/Ada/Blackwell 架构原生支持 E4M3 和 E5M2 两种 FP8 格式，几乎无损兼顾训练与推理。
\end{itemize}

\vspace{0.8em}\hrule\vspace{0.8em}

\section{[执行] 运行量化实验}

运行 \textbf{\texttt{quant\_basics.py}}：

\begin{lstlisting}[language=bash]
python3 lab11_quantization/quant_basics.py
\end{lstlisting}

对比对称量化、非对称量化以及 Group-wise INT4 分组量化的重构误差（MSE）与信噪比（SNR）。

\vspace{0.8em}\hrule\vspace{0.8em}

\section{[实验] 进阶：从零实现 SmoothQuant / AWQ / GPTQ}

原理推导见 \textbf{第 11 篇：量化}。

\begin{lstlisting}[language=bash]
.venv/bin/python lab11_quantization/advanced_quant_numpy.py
\end{lstlisting}

\begin{center}
\small
\begin{tabularx}{\textwidth}{p{3cm} p{4.5cm} X}
\toprule
\textbf{实验} & \textbf{本机结果} & \textbf{说明} \\
\midrule
每多 1 位 SNR +6 dB & INT7$\rightarrow$INT8 +6.1 dB & 误差方差 $S^2/12$ \\
激活离群值 & 6 个离群通道让普通通道 INT8 SNR 只剩 10 dB & LLM.int8() 的发现 \\
SmoothQuant & W8A8 输出 SNR 20.2 $\rightarrow$ 32.4 dB（$\alpha$=0.5） & 恒等变换，零推理开销 \\
AWQ & W4 分组量化输出 SNR +2.9 dB，关键通道误差降为 1/2.6 & 放大关键通道再量化 \\
GPTQ & 输出 SNR 18.2 $\rightarrow$ \textbf{29.6 dB}；但权重本身的误差反而更大 & 优化 $\|WX - \hat WX\|$ 而非 $\|W - \hat W\|$；输入不相关时毫无收益 \\
\bottomrule
\end{tabularx}
\end{center}

\textbf{动手练习}
\begin{enumerate}
  \item 把实验 5 的 GPTQ 改成\textbf{分组}量化（每 128 列一组各自一个 scale），与 AWQ 在同一份数据上比较。
  \item 在 Lab 07b 的 Qwen2.5-0.5B 上，对所有线性层做 W4 分组 RTN 量化（反量化回 BF16 再计算），用一段文本测 PPL 变化；再换成你实现的 GPTQ，看能挽回多少。
  \item 用第 8 篇的带宽公式估算：Qwen2.5-7B 在 5060 Ti 上从 BF16 换成 W4A16，Decode 理论上限从多少 token/s 变成多少？为什么 Prefill 不会同样变快？
\end{enumerate}



\section{实验核心源码精选与解析}

本实验配套有完整可运行的工业级验证代码，重点实现细节如下：


\subsection{源码实现：\texttt{quant\_basics.py}}

\lstinputlisting[language=Python, caption={\texttt{quant\_basics.py}}]{code/lab11_quantization_quant_basics.py}


\subsection{源码实现：\texttt{advanced\_quant\_numpy.py}}

\lstinputlisting[language=Python, caption={\texttt{advanced\_quant\_numpy.py}}]{code/lab11_quantization_advanced_quant_numpy.py}


\chapter{Continuous Batching（持续批处理）动态调度仿真}

\section{Lab 12a：显存调度与连续批处理 (Continuous Batching)}

\section{[施工] 为什么静态批处理（Static Batching）在 LLM 中失效？}

在传统的计算机视觉或文本分类任务中，模型输入输出维度是固定的，\textbf{静态批处理（Static Batching）}非常有效。

但在大语言模型生成场景中，每个请求的：
\begin{enumerate}
  \item \textbf{输入 Prompt 长度极不一致}（有的几十字，有的几千字）。
  \item \textbf{生成长度（Output Tokens）难以事先预知}（有的遇到 \texttt{$<$eos$>$} 10 步就结束，有的需要生成 1000 步）。
\end{enumerate}

\subsection{静态 Batching 的两大致命缺陷：}

\begin{lstlisting}[numbers=none]
Request 1: [Prompt: 10] [=== Gen: 20 ===] -> 结束! (闲置等待...) [PADDING PADDING PADDING]
Request 2: [Prompt: 30] [=================== Gen: 100 ===================] -> 最长请求!
Request 3: [Prompt: 5 ] [= Gen: 10 =] -> 结束! (闲置等待...) [PADDING PADDING PADDING PADDING]
\end{lstlisting}

\begin{enumerate}
  \item \textbf{显存与算力气泡（Bubble）}：短请求早早结束，但由于必须作为一个 Batch 统一返回，短请求必须不断 padding 空 token 陪跑，浪费大量算力与带宽。
  \item \textbf{首字延迟排队灾难}：新到达的请求必须等待当前 Batch 中\textbf{最慢的那一个请求完全结束}，才能开始下一批，导致极高的排队排班延迟。
\end{enumerate}

\vspace{0.8em}\hrule\vspace{0.8em}

\section{[加速] 连续批处理 (Continuous Batching / In-flight Batching)}

由 OSDI'22 论文 \textit{Orca} 提出，并在 \textbf{vLLM} 和 \textbf{TensorRT-LLM} 中发扬光大：

\begin{noteBox}[核心导读与背景]
\textbf{核心思想}：调度粒度从「请求级（Request-level）」细化为「迭代步进级（Iteration-level）」。
\end{noteBox}

\begin{lstlisting}[numbers=none]
Iteration t:   [Req 1: decode step 15] [Req 2: decode step 40] [Req 3: decode step 8]
Iteration t+1: [Req 1: 触发 EOS! 退出]   [Req 2: decode step 41] [Req 3: decode step 9]
                ⬇ (槽位立即释放，新请求无缝插入！)
Iteration t+2: [Req 4: 快速 Prefill!]   [Req 2: decode step 42] [Req 3: decode step 10]
\end{lstlisting}

\begin{itemize}
  \item \textbf{立即退场}：任何请求一旦生成结束标记（EOS）或达到最大长度，当步立刻释放槽位与 KV Cache。
  \item \textbf{即时插入}：等待队列中的新请求无需等待其他长请求完成，在下一步前向迭代中直接加入当前 Batch 进行 Prefill。
  \item \textbf{吞吐飞跃}：消除了绝大部分 Padding 气泡，吞吐量显著提升（Orca、vLLM 等论文中常见数倍提升，具体取决于负载）。本仿真会同时打印"Padding 无效计算占比"和"槽位占用率"：前者对连续批处理按构造恒为 0，后者才是两种方案可比的指标。
\end{itemize}

\vspace{0.8em}\hrule\vspace{0.8em}

\section{[文档] 内存革命：PagedAttention 原理}

即使有了 Continuous Batching，连续的内存分配依然会造成严重浪费：
\begin{itemize}
  \item 在传统显存管理中，框架必须为请求预先分配一段\textbf{连续的固定最大长度显存}（如预留 4096 tokens）。
  \item 但大多数请求只生成了 100 个 token。vLLM 论文测到 60\%\textasciitilde{}80\% 的 KV 显存被浪费——这\textbf{不只是内部碎片}，还包括按 max\_len 预留的闲置和外部碎片，三者要分开看（第 12 篇 §3）。
\end{itemize}

\textbf{vLLM 的解决方案}：
借鉴操作系统的\textbf{虚拟内存分页机制（Virtual Memory Paging）}：
\begin{itemize}
  \item 将 KV Cache 切分成固定大小的物理块（Physical Blocks，例如每块 16 个 tokens）。
  \item 物理显存块\textbf{无需连续}，通过一个逻辑块表（Block Table）映射。
  \item 只有当生成了新的 16 个 token 时，才向物理池申请分配 1 个新块。
  \item \textbf{成效}：分页把「最后一个块没写满」以外的浪费基本消掉；单序列的上界是 $15/S$，序列长度 $S \gtrsim 400$ token 时才低于 4\%（短对话达不到，见第 12 篇 §3）；实际收益是同一张卡能承载成倍增加的并发请求数。
\end{itemize}

\vspace{0.8em}\hrule\vspace{0.8em}

\section{[执行] 运行仿真实验}

运行 \textbf{\texttt{continuous\_batching\_sim.py}}：

\begin{lstlisting}[language=bash]
python3 lab12a_continuous_batching/continuous_batching_sim.py
\end{lstlisting}

对比静态批处理与连续批处理在处理相同真实负载时的总吞吐量（Throughput）、平均等待延迟（Latency）与算力浪费（Bubble Ratio）。

\vspace{0.8em}\hrule\vspace{0.8em}

\section{[链接] 继续深入}

\begin{itemize}
  \item \textbf{为什么"每一步耗时视为常数"这个假设成立？} 本仿真假设无论批里有几个请求，一步的耗时都一样。这来自第 8 篇的推导：Decode 线性层的算术强度 $\approx$ batch size，batch 远小于拐点时，一步的耗时由"读一遍权重"决定，与 batch 几乎无关。见 \textbf{第 8 篇 §4} 与 \texttt{lab08\_memory\_and\_roofline/measure\_gpu\_roofline.py} 实验 4。
  \item \textbf{本仿真没有模拟的}：Prefill 的开销、显存上限、长 Prefill 对 Decode 的干扰。这些在 \textbf{第 12 篇} 与 \textbf{Lab 12b}（分页 KV Cache、前缀缓存、Chunked Prefill）中补齐。
\end{itemize}



\section{实验核心源码精选与解析}

本实验配套有完整可运行的工业级验证代码，重点实现细节如下：


\subsection{源码实现：\texttt{continuous\_batching\_sim.py}}

\lstinputlisting[language=Python, caption={\texttt{continuous\_batching\_sim.py}}]{code/lab12a_continuous_batching_continuous_batching_sim.py}


\chapter{PagedAttention（分页显存分配机制）仿真}

\section{Lab 12b：PagedAttention、写时复制、前缀缓存与 Chunked Prefill}

\begin{noteBox}[核心导读与背景]
对应理论：\textbf{第 12 篇：推理服务系统}
前置：第 6 篇、Lab 12a、第 8 篇、第 9 篇（Online Softmax）。
依赖：NumPy。
\end{noteBox}

\begin{noteBox}[核心导读与背景]
{}[计时] 以下计时类数字是本机某一次运行的\textbf{量级示例}（显卡空闲、已预热）。它是\textbf{第 8 篇耗时模型的推算}，不是真实 kernel 的实测，请看趋势和比例。
\end{noteBox}

\section{运行}

\begin{lstlisting}[language=bash]
.venv/bin/python lab12b_paged_attention/paged_kv_cache_sim.py
\end{lstlisting}

\begin{center}
\small
\begin{tabularx}{\textwidth}{p{3.5cm} X}
\toprule
\textbf{实验} & \textbf{你会看到} \\
\midrule
1. 分页注意力 & 三个请求交替写入，物理块号杂乱（如 \texttt{[49, 26, 32, 50]}），逐块 Online Softmax 结果与连续存储误差 \textasciitilde{}1e-16 \\
2. 显存利用率 & 同样 64K token 的显存：预留 \texttt{max\_len} 只能服务 32 个请求（浪费 78\%），分页可服务 144 个（浪费 1.7\%） \\
3. 写时复制 & 并行采样 n=4：共享提示词块只存一份，只有"共享且未写满"的最后一块被复制，而且只复制 n−1 次（最后一个子序列可原地写入），省 50\% \\
4. 前缀缓存 & 1000 token 系统提示词：84\% 的 token 免去 Prefill；改掉第 1 个 token 后缓存全部失效 \\
5. Chunked Prefill & 6000 token 长提示词插队：不切块时所有 Decode 用户卡顿 390ms；切成 256 一块后最长卡顿 22ms（代价是 TTFT 从 390ms 涨到 506ms） \\
\bottomrule
\end{tabularx}
\end{center}

\section{代码结构}

\begin{itemize}
  \item \texttt{BlockAllocator}：空闲链表 + 引用计数（$\approx$ 操作系统的物理页框管理）
  \item \texttt{PagedKVCache}：物理 KV 池 \texttt{[num\_blocks, 16, d]} + 每个序列一张块表；\texttt{append} 里实现按需分配与写时复制；\texttt{attention} 按块表逐块读取并用 Online Softmax 累加
\end{itemize}

这就是 vLLM 的 \texttt{BlockManager} 与 PagedAttention kernel 的核心逻辑，去掉了 GPU 与多层细节。

\section{动手练习}

\begin{enumerate}
  \item 给 \texttt{BlockAllocator} 加上 LRU 淘汰：请求结束后块不立即释放而是进入"可回收缓存"，把实验 4 的前缀缓存和实验 2 的分配器真正接在一起。
  \item 实现抢占：显存不足时选择最晚到达的请求，分别实现"换出到 CPU"与"丢弃后重算"，用第 8 篇的耗时模型比较两者代价（提示：PCIe 带宽约 25 GB/s，Prefill 的算力受限耗时见实验 5）。
  \item 把块大小改成 1、16、128，观察实验 2 的浪费率与"块表长度"之间的权衡（块太小，块表和 kernel 里的间接寻址开销变大）。
  \item （阅读）对照 vLLM 源码中 KV Cache 管理与调度器相关的模块（在官方文档 Design 部分可以找到架构说明），找到与本文件中每个函数对应的地方。
\end{enumerate}



\section{实验核心源码精选与解析}

本实验配套有完整可运行的工业级验证代码，重点实现细节如下：


\subsection{源码实现：\texttt{paged\_kv\_cache\_sim.py}}

\lstinputlisting[language=Python, caption={\texttt{paged\_kv\_cache\_sim.py}}]{code/lab12b_paged_attention_paged_kv_cache_sim.py}


\chapter{多卡张量并行（Tensor Parallelism）切分与 All-Reduce 仿真}

\section{Lab 13：多卡并行推理仿真}

\begin{noteBox}[核心导读与背景]
对应理论：\textbf{第 13 篇：多卡并行推理——张量并行、流水线并行与专家并行}
前置：第 3 篇（矩阵乘形状）、第 8 篇（带宽与耗时模型）。
依赖：NumPy。不需要真的有多张显卡——用列表里的多个数组模拟多张卡，重点是看清"切法"的数学。
\end{noteBox}

\begin{noteBox}[核心导读与背景]
{}[计时] 以下计时类数字是本机某一次运行的\textbf{量级示例}（显卡空闲、已预热）。它是\textbf{第 8 篇耗时模型的推算}，不是真实 kernel 的实测，请看趋势和比例。
\end{noteBox}

\section{运行}

\begin{lstlisting}[language=bash]
.venv/bin/python lab13_parallelism/tensor_parallel_sim.py
\end{lstlisting}

\begin{center}
\small
\begin{tabularx}{\textwidth}{p{3.5cm} X}
\toprule
\textbf{实验} & \textbf{结果} \\
\midrule
1. TP SwiGLU（先列后行） & TP=2/4/8 与单卡误差 \textasciitilde{}1e-15，每卡权重 1/n \\
2. 反例（先行切） & 在部分和上做 SiLU，相对误差 \textbf{52\%}；要算对必须多一次通信 \\
3. TP 注意力（按头切） & 与单卡一致，每卡只存 1/n 的 KV Cache \\
4. Ring All-Reduce & 每卡发送 $2(n-1)/n$ 倍数据：2 卡 1.0、4 卡 1.5、8 卡 1.75 \\
5. 70B Decode 延迟模型 & NVLink TP=8 约 6.9x；PCIe TP=8 约 4.9x（通信占近 40\%）；PP 不降延迟 \\
\bottomrule
\end{tabularx}
\end{center}

\section{动手练习}

\begin{enumerate}
  \item 在实验 3 中实现 GQA（8 个 Q 头、2 个 KV 头）+ TP=4：KV 头少于卡数时，每个 KV 头要复制到哪几张卡上？KV Cache 总占用变成多少？
  \item 给实验 5 加上 batch 维度：batch=64 时每次 All-Reduce 的数据量变成多少？通信是否从"延迟主导"变成"带宽主导"？
  \item 模拟专家并行：8 个专家分到 4 张卡，随机路由 1024 个 token（top-2），统计每张卡收到的 token 数，计算最忙的卡比平均多多少（负载不均衡）。再让路由分布偏向少数专家，看看会怎样。
  \item 用实验 4 的环形结构实现 Ring Attention：每张卡持有一段 KV，环形传递 KV 块，用第 9 篇的合并公式累加，验证与完整注意力一致。
\end{enumerate}



\section{实验核心源码精选与解析}

本实验配套有完整可运行的工业级验证代码，重点实现细节如下：


\subsection{源码实现：\texttt{tensor\_parallel\_sim.py}}

\lstinputlisting[language=Python, caption={\texttt{tensor\_parallel\_sim.py}}]{code/lab13_parallelism_tensor_parallel_sim.py}


\chapter{生产级推理引擎（vLLM / SGLang / TRT-LLM）压测实践}

\section{毕业设计：工业级推理引擎实践（vLLM / llama.cpp）}

\begin{mathBox}[核心数学定理推导]
这是整条路线的\textbf{毕业设计}。前面的 Lab 03\textasciitilde{}13 里你已经亲手实现了推理引擎的每一个核心部件；现在去使用真正的引擎，并且\textbf{每一个观测到的数字，都要能用前面的公式解释}。

前置：第 8、10、11、12 篇；Lab 08（实测带宽）、Lab 07b（手写 Qwen）、Lab 12b（PagedAttention）。
\end{mathBox}

\section{原则：先预测，再测量，再解释}

每做一个实验，先用 \texttt{lab08\_memory\_and\_roofline/memory\_calculator.py} 和第 8 篇的公式\textbf{写下你的预测}，再运行，最后解释差距。这比跑出一个漂亮的数字重要得多。

\vspace{0.8em}\hrule\vspace{0.8em}

\section{[测试] 任务 1：用 vLLM 部署，并测出吞吐-延迟曲线}

\textbf{安装}：建议单独建一个虚拟环境（vLLM 会固定自己需要的 PyTorch 版本，可能和本仓库的 \texttt{.venv} 冲突）。RTX 50 系显卡需要支持 CUDA 12.8 的较新版本，按 \href{https://docs.vllm.ai/}{vLLM 官方安装文档} 操作。WSL2 下可以正常使用。

\begin{lstlisting}[language=bash]
vllm serve Qwen/Qwen2.5-1.5B-Instruct --max-model-len 4096 --gpu-memory-utilization 0.85
\end{lstlisting}

\textbf{先预测}：
\begin{itemize}
  \item Qwen2.5-1.5B：28 层、2 个 KV 头、head\_dim 128 $\rightarrow$ 每 token KV = $2 \times 28 \times 2 \times 128 \times 2 = 28{,}672$ 字节 $\approx$ 28 KB；
  \item 权重 BF16 约 3.1 GB，按你实测的带宽（本机空闲热态约 380 GB/s，负载下会更低），单流 Decode 上限 $\approx$ 120 token/s；
  \item 16GB $\times$ 0.85 扣掉权重和激活后，大约还剩多少 GB 给 KV？能放多少 token？
\end{itemize}

\textbf{再测量}：
\begin{enumerate}
  \item 看 vLLM 启动日志里的 KV Cache 容量（不同版本措辞不同，通常会打印 KV Cache 的 token 数或 GPU block 数与最大并发）。和你的预测一致吗？block 数 $\times$ 16 = token 数？
  \item 压测：
\end{enumerate}

\begin{lstlisting}[language=bash]
   python3 capstone_engines_practice/bench_openai_server.py --concurrency 1 4 16 64 --max-tokens 256
   
\end{lstlisting}

\begin{enumerate}
  \item 画出（或列出）"并发 $\rightarrow$ 总吞吐、TPOT P50"的关系。
\end{enumerate}

\begin{warningBox}[避坑指南与核心警示]
{}[注意] \textbf{压测本身有三个坑，不知道就会得出错误的结论}（脚本已在文件头 docstring 里写明，这里再强调一次）：
- \textbf{前缀缓存会污染跨档位的对比}：脚本各档位复用同一批问题，档位之间不清缓存，而 vLLM 默认开启前缀缓存 $\rightarrow$ 靠后的档位 TTFT 偏低。要得到干净的曲线，请在档位之间重启服务。
- \textbf{输出长度必须一致}：默认加了 \texttt{ignore\_eos: true}（vLLM/SGLang 扩展），让每个请求都恰好输出 \texttt{--max-tokens} 个 token。否则"总吞吐 tok/s"测的是模型有多话唠，不是引擎有多快。标准 OpenAI 服务会忽略这个字段，那就只能靠 \texttt{max\_tokens} 近似。
- \textbf{必须预热}：脚本每档先发 \texttt{--warmup}（默认 2）个不计入统计的请求。GPU 从空闲频率爬到稳态要时间——lab08 里同一份代码在不同时钟状态下读到过 217\textasciitilde{}380 GB/s，差近一倍。第一档不预热会显著偏慢。
\end{warningBox}

\begin{noteBox}[核心导读与背景]
{}[重点] 另外，脚本会\textbf{自检} \texttt{usage.completion\_tokens} 和实际收到的内容分片数是否一致：不一致时说明服务端一个 SSE chunk 带了多个 token（llama.cpp 常见），此时 TPOT 按分片间隔计算并打印告警——否则 TPOT 会被系统性高估。样本不足（例如 \texttt{--max-tokens 1}）时它会打印"样本不足"而不是崩溃。
\end{noteBox}

\textbf{要能解释}：
\begin{itemize}
  \item 并发 1 的单流速度为什么低于 120 token/s，但比 Lab 07b 的手写实现（约 20ms/token）快得多？（提示：CUDA Graph、融合 kernel，第 12 篇 §7）
  \item 并发从 1 到 16，总吞吐为什么几乎线性增长、TPOT 却几乎不变？到多少并发开始 TPOT 明显变差？（第 8 篇 §4：拐点约 130）
\end{itemize}

\vspace{0.8em}\hrule\vspace{0.8em}

\section{[测试] 任务 2：前缀缓存}

\begin{lstlisting}[language=bash]
python3 capstone_engines_practice/bench_openai_server.py --shared-prefix-tokens 3000 --concurrency 4 --max-tokens 64
\end{lstlisting}

分别在开启和关闭前缀缓存时运行（新版 vLLM 默认开启；关闭参数通常是 \texttt{--no-enable-prefix-caching}，以你的版本 \texttt{vllm serve --help} 为准）。

\textbf{预测}：3000 token 的 Prefill 按第 8 篇 $2N \cdot S / (\pi \cdot \text{MFU})$ 估算要多少毫秒？命中缓存后 TTFT 应该降到多少？

\vspace{0.8em}\hrule\vspace{0.8em}

\section{[测试] 任务 3：长 Prefill 对 Decode 的干扰（Chunked Prefill）}

同时跑两个压测：一个并发 16、短问题、长输出（模拟正在聊天的用户）；另一个偶尔发送 3000+ token 的长提示词。观察第一个的 \textbf{TPOT P99} 是否出现尖刺。
调整 \texttt{--max-num-batched-tokens}（每步 token 预算）重复实验，用 Lab 12b 实验 5 的模型解释你看到的现象。

\vspace{0.8em}\hrule\vspace{0.8em}

\section{[测试] 任务 4：llama.cpp + GGUF 量化}

\begin{lstlisting}[language=bash]
git clone https://github.com/ggml-org/llama.cpp && cd llama.cpp
cmake -B build -DGGML_CUDA=ON && cmake --build build --config Release -j
# 从 Hugging Face 上 Qwen 官方的 GGUF 仓库下载 Qwen2.5-7B-Instruct 的 Q4_K_M 版本
./build/bin/llama-bench -m <你的 gguf 文件> -ngl 99
\end{lstlisting}

\texttt{llama-bench} 会报告 \texttt{pp512}（Prefill 512 token 的吞吐）和 \texttt{tg128}（Decode 的吞吐）。

\textbf{预测}：
\begin{itemize}
  \item Q4\_K\_M 平均约 4.8 bit/参数，7.6B 参数约 4.6 GB $\rightarrow$ \texttt{tg} 上限 $\approx$ 380 / 4.6 $\approx$ 80 token/s；
  \item \texttt{pp} 是算力受限：约 $\frac{48 \times 10^{12} \times \text{MFU}}{2 \times 7.6 \times 10^9}$ token/s。
  \item BF16 的 7B（15.2 GB）在 16GB 卡上几乎放不下——这就是你这张卡上量化的第一价值（第 11 篇）。
\end{itemize}

\textbf{要能解释}：为什么量化让 \texttt{tg} 提升了约 3 倍，而 \texttt{pp} 提升远没有那么多，甚至可能变慢？（提示：W4A16 在计算前要反量化，Prefill 是算力受限的）

\vspace{0.8em}\hrule\vspace{0.8em}

\section{[测试] 任务 5：投机解码}

在 vLLM（配置草稿模型或 n-gram 投机）或 llama.cpp（\texttt{llama-speculative}，目标 Qwen2.5-7B + 草稿 Qwen2.5-0.5B）上开启投机解码，测单流速度与接受率。
用第 10 篇的公式 $\frac{1 - \alpha^{K+1}}{(1 - \alpha)(Kc + 1)}$ 代入实测接受率和两模型的耗时比，预测加速比并与实测对比。再把并发调到 32，投机解码还有收益吗？为什么？

\vspace{0.8em}\hrule\vspace{0.8em}

\section{[阅读] 任务 6：源码阅读路线}

读源码的顺序很重要——从小到大：

\begin{enumerate}
  \item \textbf{\href{https://github.com/GeeeekExplorer/nano-vllm}{nano-vllm}}（约 1200 行 Python，实现了 vLLM 的核心：调度、块管理、前缀缓存、CUDA Graph、张量并行）。对照表：
\end{enumerate}

\begin{center}
\small
\begin{tabularx}{\textwidth}{p{3.5cm} X}
\toprule
\textbf{nano-vllm 里的部分} & \textbf{本仓库对应} \\
\midrule
调度器（Prefill / Decode、抢占） & Lab 12a、第 12 篇 §2、§3.4 \\
块管理器（引用计数、哈希前缀缓存） & Lab 12b 的 \texttt{BlockAllocator}、实验 4 \\
模型实现（Qwen） & Lab 07b 的 \texttt{QwenFromScratch} \\
注意力（调用 FlashAttention 的变长 / 分页接口） & Lab 09、第 9 篇 \\
CUDA Graph 捕获 & Lab 07b 性能分析、练习 5 \\
张量并行的线性层 & Lab 13、第 13 篇 \\
\bottomrule
\end{tabularx}
\end{center}

\begin{enumerate}
  \item \textbf{\href{https://github.com/pytorch-labs/gpt-fast}{gpt-fast}}（不到 1000 行，用 \texttt{torch.compile} + 量化 + 投机解码把单流速度逼近带宽极限）。
  \item \textbf{vLLM 本体}：先读官方文档的架构设计部分，再按"入口 $\rightarrow$ 引擎核心 $\rightarrow$ 调度器 $\rightarrow$ KV Cache 管理 $\rightarrow$ 模型执行器 $\rightarrow$ 注意力后端"的顺序读。
  \item \textbf{llama.cpp}：从 \texttt{ggml} 的计算图与量化格式（\texttt{Q4\_K} 等）入手，看它如何在 CPU / GPU 上做融合的反量化矩阵乘。
\end{enumerate}

\vspace{0.8em}\hrule\vspace{0.8em}

\section{[OK] 毕业标准}

完成下面这件事，你就真正走完了这条路线：

\begin{mathBox}[核心数学定理推导]
选定一个模型和你的 5060 Ti，写一份一页纸的报告：\textbf{预测}它在 vLLM 上的单流速度、最大并发、TTFT（带/不带前缀缓存）、W4 量化后的速度，给出每个预测的公式；然后\textbf{实测}，逐项解释预测与实测之间的差距来自哪里。
\end{mathBox}

如果每一个差距你都能说清楚是哪一篇、哪一个公式、哪一个工程因素造成的，你就不再是"库的使用者"了。


\section{实验核心源码精选与解析}

本实验配套有完整可运行的工业级验证代码，重点实现细节如下：


\subsection{源码实现：\texttt{bench\_openai\_server.py}}

\lstinputlisting[language=Python, caption={\texttt{bench\_openai\_server.py}}]{code/capstone_engines_practice_bench_openai_server.py}



\end{document}